% Induction électromagnétique — Physique 1ère TI — Cours signé OnBuch+
% Programme officiel MINESEC. Compiler avec : tectonic N16-induction-electromagnetique.tex
\newcommand{\DOCMATIERE}{Physique}
\newcommand{\DOCNIVEAU}{1ère TI}
\newcommand{\DOCMODULE}{Module 4 : Aspects énergétiques des circuits électriques}
\newcommand{\DOCLECON}{N16}
\newcommand{\DOCTITRE}{Induction électromagnétique}
% OnBuch+ — préambule « cours signés » ENRICHI (compilé avec tectonic / XeLaTeX)
% Système visuel « Pop » d'OnBuch+ : fond crème (jamais blanc), bordures encre
% épaisses, ombres « dures » décalées, titres Archivo Black en capitales.
% Version MATHÉMATIQUES : graphes (pgfplots/tikz), boîtes pédagogiques
% (définition, propriété, méthode, attention, le savais-tu, exercice résolu,
% exercice, TP, bilan), page de garde et signature OnBuch+ ancrée en bas.
%
% Macros attendues dans main.tex AVANT \input{preamble} :
%   \DOCMATIERE  \DOCNIVEAU  \DOCMODULE  \DOCLECON (numéro)  \DOCTITRE
\documentclass[11pt]{article}

\usepackage{fontspec}
\usepackage[french]{babel}
\usepackage[a4paper,margin=2cm,top=3.4cm,bottom=2.8cm,headheight=1.4cm,footskip=1.3cm]{geometry}
\usepackage{amsmath,amssymb,amsfonts}
\usepackage{mathtools} % \xmapsto, \coloneqq... souvent utilisés par les modèles
\usepackage{cancel} % simplifications barrées (bilans de forces)
% \abs{z} (module, valeur absolue) et \norm{u} (norme), utilisés par les modèles
\providecommand{\abs}{}\let\abs\relax\DeclarePairedDelimiter{\abs}{\lvert}{\rvert}
\providecommand{\norm}{}\let\norm\relax\DeclarePairedDelimiter{\norm}{\lVert}{\rVert}
\usepackage[table]{xcolor} % [table] : fournit aussi \rowcolors (alternance de lignes)
\usepackage{pagecolor}
\usepackage{tikz}
\usetikzlibrary{arrows.meta,positioning,calc,shapes.geometric,shapes.misc,shapes.arrows,decorations.pathmorphing,decorations.pathreplacing,decorations.markings,patterns,shadows,mindmap,trees,fit,backgrounds,matrix,intersections,angles,quotes}
\usepackage{pgfplots}
\usepackage[version=4]{mhchem} % équations nucléaires \ce{^{235}_{92}U}
\usepackage[european,siunitx]{circuitikz} % schémas électriques
\pgfplotsset{compat=1.18}
\usepgfplotslibrary{fillbetween,groupplots} % aires entre courbes (calcul intégral)
\usepackage{enumitem}
\usepackage{fancyhdr}
% [most] charge toutes les bibliothèques tcolorbox (listings, minted, raster,
% documentation...) et gonfle inutilement la mémoire TeX sur un document long
% (~300 boîtes) : on ne charge que ce qui est réellement utilisé.
\usepackage[skins,breakable]{tcolorbox}
\usepackage{graphicx}
\usepackage{colortbl}
\usepackage{booktabs}
\usepackage{tabularx}
\usepackage{diagbox} % case d'en-tête diagonale des tableaux à double entrée
% Colonnes extensibles centrée / gauche / droite (C, L, R), souvent utilisées par les modèles
\newcolumntype{C}{>{\centering\arraybackslash}X}
\newcolumntype{L}{>{\raggedright\arraybackslash}X}
\newcolumntype{R}{>{\raggedleft\arraybackslash}X}
\usepackage{array}
\usepackage{multicol}
\usepackage{siunitx}
% parse-numbers=false : accepte aussi \SI{2,5 \times 10^{-3}}{...}
\sisetup{locale=FR,output-decimal-marker={,},group-minimum-digits=4,parse-numbers=false}
\DeclareSIUnit\ohm{\ensuremath{\symup{\Omega}}} % U+2126 absent de la police
\DeclareSIUnit\division{div} % réglages d'oscilloscope (ms/div, V/div)
\DeclareSIUnit\tr{tr} % tours (vitesse de rotation, ex. \SI{1500}{\tr\per\minute})
\DeclareSIUnit\ch{ch} % cheval-vapeur (puissance moteur, unité usuelle hors SI)
\DeclareSIUnit\franccfa{FCFA} % franc CFA (coûts, contexte camerounais)
\DeclareSIUnit\dioptre{\ensuremath{\delta}} % vergence d'une lentille (dioptrie)
\usepackage{qrcode}
% \degree : commande fréquemment utilisée par les modèles pour un angle ou une
% température en dehors de \SI{}{} (non fournie par les paquets chargés).
\providecommand{\degree}{^{\circ}}
% \qed (fin de démonstration), utilisé par les modèles sans amsthm
\providecommand{\qed}{\hfill$\square$}
% \barre : double barre des valeurs interdites dans un tableau de variations (tkz-tab)
\providecommand{\barre}{\ensuremath{\|}}
% environnement proof (amsthm non chargé) utilisé par les modèles
\ifdefined\proof\else\newenvironment{proof}[1][Démonstration]{\par\noindent\textit{#1.}\ }{\qed\par}\fi
\usepackage{lastpage}
\usepackage{hyperref}
\hypersetup{hidelinks}

% ============================================================
% Palette « Pop » OnBuch+ — mêmes hex que la classe P (plus_ui.dart)
% ============================================================
\definecolor{poporange}{HTML}{F59321}
\definecolor{popdark}{HTML}{E07A0C}
\definecolor{popcream}{HTML}{FFF6E8}
\definecolor{popink}{HTML}{1C1714}
\definecolor{poppurple}{HTML}{7A5AE0}
\definecolor{popgreen}{HTML}{2E9E6A}
\definecolor{poppink}{HTML}{E0457B}
\definecolor{popblue}{HTML}{2F7DE1}
\definecolor{popgold}{HTML}{C98A12}
\definecolor{popmuted}{HTML}{8A7F76}
\definecolor{popline}{HTML}{EADFD2}
\definecolor{popgray}{HTML}{8A7F76}
\colorlet{popgrey}{popgray}
% Alias tolérants (le modèle nomme parfois les couleurs différemment)
\colorlet{popwhite}{white}
\colorlet{popblack}{popink}
\colorlet{popred}{poppink}
\colorlet{popyellow}{popgold}
% Teintes claires pour fonds de tableaux / zones
\colorlet{popblueL}{popblue!10!white}
\colorlet{poporangeL}{poporange!14!white}
\colorlet{popgreenL}{popgreen!12!white}
\colorlet{poppurpleL}{poppurple!12!white}
\colorlet{poppinkL}{poppink!12!white}
\colorlet{popgoldL}{popgold!14!white}

\pagecolor{popcream}

% ============================================================
% Typographie
% ============================================================
\defaultfontfeatures{Ligatures=TeX}
\setmainfont{PlusJakartaSans-Medium.ttf}[Path=../../../fonts/,BoldFont=PlusJakartaSans-Bold.ttf]
% Maths en sans-serif (Fira Math) avec chiffres et lettres droites de Plus
% Jakarta, pour que les formules se fondent dans le texte.
\usepackage[math-style=ISO,bold-style=ISO]{unicode-math}
\setmathfont{FiraMath-Regular.otf}[Path=../../../fonts/]
\setmathfont{PlusJakartaSans-Medium.ttf}[Path=../../../fonts/,range={up/num,up/{Latin,latin},\mathup}]
\usepackage{icomma} % virgule décimale sans espace en mode math
% Caractères mathématiques Unicode absents des polices de texte (Plus Jakarta,
% Archivo Black) : rendus par leur équivalent mathématique, en texte comme en
% formule — sinon ils s'affichent en carré vide (ex. « Arithmétique dans ℤ »).
\usepackage{newunicodechar}
\newunicodechar{ℝ}{\ensuremath{\mathbb{R}}}
\newunicodechar{ℤ}{\ensuremath{\mathbb{Z}}}
\newunicodechar{ℂ}{\ensuremath{\mathbb{C}}}
\newunicodechar{ℕ}{\ensuremath{\mathbb{N}}}
\newunicodechar{ℚ}{\ensuremath{\mathbb{Q}}}
\newunicodechar{𝔻}{\ensuremath{\mathbb{D}}}
\newunicodechar{α}{\ensuremath{\alpha}}
\newunicodechar{β}{\ensuremath{\beta}}
\newunicodechar{γ}{\ensuremath{\gamma}}
\newunicodechar{δ}{\ensuremath{\delta}}
\newunicodechar{ε}{\ensuremath{\varepsilon}}
\newunicodechar{θ}{\ensuremath{\theta}}
\newunicodechar{λ}{\ensuremath{\lambda}}
\newunicodechar{μ}{\ensuremath{\mu}}
\newunicodechar{σ}{\ensuremath{\sigma}}
\newunicodechar{τ}{\ensuremath{\tau}}
\newunicodechar{φ}{\ensuremath{\varphi}}
\newunicodechar{ψ}{\ensuremath{\psi}}
\newunicodechar{ω}{\ensuremath{\omega}}
\newunicodechar{Σ}{\ensuremath{\Sigma}}
% majuscules produites par \MakeUppercase dans les titres (α-aminés -> Α)
\newunicodechar{Α}{\ensuremath{\alpha}}
\newunicodechar{Β}{\ensuremath{\beta}}
\newunicodechar{Γ}{\ensuremath{\Gamma}}
\newunicodechar{Δ}{\ensuremath{\Delta}}
\newunicodechar{Ε}{\ensuremath{\varepsilon}}
\newunicodechar{Θ}{\ensuremath{\Theta}}
\newunicodechar{Λ}{\ensuremath{\Lambda}}
\newunicodechar{Μ}{\ensuremath{\mu}}
\newunicodechar{Τ}{\ensuremath{\tau}}
\newunicodechar{Φ}{\ensuremath{\Phi}}
\newunicodechar{Ψ}{\ensuremath{\Psi}}
\newunicodechar{Ω}{\ensuremath{\Omega}}
\newunicodechar{↦}{\ensuremath{\mapsto}}
\newunicodechar{∈}{\ensuremath{\in}}
\newunicodechar{∉}{\ensuremath{\notin}}
\newunicodechar{∧}{\ensuremath{\wedge}}
\newunicodechar{⇔}{\ensuremath{\Leftrightarrow}}
\newunicodechar{⟺}{\ensuremath{\Longleftrightarrow}}
\newunicodechar{⇒}{\ensuremath{\Rightarrow}}
\newunicodechar{⟹}{\ensuremath{\Longrightarrow}}
\newunicodechar{∘}{\ensuremath{\circ}}
\newunicodechar{∪}{\ensuremath{\cup}}
\newunicodechar{∩}{\ensuremath{\cap}}
\newunicodechar{≡}{\ensuremath{\equiv}}
\newunicodechar{⊂}{\ensuremath{\subset}}
\newunicodechar{⊥}{\ensuremath{\perp}}
\newunicodechar{∃}{\ensuremath{\exists}}
\newunicodechar{∀}{\ensuremath{\forall}}
\newunicodechar{∛}{\ensuremath{\sqrt[3]{\,}}}
\newunicodechar{∜}{\ensuremath{\sqrt[4]{\,}}}
\newunicodechar{⃗}{\ensuremath{{}^{\to}}}
\newunicodechar{ⁿ}{\ifmmode^{n}\else\textsuperscript{\ensuremath{n}}\fi}
\newunicodechar{ˣ}{\ifmmode^{x}\else\textsuperscript{\ensuremath{x}}\fi}
\newunicodechar{ᵇ}{\ifmmode^{b}\else\textsuperscript{\ensuremath{b}}\fi}
\newunicodechar{ʳ}{\ifmmode^{r}\else\textsuperscript{\ensuremath{r}}\fi}
\newunicodechar{ᵘ}{\ifmmode^{u}\else\textsuperscript{\ensuremath{u}}\fi}
\newunicodechar{ʸ}{\ifmmode^{y}\else\textsuperscript{\ensuremath{y}}\fi}
\newunicodechar{ᵃ}{\ifmmode^{a}\else\textsuperscript{\ensuremath{a}}\fi}
\newunicodechar{ᵉ}{\ifmmode^{e}\else\textsuperscript{\ensuremath{e}}\fi}
\newunicodechar{ᵏ}{\ifmmode^{k}\else\textsuperscript{\ensuremath{k}}\fi}
\newunicodechar{ᵖ}{\ifmmode^{p}\else\textsuperscript{\ensuremath{p}}\fi}
\newunicodechar{ᴺ}{\ifmmode^{N}\else\textsuperscript{\ensuremath{N}}\fi}
\newunicodechar{ᶜ}{\ifmmode^{c}\else\textsuperscript{\ensuremath{c}}\fi}
\newunicodechar{ʰ}{\ifmmode^{h}\else\textsuperscript{\ensuremath{h}}\fi}
\newunicodechar{ᵠ}{\ifmmode^{q}\else\textsuperscript{\ensuremath{q}}\fi}
\newunicodechar{ₙ}{\ifmmode_{n}\else\textsubscript{\ensuremath{n}}\fi}
\newunicodechar{ₐ}{\ifmmode_{a}\else\textsubscript{\ensuremath{a}}\fi}
\newunicodechar{ᵢ}{\ifmmode_{i}\else\textsubscript{\ensuremath{i}}\fi}
\newunicodechar{ᵣ}{\ifmmode_{r}\else\textsubscript{\ensuremath{r}}\fi}
\newunicodechar{ₖ}{\ifmmode_{k}\else\textsubscript{\ensuremath{k}}\fi}
\newunicodechar{ᵦ}{\ifmmode_{\beta}\else\textsubscript{\ensuremath{\beta}}\fi}
\newunicodechar{ₓ}{\ifmmode_{x}\else\textsubscript{\ensuremath{x}}\fi}
\newfontfamily\popdisplay{ArchivoBlack-Regular.ttf}[Path=../../../fonts/]
\newfontfamily\poplogo{SpaceGrotesk-Bold.ttf}[Path=../../../fonts/]
\newfontfamily\popsemi{PlusJakartaSans-SemiBold.ttf}[Path=../../../fonts/]
\newfontfamily\popxbold{PlusJakartaSans-ExtraBold.ttf}[Path=../../../fonts/]
\newfontfamily\popxboldls{PlusJakartaSans-ExtraBold.ttf}[Path=../../../fonts/,LetterSpace=3.0]

\setlist[enumerate]{leftmargin=1.7em,itemsep=5pt,topsep=4pt}
\setlist[itemize]{leftmargin=1.7em,itemsep=4pt,topsep=4pt,label={\color{poporange}\textbullet}}
\setlength{\parindent}{0pt}
\setlength{\parskip}{5pt}
\renewcommand{\arraystretch}{1.25}

% pgfplots : style Pop par défaut
\pgfplotsset{
  every axis/.append style={axis line style={popink,line width=1pt},tick style={popink},
    grid style={popline,line width=0.5pt},label style={font=\small},tick label style={font=\footnotesize}},
  popcurve/.style={poporange,line width=1.8pt,no markers,smooth},
}
% Même style disponible dans un \draw TikZ simple (hors pgfplots)
\tikzset{popcurve/.style={poporange,line width=1.8pt}}
\tikzset{popgrid/.style={popline,very thin}}
\colorlet{popcurve}{poporange} % le nom du style est parfois utilisé comme couleur

% ============================================================
% Pill / squiggle
% ============================================================
\newcommand{\poppill}[3]{%
  \tikz[baseline=-0.55ex]\node[draw=popink,line width=1.2pt,rounded corners=2.6pt,%
    fill=#2,inner xsep=6.5pt,inner ysep=3pt]{%
    \popxboldls\fontsize{7}{7}\selectfont\color{#3}\MakeUppercase{#1}};%
}
\newcommand{\popsquiggle}[1]{%
  \tikz[baseline=-0.4ex]{
    \draw[#1,line width=2.6pt,line cap=round]
      plot [smooth] coordinates {(0,0.09) (0.34,0.01) (0.68,0.21) (1.02,0.01) (1.36,0.10)};
  }%
}
% Mot-clé surligné (marqueur orange sous le texte)
\newcommand{\cle}[1]{{\popxbold\color{popdark}#1}}

% ============================================================
% En-tête / pied
% ============================================================
\pagestyle{fancy}
\fancyhf{}
\renewcommand{\headrulewidth}{0pt}
\renewcommand{\footrulewidth}{0pt}
\lhead{\raisebox{-0.15ex}{\poplogo\fontsize{13}{13}\selectfont\color{popink}OnBuch\color{poporange}+}}
\rhead{\poppill{\DOCMATIERE\ \textbullet\ \DOCNIVEAU\ \textbullet\ Leçon \DOCLECON}{white}{popink}}
\lfoot{\popxbold\fontsize{7.6}{7.6}\selectfont\color{popmuted}\MakeUppercase{Cours signé OnBuch+}\ \textbullet\ {\popsemi\DOCTITRE}}
\rfoot{\tikz[baseline=-0.6ex]\node[rounded corners=8.5pt,fill=popink,minimum height=17pt,minimum width=17pt,inner xsep=5pt,inner sep=0]{\popxbold\fontsize{8}{8}\selectfont\color{popcream}\thepage\,/\,\pageref*{LastPage}};}

% ============================================================
% Titres
% ============================================================
\newcommand{\chaptitre}[2]{%
  \begin{tcolorbox}[enhanced,colback=poporange,colframe=popink,boxrule=2.6pt,arc=11pt,
      drop shadow=popink,left=15pt,right=15pt,top=12pt,bottom=11pt,before skip=3pt,after skip=18pt]
    \poppill{#2}{popink}{popcream}\par\vspace{9pt}
    {\raggedright\hyphenpenalty=10000\exhyphenpenalty=10000\popdisplay\fontsize{22}{24}\selectfont\color{popcream}\MakeUppercase{#1}\par}
  \end{tcolorbox}
}
% Section numérotée : pastille orange avec le numéro + titre Archivo
\newcounter{popsec}
\newcommand{\coursec}[1]{%
  \stepcounter{popsec}\setcounter{popsub}{0}%
  \par\vspace{16pt}\noindent
  \begin{minipage}{\linewidth}
  \tikz[baseline=-0.7ex]\node[draw=popink,line width=1.6pt,fill=poporange,rounded corners=4pt,minimum size=22pt,inner sep=0]{\popdisplay\fontsize{12}{12}\selectfont\color{popcream}\thepopsec};%
  \hspace{8pt}{\popdisplay\fontsize{15}{17}\selectfont\color{popink}\MakeUppercase{#1}}\par
  \vspace{4pt}\noindent{\color{poporange}\rule{36pt}{3.6pt}}
  \end{minipage}\par\nopagebreak\vspace{8pt}
}
\newcounter{popsub}
\newcommand{\courssub}[1]{%
  \stepcounter{popsub}%
  \par\vspace{9pt}\noindent{\popxbold\fontsize{11.5}{13}\selectfont\color{popink}{\color{poporange}\thepopsec.\thepopsub}\hspace{6pt}#1}\par\nopagebreak\vspace{4pt}
}

% ============================================================
% Boîtes pédagogiques — même recette : fond blanc, bordure épaisse colorée,
% ombre dure encre, badge plein. Titre optionnel [#1].
% ============================================================
\tcbset{popbase/.style={breakable,enhanced,colback=white,boxrule=2.3pt,arc=9pt,
  shadow={3pt}{-3pt}{0pt}{popink},left=14pt,right=14pt,top=11pt,bottom=11pt,before skip=11pt,after skip=15pt,
  fontupper=\popsemi\fontsize{10.6}{14.5}\selectfont\color{popink},
  fontlower=\popsemi\fontsize{10.6}{14.5}\selectfont\color{popink}}}
% Coche dessinée (le glyphe ✓ n'existe pas dans les polices du document)
\newcommand{\popcheck}[1]{\tikz[baseline=-0.5ex]\draw[#1,line width=1.6pt,line cap=round,line join=round](-0.13,0.0)--(-0.04,-0.09)--(0.13,0.1);}
\providecommand{\coche}{\popcheck{popgreen}}
\providecommand{\resultat}[1]{\par\smallskip\noindent\fcolorbox{popgreen}{popgreenL}{\parbox{\dimexpr\linewidth-2\fboxsep-2\fboxrule\relax}{\textbf{Résultat.} #1}}\par\smallskip}
\newcommand{\popboxhead}[3]{\poppill{#1}{#2}{white}\ifx\relax#3\relax\else\hspace{6pt}{\popxbold\fontsize{10.6}{12}\selectfont\color{popink}#3}\fi\par\vspace{7pt}}

\newtcolorbox{aretenir}[1][]{popbase,colframe=popgreen,before upper={\popboxhead{À retenir}{popgreen}{#1}}}
\newtcolorbox{exemplebox}[1][]{popbase,colframe=poporange,before upper={\popboxhead{Exemple}{poporange}{#1}}}
\newtcolorbox{definition}[1][]{popbase,colframe=popblue,before upper={\popboxhead{Définition}{popblue}{#1}}}
\newtcolorbox{propriete}[1][]{popbase,colframe=poppurple,before upper={\popboxhead{Propriété}{poppurple}{#1}}}
\newtcolorbox{methode}[1][]{popbase,colframe=popgold,before upper={\popboxhead{Méthode}{popgold}{#1}}}
\newtcolorbox{attention}[1][]{popbase,colframe=poppink,before upper={\popboxhead{Attention}{poppink}{#1}}}
\newtcolorbox{experience}[1][]{popbase,colframe=popblue,colback=popblueL,before upper={\popboxhead{Expérience}{popblue}{#1}}}
% « Le savais-tu ? » : carte violette pleine (contexte, culture, Cameroun)
\newtcolorbox{savaistu}[1][]{popbase,colback=poppurple,colframe=popink,
  fontupper=\popsemi\fontsize{10.4}{14.2}\selectfont\color{white},
  before upper={\poppill{Le savais-tu\,?}{white}{poppurple}\ifx\relax#1\relax\else\hspace{6pt}{\popxbold\color{white}#1}\fi\par\vspace{7pt}}}
% Exercice résolu : énoncé en haut, \tcblower puis la solution sur fond crème
\newcounter{popexo}\newcounter{popexr}
\newtcolorbox{exoresolu}[1][]{popbase,colframe=poporange,colbacklower=popcream,
  segmentation style={popink,line width=1.2pt,dashed},
  before upper={\stepcounter{popexr}\popboxhead{Exercice résolu \thepopexr}{poporange}{#1}},
  before lower={\poppill{Solution}{popink}{popcream}\par\vspace{6pt}}}
% Exercice (non résolu en place) — la correction est dans la partie Corrigés
\newtcolorbox{exercice}[1][]{popbase,colframe=popink,
  before upper={\stepcounter{popexo}\popboxhead{Exercice \thepopexo}{popink}{#1}}}
% Corrigé
\newtcolorbox{corrige}[1][]{popbase,colframe=popgreen,colback=popgreenL,
  before upper={\popboxhead{Corrigé}{popgreen}{#1}}}
% Objectifs / prérequis / bilan
\newtcolorbox{objectifs}{popbase,colframe=popink,colback=poporangeL,
  before upper={\popboxhead{Objectifs de la leçon}{popink}{}}}
\newtcolorbox{prerequis}{popbase,colframe=popink,colback=popblueL,
  before upper={\popboxhead{Prérequis}{popblue}{}}}
\newtcolorbox{bilan}{popbase,colframe=popink,colback=white,boxrule=2.8pt,
  before upper={\popboxhead{Fiche bilan}{poporange}{L'essentiel à connaître pour l'examen}}}
% Figure encadrée avec légende
\newtcolorbox{popfigure}[1][]{enhanced,breakable=false,colback=white,colframe=popline,boxrule=1.4pt,arc=8pt,
  left=10pt,right=10pt,top=10pt,bottom=8pt,before skip=10pt,after skip=12pt,halign=center,
  lower separated=false,halign lower=center,
  fontlower=\popsemi\fontsize{9}{11.5}\selectfont\color{popmuted},
  #1}
\newcommand{\legende}[1]{\par\vspace{6pt}{\popsemi\fontsize{9}{11.5}\selectfont\color{popmuted}#1}}

% ============================================================
% Signature OnBuch+
% ============================================================
\newcommand{\poplogoinline}[1]{{\poplogo\fontsize{#1}{#1}\selectfont\color{popink}OnBuch\color{poporange}+}}
\newcommand{\signatureonbuch}{%
  \begin{tcolorbox}[enhanced,colback=popink,colframe=popink,boxrule=2.2pt,arc=10pt,
      drop shadow=poporange,left=14pt,right=14pt,top=10pt,bottom=10pt,before skip=0pt,after skip=0pt]
    \tikz[baseline=-0.7ex]\node[circle,fill=popgreen,draw=popcream,line width=1.3pt,minimum size=22pt,inner sep=0]{\popcheck{white}};%
    \hspace{9pt}\begin{minipage}[c]{0.62\linewidth}
      {\popxbold\fontsize{12}{13}\selectfont\color{popcream}Signé\ \textbullet\ {\poplogo OnBuch\color{poporange}+}}\par\vspace{2pt}
      {\raggedright\popsemi\fontsize{8}{10.5}\selectfont\color{popline}\MakeUppercase{Équipe pédagogique OnBuch+}\\\MakeUppercase{Conforme au programme officiel MINESEC}\par}
    \end{minipage}\hfill
    {\poplogo\fontsize{22}{22}\selectfont\color{popcream}OnBuch\color{poporange}+}
  \end{tcolorbox}%
}

% ============================================================
% Page de garde
% #1 titre · #2 matière · #3 niveau · #4 module · #5 n° leçon · #6 durée ·
% #7 description / objectifs (texte ou itemize)
% ============================================================
\newcommand{\pagegarde}[7]{%
  \thispagestyle{empty}
  \newgeometry{a4paper,margin=1.7cm,top=1.6cm,bottom=1.4cm}%
  \begin{tikzpicture}[remember picture,overlay]
    \fill[poporange!18] ([xshift=-3.2cm,yshift=-3.6cm]current page.north east) circle (4.6cm);
    \fill[poppurple!14] ([xshift=2.4cm,yshift=7.6cm]current page.south west) circle (3.8cm);
  \end{tikzpicture}%
  {\poplogo\fontsize{30}{30}\selectfont\color{popink}OnBuch\color{poporange}+}\hfill
  \raisebox{0.6ex}{\poppill{Cours signé \textbullet\ Premium}{popink}{popcream}}\par
  \vspace{1pt}\popsquiggle{poppurple}\par
  \vspace{8pt}
  \begin{tcolorbox}[enhanced,colback=poporange,colframe=popink,boxrule=2.8pt,arc=13pt,
      drop shadow=popink,left=18pt,right=18pt,top=12pt,bottom=13pt,after skip=10pt]
    \poppill{#2 \textbullet\ #3}{popink}{popcream}\hspace{5pt}\poppill{#4}{white}{popink}\par\vspace{8pt}
    {\popdisplay\fontsize{14}{14}\selectfont\color{popink}LEÇON #5}\par\vspace{4pt}
    {\raggedright\hyphenpenalty=10000\exhyphenpenalty=10000\popdisplay\fontsize{25}{27}\selectfont\color{popcream}\MakeUppercase{#1}\par}
    \vspace{8pt}
    {\popxbold\fontsize{9.5}{9.5}\selectfont\color{popink}\MakeUppercase{Durée conseillée : #6}}
  \end{tcolorbox}
  \begin{tcolorbox}[enhanced,colback=white,colframe=popink,boxrule=2.2pt,arc=10pt,
      drop shadow=popink,left=16pt,right=16pt,top=10pt,bottom=10pt,after skip=8pt]
    \poppill{Dans cette leçon}{poppurple}{white}\par\vspace{5pt}
    \popsemi\fontsize{9.3}{12.6}\selectfont\color{popink}#7
  \end{tcolorbox}
  \noindent\begin{minipage}[t]{0.487\linewidth}
    \begin{tcolorbox}[enhanced,colback=popgreen,colframe=popink,boxrule=2.2pt,arc=10pt,
        drop shadow=popink,left=11pt,right=11pt,top=10pt,bottom=10pt,height=3.1cm,valign=center]
      \tcbox[colback=white,colframe=popink,boxrule=1.3pt,arc=5pt,boxsep=0pt,nobeforeafter,
        left=3pt,right=3pt,top=3pt,bottom=3pt,valign=center]{\qrcode[height=1.9cm]{https://play.google.com/store/apps/details?id=com.luvvix.onbuch&hl=fr}}%
      \hspace{8pt}\begin{minipage}[c]{0.5\linewidth}
        {\popdisplay\fontsize{11}{12.5}\selectfont\color{white}\MakeUppercase{Télécharge OnBuch}}\par\vspace{4pt}
        {\raggedright\popsemi\fontsize{8.4}{11}\selectfont\color{white}Gratuit \textbullet\ Google Play\\Tuteur IA, annales, concours, résultats\par}
      \end{minipage}
    \end{tcolorbox}
  \end{minipage}\hfill
  \begin{minipage}[t]{0.487\linewidth}
    \begin{tcolorbox}[enhanced,colback=poppurple,colframe=popink,boxrule=2.2pt,arc=10pt,
        drop shadow=popink,left=11pt,right=11pt,top=10pt,bottom=10pt,height=3.1cm,valign=center]
      \tikz[baseline=-0.7ex]\node[circle,fill=white,draw=popink,line width=1.3pt,minimum size=1.6cm,inner sep=0]{\popdisplay\fontsize{24}{24}\selectfont\color{poppurple}+};%
      \hspace{8pt}\begin{minipage}[c]{0.6\linewidth}
        {\popdisplay\fontsize{11}{12.5}\selectfont\color{white}\MakeUppercase{Passe à OnBuch+}}\par\vspace{4pt}
        {\raggedright\popsemi\fontsize{8.4}{11}\selectfont\color{white}Tous les cours signés, Tuteur IA illimité, fiches PDF — onglet OnBuch+ de l'app\par}
      \end{minipage}
    \end{tcolorbox}
  \end{minipage}
  \vspace{8pt}
  \popcontenu
  \vfill
  \signatureonbuch
  \restoregeometry
  \newpage
}

% Rangée « Ce cours contient » (page de garde)
\newcommand{\poptile}[3]{% #1 couleur #2 gros texte #3 légende
  \begin{tcolorbox}[enhanced,colback=#1,colframe=popink,boxrule=2pt,arc=9pt,shadow={2.5pt}{-2.5pt}{0pt}{popink},
      left=6pt,right=6pt,top=9pt,bottom=9pt,halign=center,valign=center,height=1.75cm,width=\linewidth,nobeforeafter]
    {\popdisplay\fontsize{12.5}{13}\selectfont\color{white}\MakeUppercase{#2}}\par\vspace{3pt}
    {\centering\popsemi\fontsize{7.8}{9.5}\selectfont\color{white}#3\par}
  \end{tcolorbox}}
\newcommand{\popcontenu}{%
  \par\noindent{\popxboldls\fontsize{7.5}{7.5}\selectfont\color{popmuted}CE COURS SIGNÉ CONTIENT}\par\vspace{7pt}
  \noindent\begin{minipage}[t]{0.235\linewidth}\poptile{popblue}{Cours}{complet, illustré, pas à pas}\end{minipage}\hfill
  \begin{minipage}[t]{0.235\linewidth}\poptile{poporange}{Méthodes}{et exercices résolus}\end{minipage}\hfill
  \begin{minipage}[t]{0.235\linewidth}\poptile{poppink}{Exercices}{type examen + corrigés}\end{minipage}\hfill
  \begin{minipage}[t]{0.235\linewidth}\poptile{popgreen}{Fiche bilan}{l'essentiel à retenir}\end{minipage}\par}

% Sommaire
\newtcolorbox{sommaire}{enhanced,colback=white,colframe=popink,boxrule=2.2pt,arc=10pt,
  drop shadow=popink,left=18pt,right=18pt,top=13pt,bottom=13pt,before skip=6pt,after skip=20pt,
  before upper={\poppill{Sommaire}{poppurple}{white}\par\vspace{9pt}\popsemi\fontsize{10.6}{14.5}\selectfont\color{popink}}}

% Signature de fin de cours — poussée tout en bas de la dernière page
\newcommand{\findecours}{%
  \par\vfill\nopagebreak
  \begin{center}\popsquiggle{poporange}\end{center}\vspace{4pt}
  \signatureonbuch
}
\AtBeginDocument{\def\llbracket{\mathopen{[\mkern-3.5mu[}}\def\rrbracket{\mathclose{]\mkern-3.5mu]}}} % intervalles d'entiers (glyphe ⟦ absent de la police)
% Glyphes absents de Fira Math : redéfinis à partir de glyphes disponibles
\AtBeginDocument{%
  \def\setminus{\mathbin{\backslash}}\let\smallsetminus\setminus
  \def\vdots{\mathord{\vbox{\baselineskip3.2pt\lineskiplimit0pt\kern4pt\hbox{.}\hbox{.}\hbox{.}}}}%
  \def\ddots{\mathinner{\mkern1mu\raise6.5pt\hbox{.}\mkern2mu\raise3.6pt\hbox{.}\mkern2mu\raise0.7pt\hbox{.}\mkern1mu}}%
  \def\triangle{\mathord{\scalebox{0.9}{$\symup{\Delta}$}}}%
  \def\nabla{\mathord{\raisebox{\depth}{\scalebox{1}[-1]{$\symup{\Delta}$}}}}%
  \def\checkmark{\popcheck{popdark}}%
  \def\star{\mathbin{\ast}}\def\bigstar{\mathord{\ast}}%
  \def\diamond{\mathbin{\scalebox{0.6}{$\Diamond$}}}%
  \def\bigcup{\mathop{\mathchoice{\raisebox{-0.3em}{\scalebox{1.7}{$\cup$}}}{\scalebox{1.25}{$\cup$}}{\cup}{\cup}}\displaylimits}%
  \def\bigcap{\mathop{\mathchoice{\raisebox{-0.3em}{\scalebox{1.7}{$\cap$}}}{\scalebox{1.25}{$\cap$}}{\cap}{\cap}}\displaylimits}%
  \def\lesseqqgtr{\mathrel{\lessgtr}}\def\bigoplus{\mathop{\oplus}\displaylimits}%
}
\providecommand{\ssi}{si et seulement si} % abréviation fréquente
\AtBeginDocument{\providecommand{\wideparen}{\overparen}} % arc de cercle

%% FIGFIX-BEGIN v4 — correctifs globaux des figures (docs/onbuch-generation/tools/figfix)
\usepackage{adjustbox}\usepackage{etoolbox}
% toute figure TikZ/pgfplots plus large que la ligne est réduite à la largeur disponible
\BeforeBeginEnvironment{tikzpicture}{\begin{adjustbox}{max width=\linewidth}}
\AfterEndEnvironment{tikzpicture}{\end{adjustbox}}
% légendes pgfplots lisibles par-dessus les courbes
\pgfplotsset{every axis/.append style={legend style={fill=white,fill opacity=0.92,text opacity=1,draw=black!15,inner sep=3pt}}}
% étiquettes d'arêtes (syntaxe edge["..."]) avec fond blanc
\tikzset{every edge quotes/.append style={fill=white,inner sep=1.2pt}}
%% FIGFIX-END
\begin{document}

\pagegarde{Induction électromagnétique}{Physique}{1ère TI}{Module 4 : Aspects énergétiques des circuits électriques}{N16}{5 h}{Cette leçon introduit le phénomène d'induction électromagnétique : définition du flux magnétique, loi de Lenz, expression de la force électromotrice induite et phénomène d'auto-induction dans une bobine. Elle fournit les outils de calcul (flux, f.é.m. induite, inductance) indispensables pour comprendre alternateurs, transformateurs et l'interprétation énergétique des circuits.
\vspace{4pt}
\begin{multicols}{2}\small
\begin{itemize}
\item À la fin de cette leçon, je sais décrire et interpréter les expériences mettant en évidence l'induction électromagnétique (aimant-bobine, barreau sur rails).
\item À la fin de cette leçon, je sais distinguer le circuit inducteur du circuit induit dans une situation donnée.
\item À la fin de cette leçon, je sais définir le flux magnétique et calculer sa valeur à travers une surface plane (cas uniforme, surface orientée).
\item À la fin de cette leçon, je sais énoncer la loi de Lenz et l'utiliser pour déterminer le sens du courant induit.
\item À la fin de cette leçon, je sais appliquer la loi de Faraday e = -dΦ/dt pour déterminer la valeur d'une f.é.m. induite.
\item À la fin de cette leçon, je sais mettre en évidence expérimentalement le phénomène d'auto-induction (étincelle de rupture, retard à l'établissement du courant).
\end{itemize}\end{multicols}}

\begin{sommaire}
\begin{multicols}{2}
\begin{enumerate}
\item Mise en évidence expérimentale de l'induction
\item Flux magnétique
\item Loi de Lenz et sens du courant induit
\item F.é.m. induite : loi de Faraday
\item Auto-induction et inductance
\item Applications et bilan énergétique
\item Activité d'intégration
\item Méthodes et exercices résolus
\item Exercices
\item Corrigés des exercices
\item Fiche bilan
\end{enumerate}
\end{multicols}
\end{sommaire}

\begin{objectifs}
\begin{itemize}
\item À la fin de cette leçon, je sais décrire et interpréter les expériences mettant en évidence l'induction électromagnétique (aimant-bobine, barreau sur rails).
\item À la fin de cette leçon, je sais distinguer le circuit inducteur du circuit induit dans une situation donnée.
\item À la fin de cette leçon, je sais définir le flux magnétique et calculer sa valeur à travers une surface plane (cas uniforme, surface orientée).
\item À la fin de cette leçon, je sais énoncer la loi de Lenz et l'utiliser pour déterminer le sens du courant induit.
\item À la fin de cette leçon, je sais appliquer la loi de Faraday e = -dΦ/dt pour déterminer la valeur d'une f.é.m. induite.
\item À la fin de cette leçon, je sais mettre en évidence expérimentalement le phénomène d'auto-induction (étincelle de rupture, retard à l'établissement du courant).
\item À la fin de cette leçon, je sais définir l'inductance L d'une bobine, citer son unité et calculer la f.é.m. d'auto-induction e = -L di/dt.
\item À la fin de cette leçon, je sais interpréter énergétiquement l'induction (conversion mécanique → électrique) dans un alternateur.
\end{itemize}
\end{objectifs}

\begin{prerequis}
\begin{itemize}
\item Champ magnétique : vecteur B, spectre d'un aimant et d'une bobine, règle du tire-bouchon / de la main droite (leçon sur les champs magnétiques).
\item Intensité du courant, tension, loi d'Ohm et conventions récepteur/générateur (Seconde et début de Première).
\item Notion de dérivée et de taux de variation d'une grandeur (mathématiques, Première C).
\item Produit scalaire de deux vecteurs et angle entre deux directions (mathématiques).
\item Bilan énergétique dans un circuit : puissance P = UI, conversions d'énergie (Module 4, leçons précédentes).
\end{itemize}
\end{prerequis}

\newpage

\coursec{Mise en évidence expérimentale de l'induction}

À Yaoundé comme à Ngaoundéré, quand ENEO coupe le courant, les groupes électrogènes prennent le relais : comment un moteur qui tourne produit-il de l'électricité sans aucune pile ? De même, le chargeur de téléphone posé sur une table transforme la tension du secteur grâce à un transformateur qui n'a pourtant aucun contact électrique direct entre ses deux bobines. Plus intriguant encore : approchez un aimant d'une bobine reliée à un ampèremètre, et l'aiguille dévie alors qu'aucune source n'est branchée ! D'où vient ce courant ? C'est le phénomène d'\cle{induction électromagnétique}, base de toute la production d'énergie électrique au Cameroun (barrages de Songloulou, Édéa), que cette leçon va expliquer puis quantifier.

\textbf{Problématique :}\par Dans quelles conditions un circuit électrique \emph{sans générateur} peut-il être le siège d'une tension et d'un courant ? Nous allons répondre par l'expérience, en dégageant les constats qui fonderont toute la leçon.

\courssub{Expérience 1 : aimant en mouvement devant une bobine}

\begin{experience}[Aimant droit et bobine reliée à un galvanomètre]
\textbf{Matériel :} un aimant droit (pôles N et S), une bobine de plusieurs centaines de spires, un galvanomètre à zéro central (ou un millivoltmètre), des fils de connexion.

\textbf{Protocole :} la bobine est reliée au galvanomètre, \emph{sans aucun générateur}. On approche puis on éloigne le pôle Nord de l'aimant de la bobine, puis on recommence avec le pôle Sud.

\textbf{Observations :}
\begin{itemize}
\item Quand l'aimant \textbf{se déplace} (approche ou recul), l'aiguille du galvanomètre dévie : un courant circule dans la bobine.
\item Quand l'aimant est \textbf{immobile}, même très proche de la bobine, l'aiguille reste à zéro : aucun courant.
\item Le sens de déviation \textbf{s'inverse} quand on inverse le sens du mouvement (approche $\rightarrow$ recul) ou quand on permute les pôles (N $\rightarrow$ S).
\item Plus le mouvement est \textbf{rapide}, plus la déviation est grande.
\end{itemize}

\textbf{Interprétation :} le déplacement de l'aimant fait varier le champ magnétique qui traverse la bobine. Cette variation crée dans la bobine une tension — appelée \cle{force électromotrice induite} (f.é.m. induite), notée $e$ — qui fait circuler un \cle{courant induit} puisque le circuit est fermé.
\end{experience}

\begin{popfigure}
\begin{center}
\begin{tikzpicture}[scale=1.0,>=Stealth]
% Aimant
\fill[popblue!70] (-4.2,0.4) rectangle (-2.8,1.4);
\fill[poporange!80] (-2.8,0.4) rectangle (-1.4,1.4);
\node[white] at (-3.5,0.9) {\textbf{S}};
\node[white] at (-2.1,0.9) {\textbf{N}};
% Flèche de mouvement
\draw[->,very thick,popdark] (-1.2,0.9) -- (-0.2,0.9) node[midway,above]{$\vec{v}$};
% Bobine (spires)
\foreach \x in {0.4,0.7,1.0,1.3}{
  \draw[thick,popink] (\x,0.2) ellipse (0.12 and 0.85);
}
\node[popink] at (0.85,-0.9) {bobine ($N$ spires)};
% Fils vers galvanomètre
\draw[thick] (0.4,-0.65) -- (0.4,-1.6) -- (2.6,-1.6);
\draw[thick] (1.3,-0.65) -- (1.3,-2.2) -- (2.6,-2.2);
% Galvanomètre
\draw[thick,fill=popcream] (3.3,-1.9) circle (0.7);
\node at (3.3,-2.35) {\small G};
\draw[thick] (3.3,-1.9) -- (3.62,-1.45);
\draw[thick,popmuted] (2.85,-1.55) arc (150:30:0.55);
\draw[->,thick,poporange] (3.3,-1.9) -- (3.75,-1.55);
\node at (3.3,-2.85) {\small galvanomètre};
\node[popmuted] at (2.7,-1.15) {\small $0$};
\end{tikzpicture}
\end{center}
\legende{Expérience de l'aimant et de la bobine : pendant le déplacement de l'aimant (flèche $\vec{v}$), l'aiguille du galvanomètre dévie ; elle revient à zéro dès que l'aimant s'arrête.}
\end{popfigure}

\begin{attention}[Un champ constant n'induit rien]
Erreur fréquente : croire qu'un champ magnétique \emph{constant}, même très intense, produit un courant dans la bobine. C'est faux : quand l'aimant est immobile contre la bobine, le champ en son sein est fort mais \textbf{ne varie pas}, et le galvanomètre indique zéro. Seule la \textbf{variation} du champ magnétique à travers le circuit (plus précisément la variation du \emph{flux}, voir section 2) crée une f.é.m. induite.
\end{attention}

\courssub{Expérience 2 : bobine mobile devant un aimant fixe}

Reprenons le même montage, mais cette fois l'aimant est fixé sur la paillasse et c'est la bobine (solidaire du galvanomètre) que l'on déplace vers l'aimant ou qu'on en éloigne.

\textbf{Observations :} les résultats sont \emph{rigoureusement identiques} à ceux de l'expérience 1 : déviation pendant le mouvement, zéro à l'arrêt, inversion du sens de déviation avec le sens du mouvement, déviation d'autant plus grande que le mouvement est rapide.

\begin{aretenir}[Ce qui compte : la variation relative]
L'induction ne dépend pas de \emph{qui} bouge (l'aimant ou la bobine) mais de la \textbf{variation du champ magnétique à travers le circuit}, c'est-à-dire du mouvement \emph{relatif} entre la source de champ et le circuit. C'est ce principe qu'exploitent les alternateurs : dans les centrales de Songloulou et d'Édéa, ce sont des bobines et des aimants (électroaimants) en rotation relative qui produisent la tension du réseau.
\end{aretenir}

\courssub{Expérience 3 : deux bobines face à face}

Supprimons l'aimant : une bobine parcourue par un courant crée elle-même un champ magnétique. Plaçons deux bobines face à face :
\begin{itemize}
\item la bobine $B_1$ est branchée sur une pile avec un interrupteur K et un rhéostat ;
\item la bobine $B_2$ est reliée à un galvanomètre seul, sans générateur.
\end{itemize}

\begin{experience}[Induction entre deux bobines]
\textbf{Observations :}
\begin{itemize}
\item \textbf{À la fermeture de K :} l'aiguille dévie brièvement dans un sens, puis revient à zéro alors que le courant dans $B_1$ est devenu constant.
\item \textbf{K fermé, courant constant :} aucune déviation.
\item \textbf{À l'ouverture de K :} l'aiguille dévie brièvement \emph{dans le sens opposé}.
\item \textbf{K fermé, en déplaçant le curseur du rhéostat :} déviation tant que le courant varie ; le sens de déviation dépend de l'augmentation ou de la diminution du courant.
\end{itemize}
\textbf{Interprétation :} le courant variable dans $B_1$ crée un champ magnétique variable qui traverse $B_2$ : c'est cette variation qui induit la f.é.m. dans $B_2$. Un courant constant, même fort, n'induit rien.
\end{experience}

\begin{popfigure}
\begin{center}
\begin{tikzpicture}[scale=0.95,>=Stealth]
% Bobine B1 (inducteur)
\foreach \y in {0.2,0.5,0.8,1.1}{
  \draw[thick,popblue] (-2.6,\y) ellipse (0.75 and 0.12);
}
\node[popblue] at (-2.6,1.6) {\textbf{Bobine $B_1$ (inducteur)}};
% Circuit inducteur : pile + interrupteur + rhéostat
\draw[thick] (-3.35,0.2) -- (-3.35,-1.2) -- (-1.85,-1.2) -- (-1.85,0.2);
\draw[thick,popdark] (-3.35,-0.75) -- (-3.35,-0.45);
\draw[very thick,popdark] (-3.55,-0.6) -- (-3.15,-0.6);
\draw[thick,popdark] (-3.45,-0.75) -- (-3.25,-0.75);
\node[left] at (-3.45,-0.6) {\small pile};
\draw[thick] (-1.85,-1.2) -- (-1.85,-1.55);
\draw[thick] (-1.85,-1.55) -- (-1.15,-1.25);
\node[right] at (-1.35,-1.6) {\small K};
\draw[thick,decorate,decoration={zigzag,segment length=4,amplitude=1.6}] (-3.35,-1.2) -- (-1.85,-1.2);
\node[below] at (-2.6,-1.35) {\small rhéostat};
% Champ magnétique entre les bobines
\draw[->,thick,popgreen] (-1.6,0.65) -- (-0.6,0.65) node[midway,above]{$\vec{B}$};
\draw[->,thick,popgreen] (-1.6,0.35) -- (-0.6,0.35);
% Bobine B2 (induit)
\foreach \y in {0.2,0.5,0.8,1.1}{
  \draw[thick,poporange] (0.4,\y) ellipse (0.75 and 0.12);
}
\node[poporange] at (0.4,1.6) {\textbf{Bobine $B_2$ (induit)}};
% Circuit induit : galvanomètre
\draw[thick] (-0.35,0.2) -- (-0.35,-1.2) -- (1.15,-1.2) -- (1.15,0.2);
\draw[thick,fill=popcream] (0.4,-1.2) circle (0.45);
\node at (0.4,-1.2) {\small G};
\draw[->,thick,popink] (0.4,-1.2) -- (0.72,-0.92);
\node at (0.4,-1.95) {\small courant induit $i$};
\end{tikzpicture}
\end{center}
\legende{Expérience des deux bobines : $B_1$ (circuit inducteur, avec pile, interrupteur K et rhéostat) crée un champ $\vec{B}$ variable ; $B_2$ (circuit induit, sans générateur) est le siège d'un courant induit détecté par le galvanomètre G.}
\end{popfigure}

\courssub{Expérience 4 : rails de Laplace (expérience historique de Faraday)}

Un barreau conducteur MN peut glisser sur deux rails parallèles reliés à un galvanomètre. L'ensemble est placé dans un champ magnétique $\vec{B}$ uniforme, perpendiculaire au plan des rails (produit par un aimant en U).

\textbf{Observations :} tant que le barreau est immobile, rien ne se passe. Dès qu'on déplace le barreau, un courant apparaît ; son sens dépend du sens de déplacement, et son intensité croît avec la vitesse du barreau.

\textbf{Interprétation :} en se déplaçant, le barreau modifie la \emph{surface} du circuit baignée par le champ $\vec{B}$ : la quantité de champ magnétique « traversant » le circuit varie, d'où l'apparition d'une f.é.m. induite. C'est exactement le dispositif que Faraday étudia en 1831, et le principe des générateurs à courant induit par mouvement mécanique.

\courssub{Vocabulaire fondamental : inducteur, induit, f.é.m. et courant induits}

\begin{definition}[Induction électromagnétique]
On appelle \cle{induction électromagnétique} l'apparition d'une force électromotrice $e$, appelée \cle{f.é.m. induite}, dans un circuit soumis à une \textbf{variation de flux magnétique}. Si le circuit est fermé, cette f.é.m. fait circuler un \cle{courant induit} d'intensité $i$.
\end{definition}

\begin{definition}[Circuit inducteur et circuit induit]
\begin{itemize}
\item Le \cle{circuit inducteur} est le dispositif qui \textbf{crée la variation de champ magnétique} : aimant en mouvement, bobine parcourue par un courant variable.
\item Le \cle{circuit induit} est le circuit \textbf{où naît la f.é.m. induite} : il ne contient aucun générateur, et c'est lui qui est relié au galvanomètre.
\end{itemize}
Dans l'expérience 3, $B_1$ est l'inducteur et $B_2$ l'induit. Dans l'expérience 1, l'aimant joue le rôle d'inducteur et la bobine celui d'induit.
\end{definition}

\begin{methode}[Identifier l'inducteur et l'induit]
\begin{enumerate}
\item Repérer la \textbf{source de champ magnétique variable} : aimant qui bouge, bobine dont le courant change (fermeture, ouverture, rhéostat, courant alternatif). C'est le circuit \textbf{inducteur}.
\item Repérer le circuit \textbf{sans générateur} où apparaît une tension ou un courant (détectés par un galvanomètre, une lampe qui s'allume, etc.). C'est le circuit \textbf{induit}.
\item Vérifier que le champ créé par l'inducteur traverse bien le circuit induit (proximité, orientation).
\end{enumerate}
\end{methode}

\courssub{Constats clés et condition d'apparition}

Les quatre expériences convergent vers les mêmes conclusions, qu'il faut savoir énoncer au Probatoire :

\begin{aretenir}[Trois constats fondamentaux]
\begin{enumerate}
\item \textbf{Pas de variation, pas d'induction :} si le champ magnétique à travers le circuit induit est constant (aimant immobile, courant constant, barreau au repos), la f.é.m. induite est nulle.
\item \textbf{Le sens du courant induit dépend du sens de la variation :} approche ou recul de l'aimant, fermeture ou ouverture de l'interrupteur donnent des courants de sens opposés.
\item \textbf{L'effet est d'autant plus grand que la variation est rapide :} la f.é.m. induite croît avec la vitesse de variation du phénomène.
\end{enumerate}
\end{aretenir}

\textbf{Condition d'apparition (qualitative).} Dans tous les cas étudiés, c'est la \textbf{variation du flux magnétique} à travers le circuit induit qui engendre la f.é.m. induite : variation du champ $\vec{B}$ (expériences 1, 2, 3) ou variation de la surface $S$ du circuit (expérience 4). La grandeur « flux magnétique », qui synthétise $\vec{B}$, $S$ et l'orientation du circuit, sera définie rigoureusement dans la section 2, et la loi de Faraday donnera l'expression quantitative de $e$ dans la section 4.

\begin{exoresolu}[Qui est l'inducteur ? Qui est l'induit ?]
\textbf{Énoncé.} Un élève de Première C à Douala réalise le montage suivant : une bobine $A$ de 500 spires est branchée sur un générateur délivrant un courant dont l'intensité augmente régulièrement de $0$ à \SI{2,0}{A} en \SI{0,50}{s}. Une bobine $B$ de 200 spires, placée à \SI{5}{cm} de $A$ et reliée à un voltmètre, présente à ses bornes une tension de valeur moyenne \SI{12}{mV} pendant cette durée, puis une tension nulle quand le courant dans $A$ est stabilisé.
\begin{enumerate}
\item Identifier le circuit inducteur et le circuit induit. Justifier.
\item Pourquoi la tension aux bornes de $B$ s'annule-t-elle quand le courant dans $A$ est constant ?
\item Que se passerait-il si l'on faisait ensuite décroître le courant de $A$ de \SI{2,0}{A} à $0$ ?
\end{enumerate}
\tcblower
\textbf{Solution détaillée.}
\begin{enumerate}
\item La source de champ magnétique \emph{variable} est la bobine $A$, parcourue par un courant croissant : $A$ est le \textbf{circuit inducteur}. La bobine $B$ ne contient aucun générateur et présente une tension de \SI{12}{mV} : $B$ est le \textbf{circuit induit}, siège de la f.é.m. induite $e = \SI{12}{mV} = \SI{0,012}{V}$.
\item Quand le courant dans $A$ est constant, le champ magnétique qu'il crée est constant : le flux magnétique à travers $B$ ne varie plus, donc la f.é.m. induite est nulle. C'est le premier constat fondamental : \emph{pas de variation, pas d'induction}.
\item Si le courant de $A$ décroît de \SI{2,0}{A} à $0$, le flux à travers $B$ varie à nouveau, mais dans le sens inverse : une f.é.m. induite apparaît pendant la décroissance, de \textbf{signe opposé} à la précédente (le voltmètre indiquerait environ $-\SI{12}{mV}$ si la variation dure aussi \SI{0,50}{s}). Le sens du courant induit dépend du sens de la variation.
\end{enumerate}
\end{exoresolu}

\begin{savaistu}[Faraday, 1831 : la découverte qui éclaire le Cameroun]
C'est le physicien anglais Michael Faraday qui découvrit l'induction électromagnétique le 29 août 1831, avec un anneau de fer entouré de deux bobines : à la fermeture du circuit de la première, un courant bref apparaissait dans la seconde. Parti de rien — il fut apprenti relieur —, Faraday devint l'un des plus grands expérimentateurs de l'histoire. Son nom est donné à l'unité de capacité des condensateurs, le farad. Aujourd'hui, les alternateurs des barrages hydroélectriques de Songloulou et d'Édéa exploitent exactement son principe : la chute d'eau fait tourner des turbines qui entraînent des électroaimants devant des bobines fixes, et la variation de flux ainsi créée produit la tension alternative distribuée par ENEO dans tout le pays.
\end{savaistu}

\begin{exercice}[Pour s'entraîner]
Dans chacune des situations suivantes, dire si une f.é.m. induite apparaît dans la bobine $B$ reliée à un galvanomètre, et préciser le cas échéant le circuit inducteur :
\begin{enumerate}
\item Un aimant est maintenu immobile à l'intérieur de $B$.
\item On éloigne rapidement l'aimant de $B$.
\item Une bobine voisine $A$ est parcourue par un courant continu constant.
\item On ouvre l'interrupteur du circuit de la bobine voisine $A$.
\item $B$ est déplacée dans un champ magnétique uniforme, sans rotation, de sorte que le champ qui la traverse reste le même.
\end{enumerate}
\end{exercice}

\begin{corrige}[Exercice 1]
\begin{enumerate}
\item Aucune f.é.m. induite : le champ est constant (pas de variation de flux).
\item F.é.m. induite pendant le déplacement ; l'aimant en mouvement est l'inducteur, $B$ l'induit.
\item Aucune f.é.m. : courant constant donc champ constant.
\item F.é.m. induite brève à l'ouverture (le champ de $A$ s'annule rapidement) ; $A$ est l'inducteur, $B$ l'induit.
\item Aucune f.é.m. si le champ reste rigoureusement uniforme : le flux à travers $B$ ne varie pas pendant la translation.
\end{enumerate}
\end{corrige}

\coursec{Flux magnétique}

Dans la section précédente, nous avons mis en évidence expérimentalement le phénomène d'induction électromagnétique : un courant apparaît dans un circuit fermé, sans générateur, dès que quelque chose « change » autour de lui — un aimant qui s'approche, une bobine voisine dont le courant varie, une spire qui tourne. Mais ce « quelque chose » restait flou. Pour transformer l'observation en loi quantitative, il nous faut une grandeur physique qui résume, en un seul nombre, « la quantité de champ magnétique qui traverse le circuit ». Cette grandeur existe : c'est le \cle{flux magnétique}. Toute la loi de Faraday, que nous verrons à la section 4, s'exprimera avec lui. Apprenons donc d'abord à le définir et à le calculer proprement.

\courssub{Orientation d'un circuit et vecteur normal}

\textbf{Pourquoi orienter ?}\par
Imagine une rivière qui traverse un filet tendu. Pour dire « combien d'eau passe à travers le filet par seconde », il faut d'abord décider dans quel sens on compte positivement : de la rive gauche vers la rive droite, ou l'inverse. Le débit sera alors positif ou négatif selon ce choix. Pour le flux magnétique, c'est pareil : le flux est une grandeur \emph{algébrique} (elle a un signe), et son signe n'a de sens que si l'on a d'abord \cle{orienté} le circuit.

\begin{definition}[Orientation d'un circuit et vecteur normal]
Orienter un circuit fermé, c'est choisir arbitrairement un \cle{sens positif} de parcours sur son contour. À cette orientation on associe le \cle{vecteur normal} $\vec{n}$ à la surface $S$ s'appuyant sur le circuit, par la \textbf{règle du tire-bouchon} (ou de la main droite) : un tire-bouchon tourné dans le sens positif choisi progresse dans la direction et le sens de $\vec{n}$. Le vecteur $\vec{n}$ est perpendiculaire au plan de la surface, de norme $1$ (vecteur unitaire).
\end{definition}

\begin{attention}[Le choix est arbitraire, mais il faut s'y tenir]
Les deux orientations possibles sont également légitimes : elles donnent des flux de signes opposés, et les résultats physiques finaux (sens réel du courant induit, par exemple) sont les mêmes. En revanche, il est \textbf{interdit de changer d'orientation en cours de calcul} : on choisit un sens positif au début de l'exercice, on trace $\vec{n}$, et on n'y touche plus.
\end{attention}

\courssub{Définition du flux magnétique}

\textbf{Intuition.}\par
Pense à une éolienne face au vent : elle capte un maximum d'air lorsque ses pales sont perpendiculaires au vent, et rien du tout lorsqu'elle est tournée « tranche au vent ». De même, le flux magnétique mesure à quel point le champ $\vec{B}$ « traverse efficacement » la surface. Trois facteurs interviennent :
\begin{itemize}
\item l'intensité $B$ du champ : plus le champ est fort, plus le flux est grand ;
\item l'aire $S$ de la surface : plus la surface est grande, plus elle « intercepte » de champ ;
\item l'inclinaison : seule la composante de $\vec{B}$ perpendiculaire à la surface, c'est-à-dire portée par $\vec{n}$, contribue au flux. Cette composante vaut $B\cos\alpha$, où $\alpha$ est l'angle entre $\vec{B}$ et $\vec{n}$.
\end{itemize}

\begin{definition}[Flux magnétique à travers une surface plane]
Le \cle{flux magnétique} à travers une surface plane d'aire $S$, orientée par son vecteur normal $\vec{n}$, placée dans un champ magnétique \textbf{uniforme} $\vec{B}$, est :
\[ \Phi = \vec{B}\cdot\vec{n}\,S = B\,S\,\cos\alpha \]
où $\alpha$ est l'angle entre le vecteur champ magnétique $\vec{B}$ et la normale $\vec{n}$ à la surface.
\begin{itemize}
\item $B$ en tesla (T) ; $S$ en mètre carré (m$^2$) ; $\Phi$ en weber (Wb).
\end{itemize}
\end{definition}

\begin{popfigure}
\begin{center}
\begin{tikzpicture}[scale=1.05, >=Stealth]
% Surface plane (rectangle en perspective)
\fill[popblueL, opacity=0.55] (0,0) -- (3.4,0) -- (4.3,1.7) -- (0.9,1.7) -- cycle;
\draw[popblue, thick] (0,0) -- (3.4,0) -- (4.3,1.7) -- (0.9,1.7) -- cycle;
\node[popblue] at (3.9,0.45) {$S$};
% Centre
\coordinate (O) at (2.15,0.85);
\fill[popdark] (O) circle (1.6pt);
% Normale
\draw[->, very thick, popdark] (O) -- ++(0,2.5) node[above left] {$\vec{n}$};
% Champ B incliné
\draw[->, very thick, poporange] (O) -- ++(35:2.6) node[right] {$\vec{B}$};
% Projection de B sur n
\draw[->, thick, popgreen, dashed] (O) -- ++(0,{2.6*cos(35)}) node[midway, right=2pt, popgreen] {$B\cos\alpha$};
\draw[dotted, popmuted] ($(O)+(35:2.6)$) -- ($(O)+(0,{2.6*cos(35)})$);
% Angle alpha
\draw[popink, thick] ($(O)+(0,1.1)$) arc (90:35:1.1) node[midway, above right=1pt] {$\alpha$};
% Contour orienté
\draw[->, poppurple, thick] (1.1,0.15) -- (2.2,0.15);
\node[poppurple] at (1.6,-0.28) {\footnotesize sens positif choisi};
\end{tikzpicture}
\end{center}
\legende{Surface plane orientée : la normale $\vec{n}$ est liée au sens positif de parcours par la règle du tire-bouchon. Seule la composante $B\cos\alpha$ de $\vec{B}$ le long de $\vec{n}$ contribue au flux.}
\end{popfigure}

\textbf{Vérification dimensionnelle.} Le tesla vaut $1\ \mathrm{T} = 1\ \mathrm{kg\,s^{-2}\,A^{-1}}$, donc $[B\,S] = \mathrm{kg\,m^2\,s^{-2}\,A^{-1}}$. Nous verrons à la section 4 que le weber s'exprime aussi en volt-seconde ($1\ \mathrm{Wb} = 1\ \mathrm{V\,s}$), ce qui rendra la loi de Faraday homogène : une f.é.m. (en volts) sera une variation de flux (en webers) par unité de temps (en secondes).

\courssub{Unité de flux : le weber}

\begin{aretenir}[Le weber]
L'unité SI de flux magnétique est le \cle{weber} (symbole Wb), en hommage au physicien allemand Wilhelm Weber (1804--1891), collaborateur de Gauss. D'après $\Phi = B S \cos\alpha$ :
\[ 1\ \mathrm{Wb} = 1\ \mathrm{T}\cdot\mathrm{m}^2 \]
\end{aretenir}

\courssub{Cas particuliers à connaître}

\begin{propriete}[Valeurs remarquables du flux]
Pour une surface plane dans un champ uniforme :
\begin{itemize}
\item $\vec{B}$ \textbf{perpendiculaire à la surface} ($\vec{B}$ colinéaire à $\vec{n}$, $\alpha = 0$) : le flux est \textbf{maximal} en valeur absolue :
\[ \Phi = B\,S \quad (\text{ou } -BS \text{ si } \vec{B} \text{ et } \vec{n} \text{ sont opposés, } \alpha = 180^\circ) \]
\item $\vec{B}$ \textbf{parallèle à la surface} ($\vec{B} \perp \vec{n}$, $\alpha = 90^\circ$) : le flux est \textbf{nul} :
\[ \Phi = 0 \]
Le champ « rase » la surface sans la traverser.
\end{itemize}
\end{propriete}

\begin{attention}[Le piège de l'angle : le classique du Probatoire]
L'angle $\alpha$ de la formule est l'angle entre $\vec{B}$ et la \textbf{normale} $\vec{n}$, \emph{pas} l'angle entre $\vec{B}$ et le \textbf{plan} de la spire. Si l'énoncé dit « le champ fait un angle de $30^\circ$ avec le plan de la bobine », alors $\alpha = 90^\circ - 30^\circ = 60^\circ$ et $\Phi = B S \cos 60^\circ$. Confondre les deux angles revient à échanger $\cos\alpha$ et $\sin\alpha$ : erreur garantie. Deuxième piège : convertir l'aire en mètres carrés — $25\ \mathrm{cm}^2 = 25\times 10^{-4}\ \mathrm{m}^2$ (et non $25\times 10^{-2}$ !).
\end{attention}

\courssub{Flux à travers une bobine de $N$ spires}

\textbf{Intuition.}\par
Une bobine plate de $N$ spires identiques, c'est $N$ surfaces superposées traversées par le même champ. Chaque spire « récolte » le flux $B S \cos\alpha$ ; comme les spires sont montées en série, leurs contributions s'ajoutent. Le flux total capté par la bobine est donc $N$ fois plus grand.

\begin{propriete}[Flux à travers une bobine]
Pour une bobine de $N$ spires identiques, chacune d'aire $S$, placées dans un champ uniforme $\vec{B}$ faisant l'angle $\alpha$ avec la normale commune aux spires :
\[ \Phi = N\,B\,S\,\cos\alpha \]
(Admis : conséquence directe de l'additivité du flux sur les $N$ spires.)
\end{propriete}

\begin{methode}[Calculer un flux magnétique en trois étapes]
\begin{enumerate}
\item \textbf{Orienter} le circuit : choisir un sens positif, en déduire et \emph{tracer} le vecteur normal $\vec{n}$ (règle du tire-bouchon).
\item \textbf{Repérer l'angle} $\alpha$ entre $\vec{B}$ et $\vec{n}$ — jamais entre $\vec{B}$ et le plan de la spire. Faire un schéma si nécessaire.
\item \textbf{Appliquer} $\Phi = N B S \cos\alpha$ en unités SI : $B$ en tesla, $S$ en m$^2$ (convertir les cm$^2$ : multiplier par $10^{-4}$). Conclure avec le signe et l'unité : le weber.
\end{enumerate}
\end{methode}

\begin{exemplebox}[Flux à travers une bobine de transformateur]
Une bobine de $N = 200$ spires, chacune d'aire $S = 25\ \mathrm{cm}^2$, est placée dans un champ magnétique uniforme $B = 0{,}40\ \mathrm{T}$ perpendiculaire au plan des spires. Calculons le flux à travers la bobine.
\begin{itemize}
\item Le champ est perpendiculaire au plan des spires, donc parallèle à la normale : $\alpha = 0$ et $\cos\alpha = 1$.
\item Conversion : $S = 25\ \mathrm{cm}^2 = 25\times 10^{-4}\ \mathrm{m}^2 = 2{,}5\times 10^{-3}\ \mathrm{m}^2$.
\item Application : $\Phi = N B S \cos\alpha = 200 \times 0{,}40 \times 2{,}5\times 10^{-3} \times 1$.
\end{itemize}
\[ \boxed{\Phi = 0{,}20\ \mathrm{Wb}} \]
Si l'on tournait la bobine de $90^\circ$ (champ dans le plan des spires), le flux deviendrait nul.
\end{exemplebox}

\begin{exoresolu}[Bobine inclinée dans le champ terrestre]
\textbf{Énoncé.} Une bobine plate de $N = 500$ spires circulaires de rayon $r = 4{,}0\ \mathrm{cm}$ est placée dans un champ magnétique uniforme de valeur $B = 5{,}0\times 10^{-5}\ \mathrm{T}$ (ordre de grandeur du champ magnétique terrestre à Douala). La normale à la bobine fait un angle $\alpha = 60^\circ$ avec $\vec{B}$. Calculer le flux magnétique à travers la bobine.
\tcblower
\textbf{Solution détaillée.}
\begin{enumerate}
\item \textbf{Orientation} : on choisit un sens positif sur la bobine ; la normale $\vec{n}$ est donnée, avec $\alpha = 60^\circ$ entre $\vec{B}$ et $\vec{n}$.
\item \textbf{Aire d'une spire} :
\[ S = \pi r^2 = \pi \times (4{,}0\times 10^{-2})^2 = \pi \times 1{,}6\times 10^{-3} \approx 5{,}03\times 10^{-3}\ \mathrm{m}^2 \]
\item \textbf{Flux} :
\begin{align*}
\Phi &= N B S \cos\alpha = 500 \times 5{,}0\times 10^{-5} \times 5{,}03\times 10^{-3} \times \cos 60^\circ \\
&= 500 \times 5{,}0\times 10^{-5} \times 5{,}03\times 10^{-3} \times 0{,}50 \\
&\approx 6{,}3\times 10^{-5}\ \mathrm{Wb}
\end{align*}
\end{enumerate}
Le flux vaut $\Phi \approx 6{,}3\times 10^{-5}\ \mathrm{Wb}$, soit environ $63\ \mu\mathrm{Wb}$. Remarque : le flux dû au champ terrestre est faible, mais si la bobine tourne rapidement, sa \emph{variation} de flux peut produire une f.é.m. mesurable — c'est le principe des magnétos et des dynamos de vélo.
\end{exoresolu}

\courssub{Signe du flux et origines de sa variation}

Le flux est une grandeur \cle{algébrique} : son signe dépend entièrement de l'orientation choisie. Si $\vec{B}$ et $\vec{n}$ pointent « du même côté » ($\alpha < 90^\circ$), $\Phi > 0$ ; s'ils sont opposés ($\alpha > 90^\circ$), $\Phi < 0$. Changer l'orientation du circuit change le signe de $\Phi$, mais pas la physique.

L'examen de la formule $\Phi = N B S \cos\alpha$ révèle quelque chose de capital pour la suite du chapitre : le flux peut varier de \textbf{trois manières}, correspondant exactement aux trois expériences de la section 1.

\begin{center}
\begin{tabularx}{\linewidth}{|l|X|X|}
\hline
\rowcolor{popblueL} \textbf{Grandeur qui varie} & \textbf{Dispositif expérimental} & \textbf{Exemple concret} \\
\hline
$B$ (le champ) & Courant variable dans la bobine inductrice voisine ; aimant qui s'approche ou s'éloigne & Transformateur d'ENEO : le courant alternatif du primaire crée un champ variable \\
\hline
$S$ (l'aire) & Barreau conducteur glissant sur des rails : l'aire du circuit augmente ou diminue & Freinage électromagnétique, railgun pédagogique \\
\hline
$\alpha$ (l'angle) & Bobine en rotation dans un champ uniforme & Alternateur des barrages de Nachtigal et Lom Pangar, dynamo de vélo \\
\hline
\end{tabularx}
\end{center}

\textbf{Le cas de la bobine en rotation} mérite une étude détaillée, car c'est le principe de \emph{tous} les alternateurs qui produisent l'électricité du Cameroun. Une bobine tourne à vitesse angulaire constante $\omega$ dans un champ $\vec{B}$ uniforme. Si à $t = 0$ la normale est alignée avec $\vec{B}$, alors l'angle évolue selon $\alpha(t) = \omega t$, et le flux devient une fonction sinusoïdale du temps :
\[ \Phi(t) = N B S \cos(\omega t) \]
Le flux oscille entre $+NBS$ (normale alignée avec $\vec{B}$) et $-NBS$ (normale opposée à $\vec{B}$), en passant par zéro deux fois par période (quand le champ rase les spires).

\begin{popfigure}
\begin{center}
\begin{tikzpicture}
\begin{axis}[
  width=0.92\linewidth, height=5.6cm,
  axis lines=middle,
  xlabel={$t$ (s)}, ylabel={$\Phi$ (Wb)},
  xlabel style={at={(ticklabel* cs:1)}, anchor=west},
  ylabel style={at={(ticklabel* cs:1)}, anchor=south},
  xtick={0, 0.5, 1, 1.5, 2},
  xticklabels={$0$, $T/4$, $T/2$, $3T/4$, $T$},
  ytick={0.2, -0.2},
  yticklabels={$+NBS$, $-NBS$},
  ymin=-0.32, ymax=0.32,
  domain=0:2, samples=200,
  grid=both, grid style={popline, very thin},
]
\addplot[popcurve] {0.2*cos(deg(2*pi*x/2))};
\end{axis}
\end{tikzpicture}
\end{center}
\legende{Flux $\Phi(t) = NBS\cos(\omega t)$ à travers une bobine de $N = 200$ spires ($S = 25\ \mathrm{cm}^2$, $B = 0{,}40\ \mathrm{T}$, donc $NBS = 0{,}20\ \mathrm{Wb}$) tournant à vitesse angulaire constante dans un champ uniforme. Le flux est sinusoïdal : c'est sa variation continue qui engendrera la f.é.m. induite (section 4).}
\end{popfigure}

\begin{savaistu}[Du flux tournant à la prise de courant]
Dans l'alternateur du barrage de Nachtigal (420 MW, le plus grand du Cameroun), l'eau du fleuve Sanaga fait tourner une turbine solidaire d'un rotor magnétique : les bobines du stator voient alors un flux $\Phi(t)$ sinusoïdal, et la variation de ce flux y induit une tension alternative — exactement la courbe ci-dessus, à l'échelle industrielle. La fréquence du réseau camerounais, $50\ \mathrm{Hz}$, signifie que le flux dans chaque bobine effectue 50 oscillations complètes par seconde. Toute l'électrotechnique moderne repose sur la formule $\Phi = NBS\cos\alpha$ que tu viens d'apprendre !
\end{savaistu}

\begin{attention}[Flux constant $\neq$ absence de champ]
Un flux nul ne signifie pas qu'il n'y a pas de champ : si $\vec{B}$ est parallèle au plan de la spire, $B$ peut être très grand et $\Phi = 0$. Réciproquement, un flux \emph{constant} (même grand) ne produit \textbf{aucun} phénomène d'induction : seule la \emph{variation} du flux compte, comme le montrera la loi de Faraday. Retiens dès maintenant : « pas de variation de flux, pas de courant induit ».
\end{attention}

\begin{aretenir}[L'essentiel sur le flux magnétique]
\begin{itemize}
\item On \textbf{oriente} le circuit (sens positif arbitraire) ; la normale $\vec{n}$ s'en déduit par la règle du tire-bouchon.
\item Flux à travers une surface plane dans un champ uniforme : $\Phi = B S \cos\alpha$, avec $\alpha$ angle entre $\vec{B}$ et $\vec{n}$.
\item Pour $N$ spires : $\Phi = N B S \cos\alpha$.
\item Unité : le weber, $1\ \mathrm{Wb} = 1\ \mathrm{T\cdot m^2}$.
\item Flux maximal : $\vec{B} \perp$ surface ($\Phi = \pm BS$) ; flux nul : $\vec{B}$ dans le plan de la surface.
\item Le flux peut varier via $B$, via $S$ ou via $\alpha$ : c'est cette variation qui est la cause de l'induction.
\end{itemize}
\end{aretenir}

Maintenant que nous savons quantifier « combien de champ » traverse un circuit, il reste deux questions : dans quel sens circule le courant induit, et quelle est la valeur de la f.é.m. qui le produit ? La loi de Lenz (section 3) répondra à la première, la loi de Faraday (section 4) à la seconde — toutes deux s'exprimeront naturellement avec le flux $\Phi$ et sa variation.

\coursec{Loi de Lenz et sens du courant induit}

Dans la section précédente, nous avons défini le \cle{flux magnétique} $\Phi = B\,S\,\cos\theta$ à travers un circuit, et nous avons constaté expérimentalement qu'un courant apparaît dans une bobine fermée dès que ce flux \emph{varie} : approche ou éloignement d'un aimant, rotation du circuit, variation du courant dans une bobine voisine. Mais une question essentielle reste en suspens : \textbf{dans quel sens circule ce courant induit ?} C'est à cette question que répond la loi de Lenz, véritable boussole de l'électromagnétisme.

\courssub{Énoncé de la loi de Lenz}

\textbf{Une observation simple pour démarrer.} Approche le pôle Nord d'un aimant d'une bobine reliée à un galvanomètre : l'aiguille dévie. Éloigne maintenant le même aimant : l'aiguille dévie \emph{dans l'autre sens}. Le courant induit change donc de sens selon que le flux augmente ou diminue. Ce n'est pas un hasard : la nature obéit à une règle remarquablement simple, formulée en 1834 par le physicien russe Heinrich Lenz.

\begin{propriete}[Loi de Lenz (admise)]
Le sens du courant induit dans un circuit fermé est tel que, \textbf{par ses effets, il s'oppose à la cause qui lui a donné naissance}, c'est-à-dire à la \emph{variation} du flux magnétique à travers le circuit.

Concrètement, le courant induit crée son propre champ magnétique, appelé \cle{champ induit} $\vec{B}_i$, dont le sens est choisi de sorte que :
\begin{itemize}
\item si le flux inducteur \textbf{augmente}, $\vec{B}_i$ est de sens \textbf{opposé} au champ inducteur $\vec{B}$ (il le \emph{contrecarre}) ;
\item si le flux inducteur \textbf{diminue}, $\vec{B}_i$ est de \textbf{même sens} que $\vec{B}$ (il le \emph{renforce}).
\end{itemize}
Cette loi est admise au niveau Première, mais nous la justifierons énergétiquement à la fin de cette section.
\end{propriete}

\begin{attention}[Le piège classique : opposition à la variation, pas au champ]
L'erreur la plus fréquente au Probatoire est d'écrire : « le champ induit s'oppose toujours au champ inducteur ». C'est \textbf{faux}. Le champ induit $\vec{B}_i$ s'oppose à la \textbf{variation} du flux, non au champ lui-même :
\begin{itemize}
\item $B$ \emph{augmente} $\Rightarrow$ $\vec{B}_i$ opposé à $\vec{B}$ ;
\item $B$ \emph{diminue} $\Rightarrow$ $\vec{B}_i$ de même sens que $\vec{B}$ (il « freine » la diminution).
\end{itemize}
Retiens l'image : le circuit est un « conservateur » qui déteste le changement. Il lutte contre toute variation de flux, dans un sens comme dans l'autre.
\end{attention}

\courssub{Application à l'expérience aimant--bobine}

Reprenons l'expérience de la section 1 et raisonnons avec la loi de Lenz. Rappelons qu'une bobine parcourue par un courant se comporte comme un aimant : la face par laquelle le courant est vu tourner dans le \textbf{sens antihoraire} est une \textbf{face Nord} ; si le courant est vu tourner dans le \textbf{sens horaire}, c'est une \textbf{face Sud}.

\textbf{Premier cas : on approche le pôle Nord de l'aimant.} Le champ $\vec{B}$ de l'aimant, dirigé vers la bobine (un champ sort du pôle Nord), voit son intensité \emph{augmenter} à travers les spires. D'après Lenz, le champ induit $\vec{B}_i$ doit s'opposer à cette augmentation : il est donc dirigé \emph{vers l'aimant}, c'est-à-dire qu'il « sort » de la face proche de la bobine. Cette face devient un \textbf{pôle Nord}. Or deux pôles Nord se repoussent : la bobine \emph{repousse} l'aimant qu'on cherche à approcher. Magnifique cohérence : le circuit s'oppose bien à la cause (l'approche) qui le fait réagir.

\textbf{Deuxième cas : on éloigne le pôle Nord.} Le flux diminue. Le champ induit $\vec{B}_i$ prend le même sens que $\vec{B}$ pour compenser la diminution : la face proche de la bobine devient un \textbf{pôle Sud}, qui \emph{attire} l'aimant qu'on cherche à éloigner. Là encore, la bobine s'oppose à la cause (l'éloignement).

\begin{popfigure}
\begin{tikzpicture}[scale=0.92, every node/.style={scale=0.92}]
% ================= PANNEAU 1 : APPROCHE =================
\begin{scope}
% Aimant
\fill[popblue!70] (0,0.5) rectangle (1.2,1.5);
\fill[poporange!80] (1.2,0.5) rectangle (2.4,1.5);
\node[white] at (0.6,1) {\textbf{S}};
\node[white] at (1.8,1) {\textbf{N}};
% Flèche mouvement
\draw[->, very thick, popdark] (2.5,2.1) -- (3.7,2.1) node[midway, above]{approche $\vec{v}$};
% Bobine (spires)
\foreach \y in {0.4,0.8,1.2,1.6} {
  \draw[thick, popink] (4.6,\y) ellipse (0.28 and 0.34);
}
\draw[thick, popink] (4.6,1.94) -- (4.6,2.3) -- (6.3,2.3);
\draw[thick, popink] (4.6,0.06) -- (4.6,-0.3) -- (6.3,-0.3);
\draw[thick, popink] (6.3,2.3) -- (6.3,1.35);
\draw[thick, popink] (6.3,0.65) -- (6.3,-0.3);
\node[popink] at (6.3,1.0) {\footnotesize G};
\draw[thick, popink] (6.3,1.35) circle (0.35);
\draw[thick, popink] (6.3,0.65) circle (0.35);
% Champ inducteur B (vers la droite, sort du N)
\draw[->, very thick, popblue] (2.6,0.2) -- (4.2,0.2) node[midway, below]{$\vec{B}$ (croît)};
% Champ induit Bi (vers la gauche)
\draw[->, very thick, poporange] (4.3,1.85) -- (2.9,1.85) node[midway, above left=-1pt]{$\vec{B}_i$};
% Face Nord de la bobine
\node[poporange, font=\bfseries] at (5.35,1.0) {face N};
% Courant induit sur une spire (vu depuis l'aimant : antihoraire)
\draw[->, very thick, popgreen] (4.6,1.54) arc (90:200:0.28 and 0.34);
\node[popgreen] at (3.9,2.6) {$i$ antihoraire};
\node[popgreen] at (3.9,2.35) {\footnotesize (vu depuis l'aimant)};
\node at (3.4,-0.9) {\textbf{a) Le pôle Nord s'approche : répulsion}};
\end{scope}
% ================= PANNEAU 2 : ÉLOIGNEMENT =================
\begin{scope}[xshift=9.2cm]
\fill[popblue!70] (0,0.5) rectangle (1.2,1.5);
\fill[poporange!80] (1.2,0.5) rectangle (2.4,1.5);
\node[white] at (0.6,1) {\textbf{S}};
\node[white] at (1.8,1) {\textbf{N}};
\draw[->, very thick, popdark] (3.7,2.1) -- (2.5,2.1) node[midway, above]{éloignement $\vec{v}$};
\foreach \y in {0.4,0.8,1.2,1.6} {
  \draw[thick, popink] (4.6,\y) ellipse (0.28 and 0.34);
}
\draw[thick, popink] (4.6,1.94) -- (4.6,2.3) -- (6.3,2.3);
\draw[thick, popink] (4.6,0.06) -- (4.6,-0.3) -- (6.3,-0.3);
\draw[thick, popink] (6.3,2.3) -- (6.3,1.35);
\draw[thick, popink] (6.3,0.65) -- (6.3,-0.3);
\node[popink] at (6.3,1.0) {\footnotesize G};
\draw[thick, popink] (6.3,1.35) circle (0.35);
\draw[thick, popink] (6.3,0.65) circle (0.35);
% Champ inducteur B (vers la droite, décroît)
\draw[->, very thick, popblue] (2.6,0.2) -- (4.2,0.2) node[midway, below]{$\vec{B}$ (décroît)};
% Champ induit Bi (même sens que B)
\draw[->, very thick, poporange] (2.9,1.85) -- (4.3,1.85) node[midway, above]{$\vec{B}_i$};
\node[popblue, font=\bfseries] at (5.35,1.0) {face S};
\draw[->, very thick, popgreen] (4.6,0.46) arc (-90:20:0.28 and 0.34);
\node[popgreen] at (3.9,2.6) {$i$ horaire};
\node[popgreen] at (3.9,2.35) {\footnotesize (vu depuis l'aimant)};
\node at (3.4,-0.9) {\textbf{b) Le pôle Nord s'éloigne : attraction}};
\end{scope}
\end{tikzpicture}
\legende{Loi de Lenz appliquée à l'expérience aimant--bobine. a) L'approche du pôle Nord fait apparaître une face Nord sur la bobine (répulsion) : le courant induit est antihoraire vu depuis l'aimant. b) L'éloignement fait apparaître une face Sud (attraction) : le courant est horaire. Dans les deux cas, l'effet s'oppose à la cause.}
\end{popfigure}

\courssub{Détermination pratique du sens du courant induit}

Passons maintenant à la méthode systématique, celle qui te permettra de traiter n'importe quelle situation d'examen sans hésiter.

\begin{methode}[Trouver le sens du courant induit en 3 étapes]
\begin{enumerate}
\item \textbf{Analyser la cause.} Identifier le sens du champ inducteur $\vec{B}$ à travers le circuit et déterminer si le flux $\Phi$ \emph{augmente} ou \emph{diminue} (variation de $B$, de $S$ ou de l'angle $\theta$).
\item \textbf{Appliquer la loi de Lenz.} En déduire le sens du champ induit $\vec{B}_i$ : opposé à $\vec{B}$ si $\Phi$ augmente, de même sens que $\vec{B}$ si $\Phi$ diminue.
\item \textbf{Utiliser la règle du tire-bouchon (ou de la main droite).} Placer le tire-bouchon de sorte qu'il avance dans le sens de $\vec{B}_i$ : son sens de rotation donne le sens du courant induit $i$ dans les spires. Vérifier éventuellement avec la règle des faces (face Nord = courant antihoraire vu de face).
\end{enumerate}
\end{methode}

\begin{experience}[Vérification au galvanomètre]
\textbf{Matériel :} une bobine de 500 spires, un aimant droit, un galvanomètre à zéro central, des fils de connexion.

\textbf{Protocole :} relier la bobine au galvanomètre. Approcher lentement le pôle Nord de l'aimant d'une face de la bobine, noter le sens de déviation de l'aiguille. Éloigner ensuite l'aimant, puis recommencer plus vite.

\textbf{Observations :}
\begin{itemize}
\item à l'approche, l'aiguille dévie d'un côté (par exemple vers la droite) ;
\item à l'éloignement, elle dévie de l'autre côté (vers la gauche) ;
\item plus le mouvement est rapide, plus la déviation est grande ;
\item aimant immobile : aucune déviation, même si l'aimant est très proche.
\end{itemize}

\textbf{Interprétation :} le changement de sens de déviation confirme que le courant induit change de sens selon que le flux augmente ou diminue, exactement comme le prévoit la loi de Lenz. L'absence de courant à flux constant confirme que c'est la \emph{variation} de $\Phi$ qui est la cause du phénomène, et non la simple présence du champ.
\end{experience}

\begin{exemplebox}[Pôle Sud qui s'éloigne d'une bobine]
Un aimant droit présente son \textbf{pôle Sud} face à une bobine fermée. On \textbf{éloigne} l'aimant. Quel est le sens du courant induit vu depuis l'aimant ?

\textit{Résolution avec la méthode en 3 étapes :}
\begin{enumerate}
\item \textbf{Cause.} Le champ d'un aimant \emph{entre} par son pôle Sud : à travers la bobine, $\vec{B}$ est donc dirigé \emph{vers l'aimant}. En s'éloignant, l'aimant fait \emph{diminuer} le flux.
\item \textbf{Loi de Lenz.} Le flux diminue, donc $\vec{B}_i$ est de \emph{même sens} que $\vec{B}$ : dirigé vers l'aimant, c'est-à-dire \emph{entrant} dans la face proche de la bobine. Un champ qui entre par une face en fait une \textbf{face Sud} : la bobine attire l'aimant qui s'éloigne (opposition à la cause, conforme à Lenz).
\item \textbf{Tire-bouchon.} Pour que le tire-bouchon avance vers l'aimant (sens de $\vec{B}_i$), il doit tourner dans le sens \textbf{horaire} vu depuis l'aimant. Le courant induit est donc \textbf{horaire} vu depuis l'aimant.
\end{enumerate}
\textit{Vérification par la règle des faces :} face Sud $\Leftrightarrow$ courant horaire vu de face. Cohérent.
\end{exemplebox}

\courssub{Le rail de Laplace : une autre illustration de la loi de Lenz}

Considérons une barre métallique $MN$ pouvant glisser sans frottement sur deux rails parallèles reliés par un conducteur ohmique, le tout plongé dans un champ magnétique uniforme $\vec{B}$ perpendiculaire au plan des rails, \emph{entrant} (noté $\otimes$). On tire la barre vers la droite à la vitesse $\vec{v}$.

\textbf{Analyse.} Quand la barre avance, l'aire $S$ du circuit augmente, donc le flux $\Phi = B\,S$ (avec $\theta = 0$) \emph{augmente}. D'après Lenz, le champ induit $\vec{B}_i$ doit être \emph{sortant} ($\odot$) pour s'opposer à cette augmentation. La règle du tire-bouchon donne alors un courant induit circulant dans le \textbf{sens antihoraire} dans le circuit, c'est-à-dire de $N$ vers $M$ dans la barre.

\textbf{Force de Laplace.} La barre, parcourue par le courant $i$ dans le champ $\vec{B}$, subit la force de Laplace $\vec{F} = i\,\vec{\ell}\wedge\vec{B}$, où $\vec{\ell}$ est le vecteur le long de la barre dans le sens du courant. Le produit vectoriel montre que $\vec{F}$ est dirigée \emph{vers la gauche} : elle s'oppose exactement au mouvement de la barre. C'est une nouvelle confirmation de la loi de Lenz : l'effet (la force) s'oppose à la cause (le déplacement).

\begin{popfigure}
\begin{tikzpicture}[scale=1.0]
% Champ B entrant
\foreach \x in {0.5,1.5,2.5,3.5,4.5,5.5} {
  \foreach \y in {0.5,1.5,2.5} {
    \draw[popblue] (\x,\y) node {$\otimes$};
  }
}
\node[popblue] at (5.9,2.9) {$\vec{B}$ entrant};
% Rails
\draw[very thick, popdark] (0,0.3) -- (6.2,0.3);
\draw[very thick, popdark] (0,2.7) -- (6.2,2.7);
% Résistor à gauche
\draw[thick, popink] (0,0.3) -- (-0.6,0.3) -- (-0.6,1.1);
\draw[thick, popink, decorate, decoration={zigzag, segment length=4pt, amplitude=2pt}] (-0.6,1.1) -- (-0.6,1.9);
\draw[thick, popink] (-0.6,1.9) -- (-0.6,2.7) -- (0,2.7);
\node[popink] at (-1.0,1.5) {$R$};
% Barre mobile
\draw[line width=3pt, poporange] (3.2,0.2) -- (3.2,2.8);
\node[poporange, above] at (3.2,2.85) {$M$};
\node[poporange, below] at (3.2,0.15) {$N$};
% Vitesse
\draw[->, very thick, popdark] (3.2,3.3) -- (4.6,3.3) node[midway, above]{$\vec{v}$};
% Force de Laplace
\draw[->, very thick, popgreen] (3.2,1.5) -- (1.9,1.5) node[midway, above]{$\vec{F}$ (Laplace)};
% Courant induit dans la barre (de N vers M)
\draw[->, very thick, poppurple] (3.55,0.7) -- (3.55,2.3) node[midway, right]{$i$};
% Sens du courant sur le rail du haut (vers la gauche)
\draw[->, very thick, poppurple] (2.6,2.7) -- (1.4,2.7);
% Champ induit sortant au centre
\node[poporange] at (1.6,1.5) {$\odot$};
\node[poporange] at (1.6,1.15) {\footnotesize $\vec{B}_i$ sortant};
\end{tikzpicture}
\legende{Barre $MN$ glissant vers la droite sur des rails dans un champ $\vec{B}$ entrant. L'aire du circuit augmente, le flux augmente : le champ induit $\vec{B}_i$ est sortant, le courant $i$ circule de $N$ vers $M$ (sens antihoraire), et la force de Laplace $\vec{F} = i\,\vec{\ell}\wedge\vec{B}$ s'oppose au déplacement.}
\end{popfigure}

\courssub{Justification énergétique de la loi de Lenz}

Pourquoi la nature « s'oppose-t-elle » ainsi ? La réponse est profonde : \textbf{la loi de Lenz est une conséquence directe de la conservation de l'énergie}.

Raisonnons par l'absurde sur l'expérience aimant--bobine. Supposons qu'en approchant le pôle Nord, la face proche de la bobine devienne un pôle \emph{Sud} (contraire de Lenz) : la bobine attirerait l'aimant, qui accélérerait tout seul vers elle, le courant induit augmenterait, l'attraction croîtrait encore... On produirait à la fois de l'énergie cinétique et de l'énergie électrique \emph{sans rien fournir} : un mouvement perpétuel, impossible.

Avec la loi de Lenz, au contraire, la bobine \emph{repousse} l'aimant qu'on approche. Pour maintenir le mouvement, l'opérateur doit exercer une force contre cette répulsion et fournir un \textbf{travail mécanique}. Ce travail est précisément la source de l'énergie électrique apparue dans le circuit :
\[
W_{\text{mécanique fourni}} \;=\; E_{\text{électrique induite}} \;+\; E_{\text{dissipée par effet Joule}}.
\]

De même, sur les rails de Laplace, la force de Laplace $\vec{F}$ s'oppose à $\vec{v}$ : sa puissance $P = \vec{F}\cdot\vec{v} = -F\,v$ est négative (force résistante). Pour maintenir la barre à vitesse constante, l'opérateur doit fournir la puissance mécanique $+F\,v$, que le circuit convertit en puissance électrique $e\,i$, puis en chaleur dans le résistor $R$ (effet Joule). Tout se tient : rien ne se perd, rien ne se crée.

\begin{aretenir}[L'essentiel de la loi de Lenz]
\begin{itemize}
\item Le courant induit s'oppose, \textbf{par ses effets}, à la \textbf{variation de flux} qui le produit.
\item Flux qui augmente $\Rightarrow$ $\vec{B}_i$ opposé à $\vec{B}$ ; flux qui diminue $\Rightarrow$ $\vec{B}_i$ de même sens que $\vec{B}$.
\item Méthode : variation de $\Phi$ $\rightarrow$ sens de $\vec{B}_i$ $\rightarrow$ sens de $i$ par la règle du tire-bouchon.
\item Aimant--bobine : approche d'un pôle $\Rightarrow$ face de même nom (répulsion) ; éloignement $\Rightarrow$ face de nom contraire (attraction).
\item La loi de Lenz traduit la \textbf{conservation de l'énergie} : l'énergie électrique induite provient d'un travail mécanique fourni contre les forces électromagnétiques.
\end{itemize}
\end{aretenir}

\begin{exoresolu}[Sens du courant induit dans un cadre]
Un cadre rectangulaire fermé est placé dans un champ magnétique uniforme $\vec{B}$ perpendiculaire à son plan, dirigé \emph{vers l'avant} (sortant, $\odot$). On fait \emph{décroître} progressivement la valeur de $B$ de \SI{0,80}{\tesla} à \SI{0,20}{\tesla}.

1. Le flux à travers le cadre augmente-t-il ou diminue-t-il ?
2. Déterminer le sens du champ induit $\vec{B}_i$.
3. En déduire le sens du courant induit dans le cadre, vu de face.
\tcblower
\textbf{1. Variation du flux.} Le flux vaut $\Phi = B\,S\,\cos\theta$ avec $S$ et $\theta$ constants. Comme $B$ passe de \SI{0,80}{\tesla} à \SI{0,20}{\tesla}, le flux \textbf{diminue}.

\textbf{2. Sens du champ induit.} D'après la loi de Lenz, le circuit s'oppose à cette diminution : le champ induit $\vec{B}_i$ doit \emph{renforcer} le champ inducteur, donc être de \textbf{même sens} que $\vec{B}$, c'est-à-dire \textbf{sortant} ($\odot$). Attention au piège : ici $\vec{B}_i$ n'est \emph{pas} opposé à $\vec{B}$, car c'est la \emph{variation} qui compte.

\textbf{3. Sens du courant.} Règle du tire-bouchon : pour avancer dans le sens sortant (vers l'observateur), le tire-bouchon tourne dans le sens \textbf{antihoraire}. Le courant induit circule donc dans le \textbf{sens antihoraire} vu de face. Vérification : un courant antihoraire vu de face fait de cette face une face Nord, dont le champ sort bien vers l'observateur. Cohérent.
\end{exoresolu}

\begin{savaistu}[Le freinage électromagnétique : la loi de Lenz au service de la sécurité]
La loi de Lenz est à l'origine du \textbf{freinage électromagnétique}, utilisé dans les TGV, certains camions et bus, et les télésièges. Quand un disque métallique solidaire d'une roue tourne devant des électroaimants, des courants induits — appelés \textbf{courants de Foucault} — naissent dans le métal. D'après la loi de Lenz, les forces électromagnétiques associées s'opposent à la rotation du disque : le véhicule ralentit \emph{sans aucun contact mécanique}, donc sans usure de plaquettes ! L'énergie cinétique du véhicule est convertie en chaleur dans le disque, en parfaite illustration du bilan énergétique que nous avons établi. Au Cameroun, le même principe explique pourquoi il faut « forcer » pour faire tourner la manivelle d'une dynamo de lampe de poche rechargeable : plus la lampe consomme, plus la résistance électromagnétique est grande — c'est la loi de Lenz que ton bras combat !
\end{savaistu}

Maintenant que nous savons \emph{dans quel sens} circule le courant induit grâce à la loi de Lenz, il reste à répondre à une question quantitative : \emph{quelle est la valeur} de la force électromotrice induite ? C'est l'objet de la section suivante, consacrée à la loi de Faraday.

\coursec{F.é.m. induite : loi de Faraday}

La loi de Lenz, étudiée dans la section précédente, nous a appris à \emph{prévoir le sens} du courant induit : celui-ci s'oppose toujours, par ses effets, à la variation de flux qui lui a donné naissance. Mais un électricien qui dimensionne un alternateur de barrage (comme ceux de Nachtigal ou de Lom Pangar) a besoin de plus qu'un sens : il lui faut une \emph{valeur}. Quelle est la \cle{force électromotrice induite}, en volts, lorsque le flux varie d'une certaine manière ? C'est toute l'ambition de la \cle{loi de Faraday}, que nous allons énoncer, exploiter et illustrer dans cette section.

\courssub{Énoncé de la loi de Faraday}

\textbf{Intuition.} Reprenons l'expérience de la section 1 : un aimant approché d'une bobine reliée à un galvanomètre. On observe un fait capital : plus l'aimant est déplacé \emph{rapidement}, plus la déviation du galvanomètre est \emph{grande}. Si l'aimant est immobile, la déviation est nulle, quel que soit le flux. Ce n'est donc pas le flux $\Phi$ qui crée la f.é.m., mais sa \emph{vitesse de variation} : la f.é.m. induite est d'autant plus grande que le flux varie vite. C'est exactement ce que traduit la dérivée $\dfrac{d\Phi}{dt}$.

\begin{propriete}[Loi de Faraday (admise)]
Soit un circuit fermé orienté, traversé par un flux magnétique $\Phi$ variable au cours du temps. Il apparaît dans ce circuit une \cle{force électromotrice induite} :
\begin{itemize}
\item \textbf{f.é.m. moyenne} sur une durée $\Delta t$ :
\[ e_{\text{moy}} = -\frac{\Delta \Phi}{\Delta t} \]
\item \textbf{f.é.m. instantanée} :
\[ e = -\frac{d\Phi}{dt} \]
\end{itemize}
où $\Phi$ est le flux à travers le circuit \emph{orienté} (le signe de $\Phi$ dépend de l'orientation choisie de la normale $\vec{n}$, liée au sens positif du courant par la règle du tire-bouchon). L'unité de $e$ est le \textbf{volt} (V). Cette loi, établie expérimentalement par Michael Faraday en 1831, est \emph{admise} en Première : aucune démonstration n'est exigible au Probatoire.
\end{propriete}

\begin{aretenir}[Le signe « $-$ » : la mémoire de la loi de Lenz]
Le signe moins dans $e = -\dfrac{d\Phi}{dt}$ n'est pas un détail : il \emph{traduit mathématiquement la loi de Lenz}. Si le flux (algébrique) augmente ($\dfrac{d\Phi}{dt} > 0$), alors $e < 0$ : la f.é.m. tend à faire circuler un courant dans le sens \emph{négatif} de l'orientation choisie, c'est-à-dire le sens qui crée un flux opposé à l'augmentation. Si le flux diminue, $e > 0$ et le courant induit renforce le flux. En pratique, on calcule souvent la \emph{valeur absolue} $|e| = \left|\dfrac{\Delta \Phi}{\Delta t}\right|$, puis on détermine le sens du courant séparément par la loi de Lenz.
\end{aretenir}

\textbf{Vérification d'homogénéité.} Le flux s'exprime en webers : $1~\text{Wb} = 1~\text{T} \cdot \text{m}^2$. La loi de Faraday donne :
\[ [e] = \frac{[\Phi]}{[t]} = \frac{\text{Wb}}{\text{s}} \]
Or $1~\text{T} = 1~\text{V} \cdot \text{s} \cdot \text{m}^{-2}$ (conséquence de la force de Laplace $F = B I \ell$ : $[B] = \dfrac{\text{N}}{\text{A} \cdot \text{m}} = \dfrac{\text{V} \cdot \text{s}}{\text{m}^2}$). Donc :
\[ \frac{1~\text{Wb}}{1~\text{s}} = \frac{1~\text{V} \cdot \text{s} \cdot \text{m}^{-2} \times 1~\text{m}^2}{1~\text{s}} = 1~\text{V} \]
La loi est homogène : \textbf{un weber par seconde, c'est un volt}.

\begin{attention}[Flux ou variation de flux ?]
Deux erreurs reviennent chaque année au Probatoire :
\begin{enumerate}
\item \textbf{Confondre $\Phi$ et $\Delta \Phi$.} Un flux \emph{grand mais constant} ne produit \emph{aucune} f.é.m. Seule la \emph{variation} de flux compte. Un circuit plongé dans un champ intense mais permanent n'est le siège d'aucune induction.
\item \textbf{Oublier le nombre de spires $N$.} Pour une bobine de $N$ spires, le flux \emph{total} (ou flux embrassé) est $\Phi_{\text{total}} = N \times \Phi_{\text{une spire}}$. C'est ce flux total qu'il faut dériver. Oublier $N$, c'est sous-estimer la f.é.m. d'un facteur $N$ — parfois plusieurs centaines !
\end{enumerate}
On ajoutera une troisième vigilance : \textbf{convertir les unités} avant tout calcul — les surfaces sont souvent données en cm$^2$ ($1~\text{cm}^2 = 10^{-4}~\text{m}^2$) et les flux en mWb ($1~\text{mWb} = 10^{-3}$ Wb).
\end{attention}

\courssub{Lecture graphique : flux affine, f.é.m. constante}

Si le flux $\Phi(t)$ est une fonction \emph{affine} du temps (droite sur le graphique), sa dérivée est constante : la f.é.m. induite est alors constante, égale à l'opposée de la pente de la droite. C'est la situation la plus simple, et elle est très souvent exploitée dans les exercices.

\begin{popfigure}
\begin{center}
\begin{tikzpicture}
\begin{axis}[
  width=0.48\linewidth, height=5cm,
  xlabel={$t$ (s)}, ylabel={$\Phi$ (Wb)},
  xmin=0, xmax=0.5, ymin=0, ymax=0.35,
  xtick={0,0.1,0.2,0.3,0.4}, ytick={0,0.1,0.2,0.3},
  grid=both, grid style={popline!40},
  title={Flux $\Phi(t)$},
  title style={font=\small}
]
\addplot[popcurve, domain=0:0.4, samples=50] {0.3 - 0.5*x};
\addplot[mark=*, poporange, only marks] coordinates {(0,0.3) (0.4,0.1)};
\node[font=\scriptsize, popdark] at (axis cs:0.24,0.25) {pente $= -0{,}5$ Wb/s};
\end{axis}
\end{tikzpicture}\hfill
\begin{tikzpicture}
\begin{axis}[
  width=0.48\linewidth, height=5cm,
  xlabel={$t$ (s)}, ylabel={$e$ (V)},
  xmin=0, xmax=0.5, ymin=-0.2, ymax=0.8,
  xtick={0,0.1,0.2,0.3,0.4}, ytick={0,0.5},
  grid=both, grid style={popline!40},
  title={F.é.m. $e(t) = -\dfrac{d\Phi}{dt}$},
  title style={font=\small}
]
\addplot[popcurve, domain=0:0.4, samples=2] {0.5};
\addplot[mark=*, poporange, only marks] coordinates {(0,0.5) (0.4,0.5)};
\end{axis}
\end{tikzpicture}
\end{center}
\legende{Le flux décroît linéairement de $0{,}30$ Wb à $0{,}10$ Wb en $0{,}40$ s (pente $-0{,}50$ Wb/s) : la f.é.m. induite est constante, $e = +0{,}50$ V. Le signe positif traduit la loi de Lenz : la f.é.m. s'oppose à la diminution du flux.}
\end{popfigure}

\begin{exoresolu}[Variation de flux dans une bobine de 500 spires]
Le flux à travers \emph{chaque spire} d'une bobine de $N = 500$ spires passe de $\Phi_1 = \SI{0,30}{mWb} = \num{3,0e-4}$ Wb à $\Phi_2 = \SI{0,10}{mWb} = \num{1,0e-4}$ Wb en $\Delta t = \SI{0,20}{s}$, de façon régulière. Calculer la valeur absolue de la f.é.m. induite moyenne aux bornes de la bobine.
\tcblower
\textbf{Solution détaillée.}

Étape 1 : flux total embrassé par la bobine. Le flux total est la somme des flux à travers les $N$ spires :
\[ \Phi_{\text{total}} = N \, \Phi_{\text{spire}} \]
Étape 2 : variation de flux total :
\[ \Delta \Phi_{\text{total}} = N \, (\Phi_2 - \Phi_1) = 500 \times (\num{1,0e-4} - \num{3,0e-4}) = 500 \times (-\num{2,0e-4}) = -\SI{0,10}{Wb} \]
Étape 3 : loi de Faraday en valeur absolue :
\[ |e| = \left| \frac{\Delta \Phi_{\text{total}}}{\Delta t} \right| = \frac{0{,}10}{0{,}20} = \SI{0,50}{V} \]
\textbf{Vérification de cohérence} : la variation de flux par spire est $\num{2,0e-4}$ Wb en $0{,}20$ s, soit $10^{-3}$ Wb/s par spire ; multipliée par $N = 500$ spires, on retrouve bien $0{,}50$ V. Sans le facteur $N$, on aurait trouvé $1{,}0$ mV — erreur d'un facteur 500 !

Le signe de $e$ dépend de l'orientation choisie ; le sens réel du courant induit s'obtient par la loi de Lenz : le courant crée un champ qui \emph{renforce} le flux décroissant.
\end{exoresolu}

\courssub{Cas général d'une bobine : les trois façons de faire varier le flux}

Pour une bobine plane de $N$ spires, d'aire $S$ chacune, dont la normale fait l'angle $\alpha$ avec le champ $\vec{B}$ supposé uniforme, le flux total s'écrit :
\[ \Phi = N \, B \, S \cos \alpha \]
La f.é.m. induite est donc :
\[ e = -\frac{d\Phi}{dt} = -N \, \frac{d}{dt}\left( B \, S \cos \alpha \right) \]
Faire varier le flux revient à faire varier \emph{l'une au moins} des trois grandeurs $B$, $S$ ou $\alpha$. Distinguons les trois cas types.

\textbf{Cas 1 : variation du champ $B$ (bobine fixe).} Si $B = a\,t$ (champ croissant linéairement, $a > 0$), avec $S$ et $\alpha$ constants :
\[ e = -N \, S \cos \alpha \, \frac{dB}{dt} = -N \, S \, a \cos \alpha \]
La f.é.m. est \emph{constante}. C'est le principe des expériences d'induction par électroaimant alimenté par un courant variable.

\textbf{Cas 2 : variation de la surface $S$ (barreau sur rails).} C'est le cas historique et le plus formateur : il fait l'objet de la démonstration guidée ci-dessous.

\textbf{Cas 3 : variation de l'angle $\alpha$ (bobine en rotation).} Si la bobine tourne à vitesse angulaire constante $\omega$ dans un champ $\vec{B}$ uniforme, alors $\alpha = \omega t$ (en choisissant $\alpha = 0$ à $t = 0$) et :
\[ \Phi = N B S \cos(\omega t) \quad \Longrightarrow \quad e = -\frac{d\Phi}{dt} = N B S \, \omega \sin(\omega t) \]
La f.é.m. est \emph{sinusoïdale}, d'amplitude $e_{\max} = N B S \omega$ : c'est exactement le \cle{principe de l'alternateur}. Les turbo-alternateurs des barrages d'Edéa, de Song Loulou ou de Nachtigal ne font rien d'autre : l'eau fait tourner des bobines (ou des aimants) et la rotation périodique de $\alpha$ produit la tension alternative de $50$ Hz distribuée par ENEO.

\begin{experience}[Démonstration guidée : la f.é.m. d'un barreau sur rails (à reconstruire pas à pas)]
Un barreau conducteur MN, de longueur $\ell$, glisse sans frottement à vitesse constante $\vec{v}$ sur deux rails parallèles reliés par un conducteur ohmique. L'ensemble est plongé dans un champ $\vec{B}$ uniforme, perpendiculaire au plan des rails. Établissons l'expression $|e| = B \, \ell \, v$.

\textbf{Étape 1 — Surface balayée.} Pendant la durée $\Delta t$, le barreau se déplace de la distance :
\[ \Delta x = v \, \Delta t \]
La surface du circuit augmente donc de :
\[ \Delta S = \ell \, \Delta x = \ell \, v \, \Delta t \]

\textbf{Étape 2 — Variation de flux.} Le champ $B$ étant uniforme et perpendiculaire au plan du circuit ($\cos \alpha = 1$ en valeur absolue) :
\[ |\Delta \Phi| = B \, |\Delta S| = B \, \ell \, v \, \Delta t \]

\textbf{Étape 3 — Loi de Faraday.}
\[ |e| = \frac{|\Delta \Phi|}{\Delta t} = \frac{B \, \ell \, v \, \Delta t}{\Delta t} \]
soit, après simplification par $\Delta t$ :
\[ \boxed{|e| = B \, \ell \, v} \]

\textbf{Vérification d'homogénéité} : $[B][\ell][v] = \text{T} \cdot \text{m} \cdot \text{m/s} = \dfrac{\text{V} \cdot \text{s}}{\text{m}^2} \times \dfrac{\text{m}^2}{\text{s}} = \text{V}$. La formule est homogène.

\textbf{Remarque} : le sens du courant induit se retrouve par Lenz — le courant induit crée une force de Laplace sur le barreau qui \emph{s'oppose à son mouvement} : c'est le freinage électromagnétique.
\end{experience}

\begin{popfigure}
\begin{center}
\begin{tikzpicture}[scale=1.05, >=Stealth]
% Champ B (croix = rentrant)
\foreach \x in {0.4,1.2,...,6.4} {\foreach \y in {0.35,1.05,1.75} {\node[popblue, font=\scriptsize] at (\x,\y) {$\times$};}}
\node[popblue, font=\small] at (6.6,2.3) {$\vec{B}$};
% Rails
\draw[very thick, popdark] (0,0) -- (6.5,0);
\draw[very thick, popdark] (0,2.1) -- (6.5,2.1);
\draw[very thick, popdark] (0,0) -- (0,2.1);
% Résistance
\node[left, popdark] at (-0.05,1.05) {$R$};
\fill[popgold] (0,1.05) circle (0.06);
% Surface balayée hachurée
\fill[pattern=north east lines, pattern color=poporange!70] (2.2,0) rectangle (4.4,2.1);
\node[poporange, font=\scriptsize] at (3.3,1.7) {surface balayée $\Delta S$};
% Barreau position initiale (pointillés) et finale
\draw[dashed, popmuted, thick] (2.2,0) -- (2.2,2.1);
\draw[line width=2.5pt, popgreen] (4.4,0) -- (4.4,2.1);
\node[popgreen, font=\small, above] at (4.4,2.1) {M};
\node[popgreen, font=\small, below] at (4.4,0) {N};
% Longueur ell
\draw[<->, popgreen] (4.55,0) -- (4.55,2.1) node[midway, right, font=\small] {$\ell$};
% Vitesse
\draw[->, very thick, poppink] (4.4,1.05) -- (5.6,1.05) node[right, font=\small] {$\vec{v}$};
% Déplacement
\draw[<->, popmuted] (2.2,-0.35) -- (4.4,-0.35) node[midway, below, font=\scriptsize] {$\Delta x = v\,\Delta t$};
\end{tikzpicture}
\end{center}
\legende{Barreau MN glissant sur des rails en U dans un champ $\vec{B}$ rentrant ($\times$). En $\Delta t$, le barreau balaye la surface hachurée $\Delta S = \ell \, v \, \Delta t$, d'où $|e| = B\ell v$.}
\end{popfigure}

\begin{exemplebox}[Rails et barreau : calcul complet]
Un barreau de longueur $\ell = \SI{20}{cm} = \SI{0,20}{m}$ glisse à $v = \SI{5,0}{m/s}$ sur des rails dans un champ $B = \SI{0,40}{T}$. Le circuit a une résistance totale $R = \SI{2,0}{\Omega}$.
\begin{itemize}
\item F.é.m. induite : $|e| = B \ell v = 0{,}40 \times 0{,}20 \times 5{,}0 = \SI{0,40}{V}$.
\item Courant induit : $i = \dfrac{|e|}{R} = \dfrac{0{,}40}{2{,}0} = \SI{0,20}{A}$.
\item Sens (Lenz) : le flux rentrant augmente (la surface grandit), donc le courant induit crée un champ \emph{sortant} : il circule dans le sens trigonométrique vu de face.
\end{itemize}
\end{exemplebox}

\courssub{Courant induit, puissance et bilan énergétique}

\begin{propriete}[Courant induit dans un circuit fermé]
Si le circuit induit est fermé et présente une résistance totale $R$, la f.é.m. induite $e$ y fait circuler un courant d'intensité :
\[ i = \frac{e}{R} \]
(loi d'Ohm appliquée au circuit induit, qui se comporte comme un générateur de f.é.m. $e$). Le \emph{sens} de $i$ est donné par la loi de Lenz. Si le circuit est \emph{ouvert}, aucun courant ne circule, mais la f.é.m. $e$ existe bel et bien entre les bornes ouvertes.
\end{propriete}

D'où vient l'énergie électrique ainsi produite ? Reprenons le barreau sur rails. Pour maintenir le barreau à vitesse constante malgré la force de Laplace de freinage $F = B i \ell$, un opérateur doit exercer une force opposée et fournir la puissance mécanique :
\[ P_{\text{méca}} = F \, v = B \, i \, \ell \, v = (B \ell v) \, i = e \, i = P_{\text{élec}} \]
La puissance électrique $P = e\,i$ apparue dans le circuit est \emph{exactement} égale au travail mécanique fourni par unité de temps : c'est une conversion d'énergie mécanique en énergie électrique, conforme au principe de conservation de l'énergie. La loi de Lenz n'est d'ailleurs qu'une conséquence de cette conservation : si le courant induit \emph{aidait} la variation de flux au lieu de s'y opposer, on créerait de l'énergie à partir de rien !

\begin{methode}[Calculer une f.é.m. induite : la démarche en 4 étapes]
\begin{enumerate}
\item \textbf{Exprimer le flux} $\Phi(t)$ à travers le circuit orienté : $\Phi = N B S \cos\alpha$, en identifiant la grandeur qui varie ($B$, $S$ ou $\alpha$) et en convertissant toutes les données en unités SI.
\item \textbf{Dériver} $\dfrac{d\Phi}{dt}$ (ou calculer $\dfrac{\Delta \Phi}{\Delta t}$ si la variation est régulière).
\item \textbf{Appliquer la loi de Faraday} : $e = -\dfrac{d\Phi}{dt}$ ; vérifier l'homogénéité (Wb/s $=$ V) et ne pas oublier $N$.
\item \textbf{Si le circuit est fermé} : calculer $i = \dfrac{|e|}{R}$ et déterminer le sens du courant par la loi de Lenz.
\end{enumerate}
\end{methode}

\begin{exercice}[Pour s'entraîner]
Une bobine circulaire de $N = 200$ spires, d'aire $S = \SI{50}{cm^2}$ chacune, est placée dans un champ uniforme perpendiculaire à son plan, dont la valeur passe de $\SI{0,10}{T}$ à $\SI{0,50}{T}$ en $\SI{0,10}{s}$.
\begin{enumerate}
\item Calculer la variation du flux total à travers la bobine.
\item En déduire la valeur absolue de la f.é.m. induite moyenne.
\item La bobine est fermée sur une résistance $R = \SI{10}{\Omega}$. Calculer l'intensité du courant induit et préciser son sens.
\end{enumerate}
\end{exercice}

\begin{corrige}[Exercice]
1. Conversion : $S = \SI{50}{cm^2} = \num{5,0e-3}$ m$^2$. Variation de flux par spire : $\Delta \Phi_1 = S \, \Delta B = \num{5,0e-3} \times (0{,}50 - 0{,}10) = \num{2,0e-3}$ Wb. Flux total : $\Delta \Phi = N \Delta \Phi_1 = 200 \times \num{2,0e-3} = \SI{0,40}{Wb}$.

2. $|e| = \dfrac{\Delta \Phi}{\Delta t} = \dfrac{0{,}40}{0{,}10} = \SI{4,0}{V}$.

3. $i = \dfrac{|e|}{R} = \dfrac{4{,}0}{10} = \SI{0,40}{A}$. Le flux (perpendiculaire au plan) augmente : par la loi de Lenz, le courant induit crée un champ magnétique \emph{opposé} au champ extérieur ; son sens s'en déduit par la règle du tire-bouchon (ou de la main droite).
\end{corrige}

\begin{savaistu}[De Faraday à Nachtigal]
Michael Faraday découvre l'induction en 1831 avec une bobine et un aimant ; moins de cinquante ans plus tard, les premiers alternateurs industriels éclairent les villes. Aujourd'hui, le barrage de Nachtigal (420 MW) convertit l'énergie mécanique de l'eau du fleuve Sanaga en énergie électrique grâce à la loi $e = -N\,\dfrac{d\Phi}{dt}$ : des aimants tournent devant des bobines géantes, et la variation périodique du flux y engendre la tension alternative qui alimente le réseau. La même loi, à petite échelle, fait fonctionner la dynamo de certaines lampes de vélo et les chargeurs à manivelle utilisés dans les zones rurales non électrifiées.
\end{savaistu}

\begin{aretenir}[L'essentiel de la section]
\begin{itemize}
\item Loi de Faraday : $e = -\dfrac{d\Phi}{dt}$ ; le signe « $-$ » traduit la loi de Lenz.
\item Homogénéité : $1$ Wb/s $= 1$ V.
\item Bobine de $N$ spires : $e = -N\,\dfrac{d(BS\cos\alpha)}{dt}$ — ne jamais oublier $N$.
\item Barreau sur rails : $|e| = B\,\ell\,v$ ; bobine en rotation : $e = NBS\omega\sin(\omega t)$ (alternateur).
\item Circuit fermé : $i = \dfrac{e}{R}$, sens donné par Lenz ; puissance $P = e\,i$ issue du travail mécanique (conservation de l'énergie).
\end{itemize}
\end{aretenir}

\coursec{Auto-induction et inductance}

Dans la section précédente, nous avons vu comment un aimant en mouvement ou un circuit inducteur traversé par un courant variable crée un flux magnétique variable à travers un circuit voisin, induisant une f.é.m. selon la loi de Faraday. Mais que se passe-t-il si le circuit \emph{lui-même} est la source de cette variation de flux ? C'est le phénomène d'\cle{auto-induction} : un circuit traversé par un courant variable devient à la fois inducteur (il crée le champ) et induit (il subit la variation de son propre flux). C'est le cœur du fonctionnement des bobines, transformateurs, et de l'allumage de vos motos-taxi « benskin » !

\courssub{Mise en évidence expérimentale de l'auto-induction}

\begin{experience}[Retard à l'allumage et étincelle de rupture]
\textbf{Matériel :} une source de tension continue $E$ (pile ou alimentation), un interrupteur $K$, deux branches parallèles : branche 1 (lampe $L_1$ en série avec une bobine $\mathcal{B}$ de forte inductance $L$ et résistance interne $r$), branche 2 (lampe $L_2$ identique en série avec une résistance $R = r$).\\
\textbf{Protocole :}
\begin{enumerate}
\item Fermer l'interrupteur $K$ et observer l'allumage des deux lampes.
\item Ouvrir brutalement l'interrupteur $K$ et observer ce qui se produit aux bornes de la bobine.
\end{enumerate}
\textbf{Observations :}
\begin{itemize}
\item À la fermeture : la lampe $L_2$ (branche résistive) s'allume \emph{instantanément} à sa luminosité nominale. La lampe $L_1$ (branche bobine) met un \emph{temps perceptible} pour atteindre la même luminosité : il y a un \textbf{retard à l'établissement} du courant.
\item À l'ouverture : une \textbf{étincelle} jaillit aux contacts de l'interrupteur $K$ (ou aux bornes de la bobine si on l'ouvre directement), alors qu'aucune étincelle n'apparaît sur la branche purement résistive.
\end{itemize}
\textbf{Interprétation :}
\begin{itemize}
\item \emph{Fermeture :} le courant $i$ dans la bobine croît de $0$ vers sa valeur finale. Ce courant variable crée un flux magnétique variable \emph{à travers la bobine elle-même}. D'après la loi de Lenz, la f.é.m. d'auto-induction induite s'oppose à cette \emph{augmentation} : elle est orientée \emph{contraire} au sens du courant imposé par le générateur. Elle freine la montée du courant, d'où le retard.
\item \emph{Ouverture :} le courant tente de tomber brutalement à zéro. La variation $\dfrac{di}{dt}$ est très grande et négative. La f.é.m. d'auto-induction s'oppose à cette \emph{diminution} : elle tend à \emph{maintenir} le courant dans le même sens. Elle crée une tension très élevée aux bornes de la bobine (souvent plusieurs centaines de volts), ionisant l'air aux contacts de l'interrupteur : c'est l'étincelle de rupture.
\end{itemize}
\end{experience}

\begin{popfigure}
\begin{tikzpicture}[scale=0.9, transform shape, european]
% Source
\draw (0,0) to[battery1, l={$E$}, invert] (0,3);
% Interrupteur principal
\draw (0,3) -- (2,3) to[nos, l={$K$}] (4,3);
% Nœud de dérivation
\draw (4,3) -- (6,3);
% Branche haute : Lampe + Résistance
\draw (6,3) to[lamp, l={$L_2$}] (6,1.5) to[R, l={$R=r$}] (6,0);
% Branche basse : Lampe + Bobine
\draw (6,3) -- (8,3) to[lamp, l={$L_1$}] (8,1.5) to[L, l={$\mathcal{B}$ ($L$, $r$)}] (8,0);
% Retour
\draw (6,0) -- (4,0) -- (0,0);
\draw (8,0) -- (6,0);
% Annotations
\node[align=center] at (3,3.7) {Branche résistive\\ (réponse instantanée)};
\node[align=center] at (9.5,1.75) {Branche inductive\\ (retard à l'allumage)};
\end{tikzpicture}
\legende{Circuit de mise en évidence de l'auto-induction : comparaison d'une branche résistive et d'une branche inductive en parallèle.}
\end{popfigure}

\courssub{Flux propre et inductance}

Lorsqu'un courant $i$ traverse une bobine, il crée un champ magnétique $\vec{B}$ à l'intérieur de ses spires. Le flux de $\vec{B}$ à travers la surface portée par \emph{chaque} spire s'additionne. On appelle \cle{flux propre} (ou flux d'auto-induction) $\Phi_p$ la somme des flux à travers les $N$ spires : $\Phi_p = N \Phi_{\text{1 spire}}$.

\begin{definition}[\cle{Auto-induction}]
L'\textbf{auto-induction} est le phénomène d'induction d'une force électromotrice dans un circuit due à la variation du courant dans \emph{ce même circuit}. Le circuit est simultanément \emph{inducteur} (source du champ) et \emph{induit} (siège de la f.é.m. induite).
\end{definition}

Pour une bobine donnée (géométrie, nombre de spires, noyau magnétique fixés), le champ magnétique est proportionnel au courant $i$ (loi de Biot-Savart ou loi d'Ampère). Le flux propre est donc proportionnel à l'intensité du courant :

\begin{definition}[\cle{Inductance}]
L'\textbf{inductance} $L$ d'une bobine est le rapport entre le flux propre $\Phi_p$ qui la traverse et l'intensité $i$ du courant qui la génère :
\[
\Phi_p = L \cdot i
\]
$L$ ne dépend que de la géométrie de la bobine (longueur, section, nombre de spires $N^2$) et de la perméabilité magnétique $\mu$ du milieu (air ou noyau ferromagnétique).
\end{definition}

\begin{propriete}[\cle{Unité de l'inductance}]
L'unité SI de l'inductance est le \cle{henry} (symbole \textbf{H}), en l'honneur de Joseph Henry.
\[
\SI{1}{H} = \SI{1}{Wb/A} = \SI{1}{V.s/A} = \SI{1}{kg.m^2/(s^2.A^2)}
\]
Les valeurs courantes en électronique vont du microhenry ($\mu\text{H} = \SI{e-6}{H}$) au millihenry ($\text{mH} = \SI{e-3}{H}$). Les bobines de puissance (transformateurs, moteurs) peuvent atteindre plusieurs henrys.
\end{propriete}

\begin{aretenir}[Analyse dimensionnelle de l'inductance]
Vérifions : $[\Phi_p] = [B]\cdot[S] = \text{T}\cdot\text{m}^2 = \text{Wb}$. $[i] = \text{A}$. Donc $[L] = \text{Wb}/\text{A} = \text{H}$. Cohérent.
On peut aussi écrire $L = \dfrac{N\Phi}{i}$. Pour un solénoïde long de $N$ spires, longueur $l$, section $S$, perméabilité $\mu$ : $B = \mu n i = \mu \frac{N}{l} i$, $\Phi = B S = \mu \frac{N}{l} i S$, donc $L = \frac{N\Phi}{i} = \mu \frac{N^2}{l} S$. $L \propto N^2$ : doubler le nombre de spires quadruple l'inductance.
\end{aretenir}

\courssub{F.é.m. d'auto-induction et tension aux bornes d'une bobine réelle}

La f.é.m. d'auto-induction $e$ est la f.é.m. induite par la variation du flux propre. En appliquant la loi de Faraday ($e = -\dfrac{d\Phi_p}{dt}$) et la définition $\Phi_p = L i$ (en supposant $L$ constant, ce qui est vrai pour une bobine à noyau d'air ou en dehors de la saturation magnétique), on obtient :

\begin{propriete}[\cle{F.é.m. d'auto-induction}]
La f.é.m. d'auto-induction dans une bobine d'inductance $L$ traversée par un courant $i(t)$ est :
\[
e = -\frac{d\Phi_p}{dt} = -L\,\frac{di}{dt}
\]
\begin{itemize}
\item Si $\dfrac{di}{dt} > 0$ (courant croissant) : $e < 0$. La f.é.m. s'oppose à l'augmentation (loi de Lenz).
\item Si $\dfrac{di}{dt} < 0$ (courant décroissant) : $e > 0$. La f.é.m. s'oppose à la diminution (elle tend à maintenir le courant).
\item En régime permanent ($di/dt = 0$) : $e = 0$. La bobine ne manifeste plus d'effet d'auto-induction.
\end{itemize}
\emph{Note :} Cette relation est admise au Probatoire comme conséquence directe de la loi de Faraday. La démonstration rigoureuse fait appel à l'équation de Maxwell-Faraday sous forme intégrale.
\end{propriete}

Une bobine réelle possède toujours une résistance interne $r$ (fil de cuivre). En convention récepteur (le courant $i$ entre par la borne de potentiel supérieur $A$ et sort par $B$), la tension $u = u_A - u_B$ aux bornes de la bobine est la somme de la chute de tension ohmique et de la f.é.m. d'auto-induction (qui agit comme une source de tension interne) :

\begin{propriete}[\cle{Tension aux bornes d'une bobine réelle}]
En convention récepteur ($A$ borne d'entrée du courant, $B$ borne de sortie) :
\[
\boxed{u = r\,i + L\,\frac{di}{dt}}
\]
\begin{itemize}
\item Terme $r\,i$ : chute de tension résistive (effet Joule, irréversible).
\item Terme $L\,\dfrac{di}{dt}$ : tension d'auto-induction (réversible, liée au stockage d'énergie magnétique).
\end{itemize}
\emph{Cas particulier : régime permanent continu.} $\dfrac{di}{dt} = 0 \Rightarrow u = r\,i$. La bobine se comporte comme un simple conducteur ohmique de résistance $r$.
\end{propriete}

\begin{attention}[Pièges classiques — À ne pas commettre au Probatoire !]
\begin{enumerate}
\item \textbf{Oublier la résistance interne $r$.} Écrire $u = L\,\dfrac{di}{dt}$ pour une bobine réelle est \textbf{faux}. La relation $u = L\,di/dt$ ne vaut que pour une bobine \emph{idéale} ($r=0$). Tout énoncé parlant de « bobine de résistance $r$ » ou « bobine réelle » impose $u = r i + L\,di/dt$.
\item \textbf{Dire « la bobine s'oppose au courant ».} \textbf{Faux.} La bobine s'oppose à la \emph{variation} du courant ($di/dt$). En régime permanent, elle laisse passer le courant sans s'y opposer (hormis sa résistance $r$).
\item \textbf{Confondre le sens de $e$ et le sens de $u$.} $e = -L\,di/dt$ est une f.é.m. (générateur). $u = r i + L\,di/dt$ est une tension (récepteur). Le terme $L\,di/dt$ a le \emph{même signe} que $di/dt$ dans l'expression de $u$, mais le signe \emph{opposé} dans l'expression de $e$. Relisez la convention récepteur : $u = \sum (\text{chutes résistives}) - \sum (\text{f.é.m. internes orientées dans le sens du courant})$. Ici $e$ s'oppose à la variation, donc si $di/dt>0$, $e<0$ (orienté contre $i$), donc $-e = +L\,di/dt$ s'ajoute à $r i$ dans $u$.
\item \textbf{Unités :} $L$ en henry (H), $i$ en ampère (A), $t$ en seconde (s) $\Rightarrow$ $L\,di/dt$ en volt (V). Vérifiez toujours : $\text{H} = \text{V}\cdot\text{s}/\text{A}$.
\end{enumerate}
\end{attention}

\courssub{Interprétation énergétique : stockage de l'énergie magnétique}

La puissance électrique reçue par la bobine (convention récepteur) est $p = u\,i = r\,i^2 + L\,i\,\dfrac{di}{dt}$.
\begin{itemize}
\item Le terme $r\,i^2$ est la puissance dissipée par effet Joule (perte de chaleur, irréversible).
\item Le terme $L\,i\,\dfrac{di}{dt} = \dfrac{d}{dt}\left(\frac{1}{2}L i^2\right)$ est la puissance \emph{stockée} sous forme d'énergie magnétique dans le champ créé par la bobine.
\end{itemize}
L'\cle{énergie magnétique} stockée dans une bobine d'inductance $L$ traversée par un courant $i$ est :
\[
W_m = \frac{1}{2} L i^2 \qquad (\text{en joules, J})
\]

\begin{savaistu}[L'étincelle de rupture et l'allumage « benskin »]
C'est exactement ce phénomène qui fait démarrer votre moto-taxi « benskin » ! La \textbf{bobine d'allumage} est un transformateur élévateur : une bobine primaire (quelques tours, $L \sim \text{mH}$) est traversée par le courant de la batterie ($\SI{12}{V}$). Un rupteur (contact mécanique ou électronique) ouvre brutalement le circuit primaire. Le courant chute violemment ($di/dt$ énorme, négatif), créant une f.é.m. d'auto-induction $e = -L\,di/dt$ de plusieurs \textbf{centaines de volts} aux bornes du primaire. Par couplage magnétique (induction mutuelle), cette surtension est multipliée par le rapport de tours (souvent $\times 100$) dans le secondaire, produisant \textbf{\SIrange{10000}{30000}{V}} à la bougie. L'étincelle enflamme le mélange air-essence. Sans auto-induction, pas d'étincelle, pas de moteur thermique ! De même, les \textbf{parafoudres} et \textbf{varistances} protègent le réseau ENEO/SONATREL en dérivant les surtensions d'origine inductive (foudre, manœuvres de coupure).
\end{savaistu}

\begin{popfigure}
\begin{tikzpicture}
\begin{axis}[
    width=\linewidth, height=7cm,
    axis lines=middle,
    xlabel={$t$ (s)}, ylabel={$i$ (A)},
    xmin=0, xmax=5, ymin=0, ymax=2.2,
    xtick={0,0.4,1,2,3,4,5}, ytick={0,1,2},
    domain=0:5, samples=200,
    legend pos=south east,
    grid=major, grid style={popline},
]
% Courbe avec bobine : i(t) = I_max (1 - exp(-t/tau))
% E = 10 V, r = 5 ohm => I_max = 2 A. L = 2 H => tau = L/r = 0.4 s.
\addplot[popcurve, thick, popblue] {2*(1-exp(-2.5*x))};
\addlegendentry{$i(t)$ avec bobine $L=\SI{2}{H}, r=\SI{5}{\ohm}$ ($\tau=\SI{0,4}{s}$)}
% Courbe sans bobine (établissement instantané)
\addplot[poporange, thick, dashed] coordinates {(0,0) (0.01,2) (5,2)};
\addlegendentry{$i(t)$ sans bobine (résistance seule)}
% Ligne asymptote
\addplot[popmuted, thin, dashed] coordinates {(0,2) (5,2)};
% Flèche tau à 0.4s (63% de Imax = 1.264 A)
\draw[<->, popdark] (axis cs:0,1.264) -- (axis cs:0.4,1.264) node[midway, above, font=\small] {$\tau$};
\draw[popdark, dashed] (axis cs:0.4,0) -- (axis cs:0.4,1.264);
\node[popdark, font=\small] at (axis cs:0.4, -0.3) {$\tau = L/r$};
\end{axis}
\end{tikzpicture}
\legende{Établissement du courant dans une bobine réelle $(L, r)$ sous tension continue $E$ : croissance exponentielle vers $I_{\text{max}} = E/r$ avec une constante de temps $\tau = L/r$. Sans bobine, le courant saute instantanément à $E/R$.}
\end{popfigure}

\courssub{Exemple chiffré et exercice résolu}

\begin{exemplebox}[Calcul de la f.é.m. d'auto-induction]
Une bobine réelle d'inductance $L = \SI{0,20}{H}$ et de résistance interne $r = \SI{5,0}{\ohm}$ est branchée sur une source de tension continue $E = \SI{10,0}{V}$. Au cours de l'établissement, le courant croît de $0$ à $\SI{2,0}{A}$ en $\SI{0,10}{s}$ (on suppose une variation linéaire pour simplifier le calcul).
\begin{enumerate}
\item Calculer la valeur moyenne de la f.é.m. d'auto-induction $\langle e \rangle$ pendant cet intervalle.
\item Calculer la tension moyenne $\langle u \rangle$ aux bornes de la bobine pendant ce même intervalle.
\item Quelle est l'énergie magnétique stockée dans la bobine lorsque le courant a atteint $\SI{2,0}{A}$ ?
\end{enumerate}
\end{exemplebox}

\begin{exoresolu}[Solution détaillée]
\textbf{Données :} $L = \SI{0,20}{H}$, $r = \SI{5,0}{\ohm}$, $\Delta i = \SI{2,0}{A} - 0 = \SI{2,0}{A}$, $\Delta t = \SI{0,10}{s}$.
\textbf{1. f.é.m. d'auto-induction moyenne :}
La variation de courant est supposée linéaire, donc $\dfrac{di}{dt} \approx \dfrac{\Delta i}{\Delta t} = \dfrac{2,0}{0,10} = \SI{20}{A/s}$.
\[
\langle e \rangle = -L\,\frac{\Delta i}{\Delta t} = -0,20 \times 20 = \SI{-4,0}{V}
\]
La grandeur de la f.é.m. est $|\langle e \rangle| = \SI{4,0}{V}$. Le signe négatif indique qu'elle s'oppose à l'augmentation du courant (loi de Lenz).
\textbf{2. Tension aux bornes de la bobine (convention récepteur) :}
Le courant moyen pendant l'intervalle est $\langle i \rangle = \dfrac{0 + 2,0}{2} = \SI{1,0}{A}$.
\[
\langle u \rangle = r\,\langle i \rangle + L\,\frac{\Delta i}{\Delta t} = 5,0 \times 1,0 + 0,20 \times 20 = 5,0 + 4,0 = \SI{9,0}{V}
\]
\textit{Vérification :} Au début ($t=0$, $i=0$), $u(0) = 0 + 4,0 = \SI{4,0}{V}$. À la fin ($t=\SI{0,10}{s}$, $i=\SI{2,0}{A}$), $u = 5,0\times 2,0 + 4,0 = \SI{14,0}{V}$. La moyenne $(4+14)/2 = \SI{9,0}{V}$ est cohérente. Noter que $u > E$ en fin d'intervalle car on a supposé une variation linéaire continue ; en réalité, $di/dt$ diminue exponentiellement et $u$ tend vers $E = r I_{\text{max}} = \SI{10}{V}$.
\textbf{3. Énergie magnétique stockée :}
\[
W_m = \frac{1}{2} L i^2 = \frac{1}{2} \times 0,20 \times (2,0)^2 = 0,10 \times 4,0 = \SI{0,40}{J}
\]
Cette énergie a été fournie par le générateur pendant la phase transitoire. À l'ouverture du circuit, elle sera restituée (en partie sous forme d'étincelle, en partie en chaleur dans l'arc et la résistance interne).
\tcblower
\textbf{Bilan :} $\langle e \rangle = \SI{-4,0}{V}$ ; $\langle u \rangle = \SI{9,0}{V}$ ; $W_m = \SI{0,40}{J}$.
\end{exoresolu}

\begin{exercice}[Application : bobine de relais]
Un relais électromécanique (très utilisé dans l'automatisme industriel camerounais, ex. commande de pompes d'irrigation) possède une bobine d'inductance $L = \SI{50}{mH}$ et de résistance $r = \SI{25}{\ohm}$. Il est alimenté sous $E = \SI{24}{V}$.
\begin{enumerate}
\item Calculer la constante de temps $\tau$ du circuit.
\item Déterminer l'intensité du courant $\SI{0,10}{ms}$ après la fermeture de l'interrupteur (on utilise $i(t) = \frac{E}{r}(1 - e^{-t/\tau})$).
\item Calculer la f.é.m. d'auto-induction $e$ à cet instant.
\item Quelle est l'énergie magnétique stockée à l'état permanent ?
\end{enumerate}
\end{exercice}

\begin{corrige}[Exercice : bobine de relais]
\textbf{1.} $\tau = \dfrac{L}{r} = \dfrac{50 \times 10^{-3}}{25} = 2,0 \times 10^{-3}\,\text{s} = \SI{2,0}{ms}$.
\textbf{2.} $I_{\text{max}} = E/r = 24/25 = \SI{0,96}{A}$.
$t = \SI{0,10}{ms} = 1,0 \times 10^{-4}\,\text{s}$.
$t/\tau = 0,10 / 2,0 = 0,05$.
$i(t) = 0,96 \times (1 - e^{-0,05}) \approx 0,96 \times (1 - 0,9512) \approx 0,96 \times 0,0488 \approx \SI{0,047}{A}$.
\textbf{3.} $\dfrac{di}{dt} = \dfrac{E}{L} e^{-t/\tau} = \dfrac{24}{50 \times 10^{-3}} e^{-0,05} = 480 \times 0,9512 \approx \SI{457}{A/s}$.
$e = -L\,\dfrac{di}{dt} = -50 \times 10^{-3} \times 457 \approx \SI{-22,8}{V}$.
(La f.é.m. d'auto-induction est forte au début car $di/dt$ est maximal).
\textbf{4.} $W_m = \frac{1}{2} L I_{\text{max}}^2 = 0,5 \times 50 \times 10^{-3} \times (0,96)^2 \approx \SI{0,023}{J} = \SI{23}{mJ}$.
\end{corrige}

\begin{methode}[Résoudre un problème d'auto-induction (circuit RL série)]
\begin{enumerate}
\item \textbf{Identifier} le circuit : source de tension $E$ (ou $u_g(t)$), résistance $R$ (incluant $r$ de la bobine), inductance $L$. Noter le sens du courant $i$ (convention récepteur sur la bobine).
\item \textbf{Écrire la loi des mailles} (somme des tensions = 0) ou le bilan de puissance. En convention récepteur sur $R$ et $L$ : $u_g = u_R + u_L = R i + (r i + L\,di/dt)$. Souvent $R$ est la résistance totale : $u_g = R_{\text{tot}} i + L\,di/dt$.
\item \textbf{Résoudre l'équation différentielle} du premier ordre : $L\,\dfrac{di}{dt} + R i = u_g(t)$.
\begin{itemize}
\item Régime permanent continu ($u_g = E$ cte) : $i_{\text{final}} = E/R$.
\item Établissement (fermeture à $t=0$, $i(0)=0$) : $i(t) = \dfrac{E}{R}(1 - e^{-t/\tau})$, $\tau = L/R$.
\item Vidange (ouverture, source court-circuitée ou $u_g=0$) : $i(t) = i_0 e^{-t/\tau}$.
\end{itemize}
\item \textbf{Calculer les grandeurs demandées} : $e = -L\,di/dt$, $u_L = r i + L\,di/dt$, $W_m = \frac{1}{2}L i^2$, énergie dissipée $\int r i^2 dt$.
\item \textbf{Vérifier} : homogénéité des unités, valeurs limites ($t=0$, $t\to\infty$), signe de $e$ conforme à Lenz.
\end{enumerate}
\end{methode}

\begin{aretenir}[\cle{Bilan de la section : Auto-induction et inductance}]
\begin{itemize}
\item \textbf{Auto-induction :} Un circuit à courant variable induit une f.é.m. \emph{dans lui-même} ($e = -d\Phi_p/dt$).
\item \textbf{Flux propre :} $\Phi_p = L i$. \textbf{Inductance :} $L = \Phi_p/i$ (en henry H). $L$ caractérise la géométrie et le noyau.
\item \textbf{F.é.m. d'auto-induction :} $e = -L\,di/dt$. Elle s'oppose à la \emph{variation} du courant (Lenz).
\item \textbf{Tension bobine réelle $(L,r)$ :} $u = r i + L\,di/dt$ (convention récepteur).
\item \textbf{Régime permanent continu :} $di/dt=0 \Rightarrow e=0$, $u = r i$ (comportement résistif pur).
\item \textbf{Énergie magnétique :} $W_m = \frac{1}{2} L i^2$. Stockée pendant la montée, restituée pendant la descente (étincelle de rupture).
\item \textbf{Constante de temps :} $\tau = L/R$ (secondes). Temps caractéristique de l'établissement/vidange.
\end{itemize}
\end{aretenir}

\coursec{Applications et bilan énergétique}

Dans la section précédente, nous avons vu qu'une bobine traversée par un courant variable crée en son sein une f.é.m. d'auto-induction $e = -L\,\dfrac{di}{dt}$ qui s'oppose aux variations du courant. Ce phénomène d'auto-induction n'est qu'un cas particulier d'un phénomène plus général : \cle{l'induction électromagnétique}. Nous savons désormais que toute variation du flux magnétique $\Phi$ à travers un circuit y fait naître une f.é.m. induite $e = -\dfrac{d\Phi}{dt}$ (loi de Faraday) dont les effets s'opposent à la cause qui leur donne naissance (loi de Lenz). Il est temps de répondre à la question qui donne tout son sens à cette leçon : \emph{à quoi cela sert-il ?} La réponse tient en une phrase : l'induction électromagnétique est le principe physique de \textbf{quasi-totalité de la production d'électricité mondiale}, des centrales de Songloulou et d'Édéa jusqu'au groupe électrogène du quartier.

\courssub{L'alternateur : convertir le mouvement en électricité}

\courssub{Principe et description}

Reprenons l'expérience fondamentale de la section 1 : lorsqu'on approche ou éloigne un aimant d'une bobine, un courant induit apparaît. Mais au lieu de déplacer l'aimant à la main, on peut faire \textbf{tourner} une bobine de $N$ spires, d'aire $S$, à vitesse angulaire constante $\omega$ dans un champ magnétique uniforme $\vec{B}$. L'angle entre $\vec{B}$ et la normale $\vec{n}$ à la bobine varie alors selon $\alpha = \omega t$, et le flux devient :
\[ \Phi = NBS\cos(\omega t) \]
C'est un flux \textbf{variable} : d'après la loi de Faraday, il naît une f.é.m. induite :
\[ e = -\dfrac{d\Phi}{dt} = NBS\omega\sin(\omega t) = E_{\max}\sin(\omega t) \quad \text{avec} \quad E_{\max} = NBS\omega \]
La f.é.m. est \cle{alternative sinusoïdale} : elle change de signe à chaque demi-tour. C'est exactement la tension que délivrent les prises de courant (au Cameroun : valeur efficace $E \approx \SI{220}{V}$, fréquence $f = \SI{50}{Hz}$).

\begin{definition}[Alternateur]
Un \cle{alternateur} est un dispositif qui convertit de l'\textbf{énergie mécanique de rotation} en \textbf{énergie électrique} grâce à l'induction électromagnétique : une bobine (ou un aimant) mise en rotation dans un champ magnétique est le siège d'une f.é.m. induite alternative, transmise au circuit extérieur par des bagues et des balais.
\end{definition}

\begin{popfigure}
\begin{center}
\begin{tikzpicture}[scale=0.95, >=Stealth]
% Aimant (pôles)
\fill[popblue!70] (-3.2,-1.4) rectangle (-2.2,1.4);
\fill[poporange!80] (2.2,-1.4) rectangle (3.2,1.4);
\node[white, font=\bfseries] at (-2.7,0) {N};
\node[white, font=\bfseries] at (2.7,0) {S};
% Champ magnétique
\foreach \y in {-0.9,-0.3,0.3,0.9}{
  \draw[->, popblue, thick] (-2.1,\y) -- (-1.3,\y);
  \draw[popblue, thick] (1.3,\y) -- (2.1,\y);
}
\node[popblue] at (-1.7,1.25) {$\vec{B}$};
% Bobine en rotation (rectangle incliné)
\begin{scope}[rotate=30]
  \draw[very thick, popdark] (-0.9,-1.1) rectangle (0.9,1.1);
  \draw[->, popgreen, very thick] (0,0) -- (0.55,0.83) node[above right=-2pt] {$\vec{n}$};
\end{scope}
\draw[->, poporange, thick] (1.5,0.9) arc[start angle=35, end angle=95, radius=1.1];
\node[poporange] at (0.9,1.85) {$\omega$};
% Axe de rotation
\draw[dashed, popmuted] (0,-1.9) -- (0,1.9);
% Bagues
\fill[popgold] (-0.18,-1.75) rectangle (0.18,-1.55);
\fill[popgold] (-0.18,-2.15) rectangle (0.18,-1.95);
\node[right, font=\small] at (0.25,-1.85) {bagues + balais};
% Circuit extérieur
\draw[thick, popdark] (0,-1.65) -- (1.2,-1.65) -- (1.2,-3.0);
\draw[thick, popdark] (0,-2.05) -- (-1.2,-2.05) -- (-1.2,-3.0);
\draw[thick, popdark] (-1.2,-3.0) -- (-0.5,-3.0);
\draw[thick, popdark] (0.5,-3.0) -- (1.2,-3.0);
\draw[thick, popdark] (-0.5,-3.0) -- (-0.5,-3.25) -- (0.5,-3.25) -- (0.5,-3.0);
\node[below, font=\small] at (0,-3.3) {circuit extérieur $R$};
\node[font=\small] at (0,-2.55) {$e(t)$};
\end{tikzpicture}
\end{center}
\legende{Schéma de principe d'un alternateur : la bobine tourne à la vitesse angulaire $\omega$ dans le champ $\vec{B}$ ; le flux $\Phi = NBS\cos(\omega t)$ varie, d'où une f.é.m. induite alternative recueillie par les bagues.}
\end{popfigure}

\begin{propriete}[Conversion d'énergie dans l'alternateur]
L'induction électromagnétique réalise la conversion d'\textbf{énergie mécanique} en \textbf{énergie électrique}. La \textbf{loi de Lenz garantit la conservation de l'énergie} : le courant induit crée, par les forces de Laplace, un couple résistant qui s'oppose à la rotation. Pour maintenir la rotation, l'opérateur (ou la turbine) doit fournir un travail mécanique \emph{au moins égal} à l'énergie électrique récupérée. On ne peut pas obtenir d'électricité « gratuitement » : plus on débite de courant, plus il est difficile de tourner.
\end{propriete}

\begin{savaistu}[Les centrales hydroélectriques du Cameroun]
Les centrales de \textbf{Songloulou} (environ $\SI{384}{MW}$), d'\textbf{Édéa} ($\SI{263}{MW}$) sur la Sanaga et de \textbf{Lagdo} ($\SI{72}{MW}$) sur la Bénoué fonctionnent toutes sur ce principe : l'eau en chute fait tourner des turbines, qui entraînent de gigantesques alternateurs. L'énergie potentielle de l'eau devient énergie mécanique de rotation, puis énergie électrique distribuée par ENEO. Le récent barrage de Nachtigal ($\SI{420}{MW}$) utilise la même physique que celle de cette leçon !
\end{savaistu}

\courssub{Le transformateur : l'induction sans contact électrique}

Un transformateur comporte deux bobines enroulées sur un même \textbf{noyau ferromagnétique} fermé, sans aucune connexion électrique entre elles :
\begin{itemize}
\item le \cle{primaire} ($N_1$ spires), alimenté par une tension \textbf{alternative} $u_1$ ;
\item le \cle{secondaire} ($N_2$ spires), connecté au récepteur.
\end{itemize}
Le courant alternatif du primaire crée un champ magnétique variable ; le noyau \textbf{canalise} ce champ (les lignes de champ suivent le fer) de sorte que pratiquement tout le flux traverse le secondaire. Ce flux variable induit au secondaire une f.é.m. : une tension $u_2$ apparaît \emph{sans aucun contact électrique} entre les deux circuits. En régime idéal :
\[ \dfrac{U_2}{U_1} = \dfrac{N_2}{N_1} \]
Si $N_2 > N_1$, le transformateur est \textbf{survolteur} ; si $N_2 < N_1$, il est \textbf{abaisseur} (comme celui du chargeur de téléphone : $\SI{220}{V}$ vers $\SI{5}{V}$).

\begin{popfigure}
\begin{center}
\begin{tikzpicture}[scale=0.9, >=Stealth]
% Noyau ferromagnétique
\draw[very thick, popdark, fill=popcream] (-2.6,-1.8) rectangle (2.6,1.8);
\draw[very thick, popdark, fill=white] (-1.4,-0.9) rectangle (1.4,0.9);
% Flux canalisé
\draw[->, popblue, thick] (0,1.35) -- (0.9,1.35);
\draw[->, popblue, thick] (0,-1.35) -- (-0.9,-1.35);
\draw[->, popblue, thick] (-2.0,0.6) -- (-2.0,-0.4);
\draw[->, popblue, thick] (2.0,-0.6) -- (2.0,0.4);
\node[popblue, font=\small] at (0.2,1.62) {flux $\Phi$ canalisé};
% Bobine primaire (gauche)
\foreach \y in {-0.75,-0.45,-0.15,0.15,0.45,0.75}{
  \draw[thick, poporange] (-2.6,\y) arc[start angle=-90, end angle=90, x radius=0.09, y radius=0.15];
}
\node[poporange, font=\small, align=center] at (-3.6,0.4) {primaire\\ $N_1$ spires};
\draw[thick, poporange] (-2.6,-0.9) -- (-2.6,-2.4) node[below, font=\small] {$u_1$ alternatif};
\draw[thick, poporange] (-2.6,0.9) -- (-2.6,2.2);
% Bobine secondaire (droite)
\foreach \y in {-0.6,-0.3,0,0.3,0.6}{
  \draw[thick, popgreen] (2.6,\y) arc[start angle=90, end angle=270, x radius=0.09, y radius=0.15];
}
\node[popgreen, font=\small, align=center] at (3.7,0.4) {secondaire\\ $N_2$ spires};
\draw[thick, popgreen] (2.6,-0.75) -- (2.6,-2.4) node[below, font=\small] {$u_2$ induite};
\draw[thick, popgreen] (2.6,0.75) -- (2.6,2.2);
\end{tikzpicture}
\end{center}
\legende{Transformateur : le courant alternatif du primaire crée un flux variable canalisé par le noyau ; ce flux induit une f.é.m. au secondaire, sans contact électrique.}
\end{popfigure}

\begin{attention}[Le transformateur ne fonctionne PAS en courant continu]
Erreur très fréquente au Probatoire : croire qu'on peut alimenter le primaire d'un transformateur avec une pile. Si le courant primaire est \textbf{continu} (constant), le flux dans le noyau est \textbf{constant}, donc $\dfrac{d\Phi}{dt} = 0$ : \textbf{aucune f.é.m. induite} au secondaire. Le transformateur exige un courant \textbf{variable} (en pratique, alternatif). C'est d'ailleurs pourquoi le transport d'électricité sur les lignes haute tension de SONATREL se fait en courant alternatif : on peut élever la tension à $\SI{225}{kV}$ pour limiter les pertes par effet Joule, puis la rabaisser près des habitations.
\end{attention}

\courssub{D'autres applications de l'induction}

\begin{itemize}
\item \textbf{Le microphone dynamique} : la voix fait vibrer une membrane solidaire d'une bobine placée dans le champ d'un aimant. Le mouvement de la bobine fait varier le flux : une tension induite, image fidèle du son, apparaît aux bornes. C'est un alternateur miniature.
\item \textbf{La cuisinière à induction} : sous la plaque vitrocéramique, une bobine parcourue par un courant alternatif de haute fréquence crée un champ variable. Ce champ induit des \textbf{courants de Foucault} dans le fond métallique de la casserole, qui chauffe par effet Joule. La plaque elle-même reste presque froide : c'est le récipient qui est la source de chaleur.
\item \textbf{Les cartes à puce sans contact} (badges, cartes bancaires) : le lecteur émet un champ magnétique variable ; ce champ induit dans la bobine de la carte une f.é.m. qui \emph{alimente} la puce et lui permet de répondre. La carte n'a pas de pile : elle puise son énergie par induction.
\item \textbf{Le freinage électromagnétique} des poids lourds et des TGV : des courants de Foucault induits dans un disque métallique créent, par la loi de Lenz, des forces qui s'opposent au mouvement — freinage sans usure mécanique.
\end{itemize}

\courssub{Bilan énergétique global}

Considérons un alternateur qui débite un courant $i$ dans un circuit extérieur. Ce courant circule dans la bobine en rotation : chaque spire est alors soumise à des \textbf{forces de Laplace} $\vec{F} = i\,\vec{\ell}\wedge\vec{B}$. D'après la loi de Lenz, ces forces créent un couple qui \emph{s'oppose à la rotation} (c'est un couple résistant). Pour entretenir le mouvement, il faut donc fournir un \textbf{travail mécanique} moteur. Le bilan de puissance s'écrit :
\[ P_{\text{méca}} = P_{\text{élec utile}} + P_{\text{pertes}} \]
où les pertes sont essentiellement l'\textbf{effet Joule} dans les bobinages ($r\,i^2$) et les pertes ferromagnétiques (courants de Foucault dans le noyau, limités par le feuilletage).

\begin{methode}[Conduire un bilan énergétique]
\begin{enumerate}
\item \textbf{Identifier la source de travail mécanique} : turbine à eau ou à vapeur, moteur thermique, bras qui pédale...
\item \textbf{Identifier la f.é.m. induite} et le circuit induit ; calculer la puissance électrique utile $P_{\text{élec}} = E\,I$ (en alternatif sinusoïdal : $P = U\,I\cos\varphi$, ou $P = UI$ si le récepteur est résistif).
\item \textbf{Évaluer les pertes} : effet Joule $rI^2$ dans les bobinages, pertes mécaniques (frottements).
\item \textbf{Écrire la conservation} : $P_{\text{méca}} = P_{\text{élec utile}} + P_{\text{pertes}}$, puis le rendement $\eta = \dfrac{P_{\text{élec utile}}}{P_{\text{méca}}} < 1$.
\end{enumerate}
\end{methode}

\begin{exemplebox}[Puissance d'un alternateur domestique]
Un groupe électrogène débite un courant efficace $I = \SI{10}{A}$ sous une tension efficace $U = \SI{220}{V}$ dans une installation résistive. La puissance électrique fournie est :
\[ P = U\,I = 220 \times 10 = \SI{2200}{W} = \SI{2,2}{kW} \]
Chaque seconde, l'alternateur doit donc recevoir \emph{au minimum} un travail mécanique $W = P \times \Delta t = \SI{2,2}{kJ}$. Si le rendement global est $\eta = 80\,\%$, le moteur thermique doit en réalité fournir $P_{\text{méca}} = \dfrac{2{,}2}{0{,}80} \approx \SI{2,75}{kW}$ : c'est l'essence du groupe qui paie la facture énergétique, conformément à la loi de Lenz.
\end{exemplebox}

\begin{exoresolu}[Bilan énergétique d'une dynamo de vélo]
Une dynamo de vélo alimente une lampe qui consomme une puissance électrique $P_{\text{élec}} = \SI{3,0}{W}$. Le rendement de la dynamo est $\eta = 60\,\%$. Le cycliste roule à $v = \SI{18}{km/h}$.
\begin{enumerate}
\item Quelle puissance mécanique supplémentaire le cycliste doit-il fournir ?
\item Quelle force résistante supplémentaire cela représente-t-il ?
\end{enumerate}
\tcblower
\textbf{1.} Le rendement s'écrit $\eta = \dfrac{P_{\text{élec}}}{P_{\text{méca}}}$, d'où :
\[ P_{\text{méca}} = \dfrac{P_{\text{élec}}}{\eta} = \dfrac{3{,}0}{0{,}60} = \SI{5,0}{W} \]
\textbf{2.} La puissance mécanique d'une force de traction est $P = F\,v$. Convertissons : $v = \SI{18}{km/h} = \SI{5,0}{m/s}$. Alors :
\[ F = \dfrac{P_{\text{méca}}}{v} = \dfrac{5{,}0}{5{,}0} = \SI{1,0}{N} \]
C'est faible, mais perceptible sur un long trajet : c'est le « prix » du travail contre les forces de Laplace, exactement comme le prévoit la loi de Lenz.
\end{exoresolu}

\courssub{Synthèse de la leçon : formules et unités}

\begin{aretenir}[Les formules essentielles de l'induction]
\begin{itemize}
\item \textbf{Flux magnétique} à travers une bobine de $N$ spires d'aire $S$, dont la normale fait l'angle $\alpha$ avec $\vec{B}$ :
\[ \Phi = NBS\cos\alpha \]
\item \textbf{Loi de Faraday} (f.é.m. induite) :
\[ e = -\dfrac{d\Phi}{dt} \quad \text{(le signe } - \text{ traduit la loi de Lenz)} \]
\item \textbf{Flux propre} d'auto-induction : $\Phi_p = L\,i$
\item \textbf{F.é.m. d'auto-induction} :
\[ e = -L\,\dfrac{di}{dt} \]
\item \textbf{Tension aux bornes d'une bobine réelle} (résistance $r$, inductance $L$), en convention récepteur :
\[ u = r\,i + L\,\dfrac{di}{dt} \]
\end{itemize}
\end{aretenir}

\begin{center}
\rowcolors{1}{popblueL}{white}
\begin{tabularx}{\linewidth}{|l|c|c|X|}
\hline
\rowcolor{popblueL} \textbf{Grandeur} & \textbf{Symbole} & \textbf{Unité SI} & \textbf{Signification} \\
\hline
Champ magnétique & $B$ & tesla (T) & Intensité du champ créé par l'inducteur \\
\hline
Flux magnétique & $\Phi$ & weber (Wb) & « Débit » de champ à travers le circuit : $\SI{1}{Wb} = \SI{1}{T.m^2}$ \\
\hline
Inductance & $L$ & henry (H) & Aptitude d'une bobine à s'opposer aux variations de $i$ \\
\hline
F.é.m. induite & $e$ & volt (V) & Tension créée par la variation de flux : $\SI{1}{V} = \SI{1}{Wb/s}$ \\
\hline
\end{tabularx}
\end{center}

\begin{attention}[Pièges de rédaction au Probatoire]
\begin{itemize}
\item Ne pas confondre $\Phi$ en weber (Wb) et $B$ en tesla (T) : $B$ est un flux \emph{par unité de surface}.
\item Dans $e = -\dfrac{d\Phi}{dt}$, c'est la \textbf{variation} du flux qui compte, pas sa valeur : un flux énorme mais constant n'induit rien.
\item Pour une bobine, ne pas oublier le facteur $N$ dans $\Phi = NBS\cos\alpha$.
\item Dans $u = ri + L\dfrac{di}{dt}$, vérifier la convention récepteur (flèches de $u$ et $i$ en sens opposés) avant d'écrire les signes.
\end{itemize}
\end{attention}

\begin{exercice}[Pour s'entraîner]
Une bobine plate de $N = 200$ spires, d'aire $S = \SI{50}{cm^2}$, tourne à $\omega = \SI{100\pi}{rad/s}$ dans un champ $B = \SI{0,20}{T}$.
\begin{enumerate}
\item Calculer la valeur maximale de la f.é.m. induite.
\item La bobine débite $I = \SI{2,0}{A}$ (valeurs efficaces) sous $U = \SI{110}{V}$ dans un récepteur résistif. Calculer la puissance électrique et le travail mécanique minimal à fournir en une minute.
\end{enumerate}
\end{exercice}

\begin{corrige}[Exercice d'application]
\textbf{1.} $E_{\max} = NBS\omega = 200 \times 0{,}20 \times 50\times10^{-4} \times 100\pi = 20\pi \approx \SI{63}{V}$.
\textbf{2.} $P = UI = 110 \times 2{,}0 = \SI{220}{W}$ ; $W = P\,\Delta t = 220 \times 60 = \SI{13,2}{kJ}$ de travail mécanique au minimum (rendement 100 \%).
\end{corrige}

En définitive, l'induction électromagnétique est bien plus qu'un phénomène de laboratoire : c'est le cœur énergétique du monde moderne. Chaque fois que tu branches un appareil, qu'un barrage de la Sanaga tourne ou qu'une carte sans contact s'active, c'est la loi de Faraday qui travaille — et la loi de Lenz qui veille, inflexible comptable de l'énergie, à ce que rien ne soit jamais gratuit.

\newpage

\coursec{Activité d'intégration}

\coursec{Induction électromagnétique}

\courssub{Objectifs de la leçon}

À l'issue de cette leçon, tu seras capable de :

\begin{itemize}
\item définir le \cle{flux magnétique} à travers une surface et en donner l'unité SI (le weber, Wb) ;
\item énoncer et appliquer la \cle{loi de Faraday} (f.é.m. induite moyenne et instantanée) ;
\item utiliser la \cle{loi de Lenz} pour déterminer le sens du courant induit et interpréter le bilan énergétique ;
\item définir l'\cle{inductance propre} $L$ d'une bobine et calculer la \cle{f.é.m. d'auto-induction} ;
\item distinguer le \cle{circuit inducteur} (source du champ) du \cle{circuit induit} (où apparaît la f.é.m.) ;
\item modéliser un alternateur simple (groupe électrogène, dynamo, barrage hydroélectrique) et en écrire la chaîne de conversion d'énergie.
\end{itemize}

\begin{experience}Mise en évidence de l'induction électromagnétique\end{experience}

\textbf{Matériel :} Un aimant barre, une bobine de fil de cuivre (nombre de spires $N$ connu), un galvanomètre sensible (ou un multimètre en mode $\mu$A/mV), des fils de connexion.

\textbf{Protocole :}
\begin{enumerate}
\item Brancher la bobine au galvanomètre (circuit fermé, sans générateur).
\item Approcher rapidement le pôle Nord de l'aimant de la bobine (axe de l'aimant confondu avec l'axe de la bobine). Observer l'aiguille.
\item Arrêter l'aimant au contact de la bobine. Observer.
\item Éloigner rapidement l'aimant. Observer.
\item Recommencer avec le pôle Sud. Inverser le sens de branchement de la bobine.
\end{enumerate}

\textbf{Observations :}
\begin{itemize}
\item L'aiguille se dévie \textbf{uniquement pendant le mouvement} de l'aimant (variation du champ $\vec{B}$ à travers la bobine).
\item Le sens de la déviation dépend : du pôle utilisé (N ou S), du sens du mouvement (approche/éloignement), du sens d'enroulement de la bobine.
\item Aucun courant n'est détecté si l'aimant est immobile (même très fort) ou si le circuit est ouvert.
\end{itemize}

\textbf{Interprétation :} La variation du flux magnétique $\Phi$ à travers la bobine crée une \cle{force électromotrice induite} $e$ qui fait circuler un \cle{courant induit} $i$. C'est le phénomène d'\cle{induction électromagnétique}. La bobine est le \cle{circuit induit} ; l'aimant (ou le circuit parcouru par un courant variable qui crée $\vec{B}$) est le \cle{circuit inducteur}.

\begin{popfigure}
\begin{center}
\begin{tikzpicture}[scale=0.9,>=Stealth]
% Aimant vertical (N en haut, S en bas)
\fill[poppinkL] (-1.5,1.5) rectangle (1.5,2.2);
\fill[popblueL] (-1.5,-2.2) rectangle (1.5,-1.5);
\node[poppink] at (0,1.85) {\textbf{N}};
\node[popblue] at (0,-1.85) {\textbf{S}};
% Champ B vertical (flèches vers le bas)
\foreach \x in {-1,0,1}{\draw[->,popblue,thick] (\x,1.4) -- (\x,-1.4);}
\node[popblue] at (1.8,0) {$\vec{B}$};
% Bobine (rectangle vertical dans le plan de la figure)
\draw[poporange,very thick] (-1.2,-1) rectangle (1.2,1);
\node[poporange] at (1.8,1.2) {Bobine ($N$ spires)};
% Axe de rotation horizontal (traits pointillés)
\draw[dashed,popdark,thick] (-2,0) -- (2,0);
\node[popdark] at (-2.3,0) {Axe $\Delta$};
% Flèche de rotation
\draw[->,popgreen,thick] (1.5,0.6) arc (30:150:1.6);
\node[popgreen] at (0,1.5) {\small rotation};
% Galvanomètre
\draw[popdark,thick] (1.2,0) -- (2.5,0);
\draw[popdark,thick] (-1.2,0) -- (-2.5,0);
\draw[popdark,thick] (2.5,0) -- (2.5,-1) -- (3.5,-1) -- (3.5,1) -- (2.5,1) -- (2.5,0);
\node at (3,-1.4) {$G$};
\node at (3,0) {Galvanomètre};
\end{tikzpicture}
\end{center}
\legende{Expérience de référence : un aimant (circuit \cle{inducteur}) crée un champ $\vec{B}$ vertical. Une bobine (circuit \cle{induit}) tourne autour d'un axe horizontal $\Delta$ perpendiculaire à $\vec{B}$. Le flux varie $\rightarrow$ f.é.m. induite $\rightarrow$ courant détecté par $G$.}
\end{popfigure}

\begin{definition}Flux magnétique\end{definition}

Soit une surface plane $\mathcal{S}$ d'aire $S$, traversée par un champ magnétique uniforme $\vec{B}$. On munit $\mathcal{S}$ d'un vecteur normal unitaire $\vec{n}$. L'angle entre $\vec{B}$ et $\vec{n}$ est noté $\theta$.

Le \cle{flux magnétique} $\Phi$ à travers cette surface est la grandeur scalaire définie par :
\[ \Phi = B\,S\,\cos\theta = \vec{B}\cdot\vec{S} \quad \text{avec } \vec{S} = S\,\vec{n}. \]

Pour une bobine de $N$ spires identiques, toutes traversées par le même flux (bobine « plate » ou solénoïde long), le flux total (ou flux enchainé) est :
\[ \Phi = N\,B\,S\,\cos\theta. \]

\begin{aretenir}[Unité du flux]
L'unité SI du flux magnétique est le \cle{weber} (symbole \textbf{Wb}).
\[ 1~\text{Wb} = 1~\text{T}\cdot\text{m}^2 = 1~\text{V}\cdot\text{s}. \]
Analyse dimensionnelle : $[\Phi] = [B][S] = \text{T}\cdot\text{m}^2$. Homogène.
\end{aretenir}

\begin{propriete}Loi de Faraday (f.é.m. induite)\end{propriete}

Lorsqu'une variation de flux magnétique $\Delta\Phi$ se produit pendant un intervalle de temps $\Delta t$ à travers un circuit fermé (le circuit induit), une \cle{force électromotrice induite} (f.é.m.) moyenne apparaît aux bornes de ce circuit :
\[ |e| = \left|\frac{\Delta\Phi}{\Delta t}\right|. \]

En régime continu (variation continue), la f.é.m. induite instantanée est :
\[ e = -\frac{d\Phi}{dt}. \]

\begin{attention}[Convention du signe]
Le signe $-$ dans l'expression instantanée $e = -d\Phi/dt$ est la traduction mathématique de la \textbf{loi de Lenz} (voir ci-dessous). Pour la valeur moyenne, on utilise souvent la valeur absolue $|e| = |\Delta\Phi|/\Delta t$ et on détermine le sens du courant par la loi de Lenz séparément.
\end{attention}

\begin{propriete}Loi de Lenz (sens du courant induit)\end{propriete}

\textbf{Énoncé :} Le courant induit, s'il peut circuler, crée un champ magnétique propre dont l'action s'oppose à la variation de flux magnétique qui est à l'origine de l'induction.

\textbf{Conséquences énergétiques :}
\begin{itemize}
\item Le courant induit ne « naît pas de rien » : il puise l'énergie nécessaire dans le travail fourni par la cause de la variation de flux (mouvement d'un aimant, rotation d'une bobine, variation d'un courant dans un circuit voisin).
\item C'est la manifestation locale de la \textbf{conservation de l'énergie}.
\end{itemize}

\begin{methode}[Déterminer le sens du courant induit (Lenz) -- 3 étapes]
\begin{enumerate}
\item \textbf{Repérer} le sens du champ $\vec{B}$ extérieur (inducteur) et la variation du flux $\Delta\Phi$ (augmentation ou diminution ?).
\item \textbf{Déduire} le sens du champ $\vec{B}_{\text{induit}}$ créé par le courant induit : il doit \emph{s'opposer} à la variation ($\vec{B}_{\text{induit}}$ s'oppose à $\vec{B}$ si $\Phi$ augmente ; $\vec{B}_{\text{induit}}$ a le même sens que $\vec{B}$ si $\Phi$ diminue).
\item \textbf{Appliquer} la règle de la main droite (pouce = sens du courant, doigts = sens du champ au centre de la bobine) pour trouver le sens de circulation du courant induit dans les spires.
\end{enumerate}
\end{methode}

\begin{definition}Auto-induction et inductance propre\end{definition}

Lorsqu'un courant variable $i(t)$ traverse une bobine, elle crée elle-même un flux magnétique $\Phi$ proportionnel à ce courant :
\[ \Phi = L\,i. \]
Le coefficient de proportionnalité $L$ est l'\cle{inductance propre} (ou simplement \cle{inductance}) de la bobine. Elle ne dépend que de sa géométrie (nombre de spires $N$, surface $S$, longueur, présence d'un noyau ferromagnétique) et non du courant.

\begin{aretenir}[Unité de l'inductance]
L'unité SI de l'inductance est le \cle{henry} (symbole \textbf{H}).
\[ 1~\text{H} = 1~\frac{\text{Wb}}{\text{A}} = 1~\frac{\text{V}\cdot\text{s}}{\text{A}}. \]
Analyse dimensionnelle : $[L] = [\Phi]/[i] = \text{Wb}/\text{A} = \text{H}$.
\end{aretenir}

\begin{propriete}F.é.m. d'auto-induction\end{propriete}

Toute variation du courant dans une bobine entraîne une variation de son propre flux, donc l'apparition d'une f.é.m. aux bornes de cette même bobine : c'est la \cle{f.é.m. d'auto-induction}.
\[ e_{\text{auto}} = -L\,\frac{di}{dt} \quad \text{(instantanée)} \qquad \text{ou} \qquad |e_{\text{auto}}| = L\,\left|\frac{\Delta i}{\Delta t}\right| \quad \text{(moyenne)}. \]

\begin{attention}[Rôle de l'inductance -- « Inertie électrique »]
D'après la loi de Lenz, $e_{\text{auto}}$ s'oppose \emph{toujours} à la variation de courant qui la produit :
\begin{itemize}
\item Si $i$ \textbf{augmente} ($di/dt > 0$), $e_{\text{auto}} < 0$ : la bobine « freine » l'établissement du courant.
\item Si $i$ \textbf{diminue} ($di/dt < 0$), $e_{\text{auto}} > 0$ : la bobine « soutient » le courant (tend à le maintenir).
\end{itemize}
C'est cet effet qui provoque l'étincelle à l'ouverture d'un circuit inductif (forte $di/dt$ négative $\rightarrow$ forte $e_{\text{auto}}$ positive) et qui lisse le courant au démarrage.
\end{attention}

\begin{methode}Calculer un flux et une f.é.m. induite\end{methode}
\begin{enumerate}
\item Identifier $N$, $B$, $S$, et l'angle $\theta(t)$ entre $\vec{B}$ et la normale $\vec{n}$.
\item Écrire $\Phi(t) = N B S \cos\theta(t)$.
\item Pour une f.é.m. \textbf{moyenne} sur $\Delta t$ : calculer $\Delta\Phi = \Phi(t_2) - \Phi(t_1)$ et appliquer $|e| = |\Delta\Phi|/\Delta t$.
\item Pour une f.é.m. \textbf{instantanée} : dériver $\Phi(t)$ par rapport au temps : $e = -d\Phi/dt$.
\item Vérifier l'homogénéité des unités (SI : T, m², s $\rightarrow$ V).
\end{enumerate}

\courssub{Activité d'intégration : Le groupe électrogène de Nkolndongo}

\textbf{Situation}

Dans le quartier Nkolndongo à Yaoundé, les délestages fréquents poussent M. Etoundi, gérant d'un atelier de couture, à installer un petit groupe électrogène. Le principe est simple : un moteur à essence entraîne en rotation une bobine placée entre les pôles d'un aimant puissant. C'est l'induction électromagnétique qui transforme le travail mécanique du moteur en énergie électrique.

M. Etoundi confie à ta classe de Première C la mission suivante : vérifier, par le calcul, que la bobine de son groupe peut alimenter correctement une ampoule, et expliquer pourquoi le moteur « peine » et consomme davantage d'essence lorsqu'on branche plus d'appareils.

\textbf{Modélisation du dispositif.} On modélise l'alternateur du groupe par une \cle{bobine plate} de $N$ spires identiques, chacune de surface $S$, tournant à vitesse constante autour d'un axe perpendiculaire aux lignes d'un champ magnétique uniforme $\vec{B}$.

\textbf{Données numériques :}
\begin{itemize}
\item surface d'une spire : $S = 40~\text{cm}^2 = \num{4,0e-3}~\text{m}^2$ ;
\item nombre de spires : $N = 100$ ;
\item intensité du champ magnétique : $B = \num{0,80}~\text{T}$ ;
\item fréquence de rotation : $f = 50$ tours par seconde, donc période $T = \dfrac{1}{f} = \SI{20}{ms}$ ;
\item durée d'un demi-tour : $\Delta t = \dfrac{T}{2} = \SI{10}{ms} = \num{1,0e-2}~\text{s}$ ;
\item résistance de l'ampoule alimentée : $R = \SI{44}{\ohm}$ (on néglige la résistance interne de la bobine) ;
\item inductance propre de la bobine : $L = \SI{50}{mH} = \num{5,0e-2}~\text{H}$ ;
\item variation de courant au démarrage : $\Delta i = \num{3,0}~\text{A}$ en $\Delta t' = \SI{5,0}{ms}$.
\end{itemize}

\begin{popfigure}
\begin{center}
\begin{tikzpicture}[scale=0.95,>=Stealth]
% Aimant : pôles horizontal (N gauche, S droit) -> B horizontal vers la droite
\fill[poppinkL] (-3.8,-1.5) rectangle (-2.4,1.5);
\fill[popblueL] (2.4,-1.5) rectangle (3.8,1.5);
\node[poppink] at (-3.1,0) {\textbf{N}};
\node[popblue] at (3.1,0) {\textbf{S}};
% Champ B : flèches horizontales
\foreach \y in {-1.2,-0.6,0,0.6,1.2}{
  \draw[->,popblue,thick] (-2.3,\y) -- (-0.8,\y);
  \draw[->,popblue,thick] (0.8,\y) -- (2.3,\y);
}
\node[popblue] at (0,1.8) {$\vec{B}$};
% Bobine VERTICALE (vue de côté : rectangle dans le plan yz -> ici plan xy de l'écran)
% Axe de rotation horizontal (perpendiculaire à B et au plan de la figure -> axe z sortant de l'écran)
% Mais pour faire simple en 2D : on dessine la bobine comme un rectangle vertical qui tourne autour d'un axe horizontal (pointillés verticaux ? Non)
% Standard : B vertical, Axe horizontal. Ou B horizontal, Axe vertical, Bobine verticale.
% Choix : B horizontal (x). Axe vertical (y). Bobine verticale (plan xz) -> vue de côté = rectangle vertical (plan xy).
% Axe = centre du rectangle, vertical.
\draw[dashed,popdark,thick] (0,-2.2) -- (0,2.2);
\node[popdark] at (0,2.5) {Axe $\Delta$ (vertical)};
% Bobine : rectangle vertical centré, tourné d'un angle alpha (ici 30 deg) autour de l'axe vertical (centre)
% En 2D (plan écran = plan horizontal), une bobine verticale vue de dessus est un segment.
% Mieux : Vue en perspective (isométrique) ou Vue de face (plan vertical).
% Faisons une VUE DE FACE (plan vertical xz -> xy sur écran).
% B horizontal (dans le plan de l'écran, vers la droite).
% Axe horizontal (perpendiculaire à l'écran, vers nous).
% Bobine : rectangle dans le plan de l'écran, tournant autour de son centre (axe horizontal sortant).
% Aimant : N en haut, S en bas -> B vertical. C'est plus standard pour un dessin 2D.
% REFAIT PROPREMENT : B VERTICAL, AXE HORIZONTAL.
\end{tikzpicture}
\end{center}
\legende{Schéma corrigé : Champ $\vec{B}$ vertical (N haut, S bas). Bobine verticale tournant autour d'un axe horizontal $\Delta$ (perpendiculaire à $\vec{B}$). Le flux $\Phi = NBS\cos(\omega t)$ varie sinusoïdalement.}
\end{popfigure}

% Re-drawing the figure properly in the final code block below.
% For now, continuing the text structure.

\courssub{Consignes}

\begin{enumerate}
\item \textbf{Flux maximal.} Exprimer le flux magnétique maximal $\Phi_{\max}$ à travers la bobine de $N$ spires en fonction de $N$, $B$ et $S$. Le calculer numériquement pour $N = 100$ spires, en précisant l'unité SI.
\item \textbf{F.é.m. induite moyenne.} Pendant un demi-tour, le flux passe de $+\Phi_{\max}$ à $-\Phi_{\max}$. Calculer la variation de flux $\Delta\Phi$, puis la valeur moyenne de la f.é.m. induite $|e| = \left|\dfrac{\Delta\Phi}{\Delta t}\right|$.
\item \textbf{Courant induit et loi de Lenz.} La bobine alimente l'ampoule de résistance $R = \SI{44}{\ohm}$. Calculer l'intensité moyenne $I$ du courant induit. En invoquant la loi de Lenz, expliquer pourquoi le courant induit s'oppose à la rotation de la bobine, et en déduire pourquoi le moteur à essence doit « forcer » davantage (et consommer plus) quand on branche plus d'appareils.
\item \textbf{Auto-induction au démarrage.} Au démarrage, le courant varie de $\Delta i = \num{3,0}~\text{A}$ en $\Delta t' = \SI{5,0}{ms}$. Calculer la valeur absolue de la f.é.m. d'auto-induction $|e_{\text{auto}}|$ dans la bobine. Conclure sur le rôle de la bobine (effet de « lissage » du courant) au démarrage du groupe.
\item \textbf{Rapport.} Rédiger un court rapport (une demi-page) à destination de M. Etoundi, comportant : le schéma du dispositif, les résultats numériques obtenus, et la chaîne de conversion d'énergie du groupe électrogène.
\end{enumerate}

\courssub{Ressources à mobiliser}

\begin{aretenir}[Formulaire de la leçon]
\begin{itemize}
\item \textbf{Flux magnétique} (spire plane, champ uniforme) : $\Phi = B\,S\,\cos\theta = \vec{B}\cdot\vec{S}$. Pour $N$ spires : $\Phi = N\,B\,S\,\cos\theta$. Unité : \textbf{Wb} ($1~\text{Wb} = 1~\text{T}\cdot\text{m}^2$).
\item \textbf{Loi de Faraday} : f.é.m. induite moyenne $|e| = \left|\dfrac{\Delta\Phi}{\Delta t}\right|$ ; instantanée $e = -\dfrac{d\Phi}{dt}$.
\item \textbf{Loi de Lenz} : Le courant induit s'oppose, par ses effets, à la variation de flux qui lui a donné naissance. (Conservation de l'énergie).
\item \textbf{Loi d'Ohm} (circuit induit fermé sur $R$) : $I = \dfrac{|e|}{R}$.
\item \textbf{F.é.m. d'auto-induction} : $e_{\text{auto}} = -L\,\dfrac{di}{dt}$, $L$ en \textbf{H} ($1~\text{H} = 1~\text{V}\cdot\text{s}/\text{A}$).
\end{itemize}
\end{aretenir}

\courssub{Démarche et corrigé commenté}

\begin{exoresolu}[Dimensionnement de la bobine du groupe électrogène]
Reprendre les cinq consignes de l'activité avec les données : $S = \num{4,0e-3}~\text{m}^2$, $N = 100$, $B = \num{0,80}~\text{T}$, $\Delta t = \SI{10}{ms}$, $R = \SI{44}{\ohm}$, $L = \SI{50}{mH}$, $\Delta i = \num{3,0}~\text{A}$ en $\SI{5,0}{ms}$.
\tcblower
\textbf{1) Flux maximal.}

Le flux à travers la bobine est $\Phi = N B S \cos\theta$. Il est maximal lorsque la normale aux spires est parallèle au champ ($\theta = 0$, $\cos\theta = 1$) :
\[ \Phi_{\max} = N\,B\,S = 100 \times \num{0,80} \times \num{4,0e-3} = \num{0,32}~\text{Wb}. \]
Vérification dimensionnelle : $[\text{T}]\times[\text{m}^2] = \text{Wb}$, homogène à un flux. \textbf{Le flux maximal vaut $\Phi_{\max} = \num{0,32}~\text{Wb}$.}

\textbf{2) F.é.m. induite moyenne.}

En un demi-tour, la bobine se retourne : le flux passe de $+\Phi_{\max}$ à $-\Phi_{\max}$. La variation de flux est donc
\[ \Delta\Phi = \Phi_{\text{final}} - \Phi_{\text{initial}} = -\Phi_{\max} - (+\Phi_{\max}) = -2\,\Phi_{\max} = -\num{0,64}~\text{Wb}. \]
La loi de Faraday donne alors la valeur moyenne de la f.é.m. induite :
\[ |e| = \left|\dfrac{\Delta\Phi}{\Delta t}\right| = \dfrac{2\,\Phi_{\max}}{\Delta t} = \dfrac{\num{0,64}}{\num{1,0e-2}} = \num{64}~\text{V}. \]
\textbf{La f.é.m. moyenne induite vaut $|e| = \num{64}~\text{V}$.} C'est une valeur moyenne sur un demi-période : en réalité, la f.é.m. est alternative sinusoïdale ($e(t) = \Phi_{\max}\,\omega\sin(\omega t)$), mais cette moyenne suffit pour dimensionner le circuit en régime établi.

\textbf{3) Courant induit et loi de Lenz.}

Le circuit induit est fermé sur l'ampoule de résistance $R = \SI{44}{\ohm}$. La loi d'Ohm donne :
\[ I = \dfrac{|e|}{R} = \dfrac{\num{64}}{\num{44}} \approx \num{1,45}~\text{A}. \]
\textbf{L'intensité moyenne du courant induit vaut environ $\num{1,5}~\text{A}$ (2 chiffres significatifs).}

\textit{Interprétation par la loi de Lenz.} Le courant induit, parcourant la bobine placée dans le champ $\vec{B}$, subit des forces de Laplace ($\vec{F} = I\,\vec{l}\wedge\vec{B}$). D'après la loi de Lenz, ces forces s'opposent à la cause de l'induction, c'est-à-dire à la rotation de la bobine : elles créent un \textbf{couple résistant} (ou couple de freinage électromagnétique) qui s'oppose au couple moteur. Pour maintenir la vitesse de rotation constante ($f = \SI{50}{Hz}$), le moteur à essence doit fournir un travail mécanique supplémentaire pour vaincre ce couple résistant. Brancher plus d'appareils revient à diminuer la résistance équivalente $R_{\text{eq}}$ du circuit induit, ce qui \textbf{augmente le courant induit} ($I = |e|/R_{\text{eq}}$), donc \textbf{augmente le couple résistant} : \textbf{le moteur doit « forcer » davantage et consomme plus d'essence}. C'est exactement la traduction de la conservation de l'énergie : l'énergie électrique fournie aux appareils ($P_{\text{élec}} = |e|\,I$) provient du travail mécanique du moteur ($P_{\text{méca}} = \mathcal{C}_{\text{résist}}\,\omega$), lui-même issu de la combustion de l'essence.

\textbf{4) Auto-induction au démarrage.}

La bobine, parcourue par un courant variable, est le siège d'une f.é.m. d'auto-induction :
\[ |e_{\text{auto}}| = L\,\left|\dfrac{\Delta i}{\Delta t'}\right| = \num{5,0e-2} \times \dfrac{\num{3,0}}{\num{5,0e-3}} = \num{5,0e-2} \times \num{600} = \num{30}~\text{V}. \]
\textbf{La f.é.m. d'auto-induction vaut $\num{30}~\text{V}$ (soit $\num{3,0}\times 10^1~\text{V}$).} D'après la loi de Lenz, elle s'oppose à l'établissement brutal du courant : la bobine « retarde » et « lisse » les variations de courant (effet d'inertie électrique). Au démarrage du groupe, cela protège l'ampoule et les appareils contre les à-coups de courant ; à l'ouverture brutale du circuit, c'est elle qui peut provoquer une étincelle de rupture aux bornes de l'interrupteur (forte $di/dt$ négative $\rightarrow$ forte $e_{\text{auto}}$ positive).

\textbf{5) Chaîne de conversion d'énergie (rapport).}
\[ \text{Énergie chimique (essence)} \xrightarrow{\text{moteur}} \text{Énergie mécanique (rotation)} \xrightarrow{\text{induction}} \text{Énergie électrique (alternative)} \xrightarrow{\text{ampoule}} \text{Énergie lumineuse + thermique}. \]
Le rapport à M. Etoundi conclut que la bobine délivre une tension moyenne de $\num{64}~\text{V}$, suffisante pour alimenter l'ampoule avec un courant d'environ $\num{1,5}~\text{A}$, et que toute augmentation du nombre d'appareils branchés se traduit par une surconsommation d'essence proportionnelle à la puissance électrique demandée.
\end{exoresolu}

\begin{attention}[Pièges fréquents au Probatoire]
\begin{itemize}
\item \textbf{Conversion d'unités :} $40~\text{cm}^2 = \num{4,0e-3}~\text{m}^2$ (division par $10^4$, pas par $100$). $50~\text{mH} = \num{5,0e-2}~\text{H}$. $10~\text{ms} = \num{1,0e-2}~\text{s}$. Toujours travailler en SI (kg, m, s, A) avant de calculer.
\item \textbf{Variation de flux sur un demi-tour :} $\Delta\Phi = -2\Phi_{\max}$ (valeur absolue $2\Phi_{\max}$), \textbf{pas} $\Phi_{\max}$. Le flux change de signe car la normale $\vec{n}$ s'inverse.
\item \textbf{Facteur $N$ :} Le flux $\Phi$ et la f.é.m. $e$ sont proportionnels au nombre de spires $N$. Oublier $N$ donne un résultat $100$ fois trop petit.
\item \textbf{Sens du courant (Lenz) :} Ne pas confondre « s'opposer à la variation du flux » et « s'opposer au champ ». Si le flux diminue, le courant induit crée un champ \textbf{de même sens} que $\vec{B}$ pour tenter de le maintenir.
\item \textbf{Auto-induction :} La f.é.m. d'auto-induction existe \textbf{uniquement} si le courant varie ($di/dt \neq 0$). En régime continu ($i$ constant), $e_{\text{auto}} = 0$ : la bobine se comporte comme un fil parfait (résistance nulle ici).
\end{itemize}
\end{attention}

\begin{savaistu}[L'induction au Cameroun : des barrages aux mototaxis]
\begin{itemize}
\item \textbf{Barrages hydroélectriques (Lom Pangar, Nachtigal, Edéa, Song Loulou) :} De gigantesques alternateurs synchrones transforment l'énergie potentielle de l'eau (barrage) $\rightarrow$ énergie cinétique (turbine) $\rightarrow$ énergie mécanique (rotation de l'alternateur) $\rightarrow$ énergie électrique (induction dans les statoriques triphasés). La puissance installée de Nachtigal est de $\approx \SI{420}{MW}$.
\item \textbf{Transport d'énergie (SONATREL/ENEO) :} Le courant alternatif $\SI{50}{Hz}$ permet d'utiliser des \textbf{transformateurs} (induction mutuelle) pour élever la tension ($\SI{225}{kV}$, $\SI{400}{kV}$) et réduire les pertes par effet Joule ($P = R I^2$) sur des centaines de km.
\item \textbf{Mototaxis (bendskin) :} Le démarreur et l'alternateur de la moto sont des machines à induction. La bobine d'allumage (transformateur élévateur) utilise l'auto-induction brutale ($di/dt$ énorme à l'ouverture du rupteur) pour générer $\sim \SI{20}{kV}$ aux bougies.
\item \textbf{Télécommunications (MTN, Orange, CAMTEL) :} Les antennes relais émettent/réçoivent des ondes électromagnétiques : induction mutuelle entre l'antenne émettrice (courant HF variable) et l'antenne réceptrice (f.é.m. induite).
\item \textbf{Hôpitaux (radiothérapie, IRM) :} Les accélérateurs de particules et les aimants supraconducteurs (IRM, champ $\sim \SI{1,5}{T}$ à $\SI{3}{T}$) reposent sur des bobines à forte inductance ($L$ grand) traversées par des courants constants (régime continu, $e_{\text{auto}}=0$) mais dont la mise sous tension/hors tension demande une gestion fine de $di/dt$.
\end{itemize}
\end{savaistu}

\courssub{Bilan de la leçon}

Cette leçon t'a permis de maîtriser l'ensemble des savoir-faire du Module 4 « Aspects énergétiques des circuits électriques » pour la série 1ère C :

\begin{itemize}
\item le \textbf{calcul du flux magnétique} ($\Phi = NBS\cos\theta$, unité Wb), en veillant scrupuleusement aux unités SI ;
\item l'application de la \textbf{loi de Faraday} pour obtenir une f.é.m. induite (moyenne $|e| = |\Delta\Phi|/\Delta t$ ou instantanée $e = -d\Phi/dt$), en distinguant clairement le circuit \cle{inducteur} (source de $\vec{B}$) du circuit \cle{induit} (siège de $e$) ;
\item l'interprétation physique et énergétique de la \textbf{loi de Lenz} : le courant induit s'oppose toujours à la cause qui le produit (conservation de l'énergie), ce qui explique concrètement le couple résistant dans un alternateur et la surconsommation de carburant sous charge ;
\item la définition de l'\textbf{inductance propre} $L = \Phi/i$ (unité H) et le calcul de la \textbf{f.é.m. d'auto-induction} ($e_{\text{auto}} = -L\,di/dt$), avec la compréhension du rôle de « régulateur/inertie » de la bobine face aux variations brutales de courant ;
\item enfin, la construction de \textbf{chaînes de conversion d'énergie} complètes (essence $\rightarrow$ mécanique $\rightarrow$ électrique $\rightarrow$ lumière/chaleur ; eau $\rightarrow$ mécanique $\rightarrow$ électrique $\rightarrow$ réseau) qui donnent tout leur sens physique au phénomène d'induction : il n'y a pas de création d'énergie, seulement des conversions réversibles (moteur $\leftrightarrow$ générateur).
\end{itemize}

Tu es désormais capable d'expliquer le fonctionnement d'un alternateur, qu'il s'agisse du groupe électrogène du quartier Nkolndongo, de la dynamo d'une mototaxi, ou des grands alternateurs des barrages de Lom Pangar et Nachtigal qui alimentent le réseau national géré par SONATREL et ENEO.

\newpage

\coursec{Méthodes et exercices résolus}

\coursec{Induction électromagnétique}

\courssub{Mise en évidence expérimentale de l'induction}

\begin{experience}[Induction par aimant mobile et par variation de courant]
\textbf{Matériel :} Aimant droit, bobine de nombreuses spires ($N \approx 500$), galvanomètre à zéro central (ou milliammètre sensible), fils de connexion. Pour la deuxième partie : deux bobines $B_1$ et $B_2$ placées face à face, générateur de tension continue, rhéostat, interrupteur, galvanomètre.
\textbf{Protocole 1 (aimant + bobine) :}
\begin{enumerate}
\item Relier la bobine au galvanomètre (circuit fermé, sans générateur).
\item Approcher rapidement le pôle nord de l'aimant de la face de la bobine, observer l'aiguille.
\item Maintenir l'aimant immobile devant la bobine, observer.
\item Éloigner rapidement l'aimant, observer.
\item Recommencer avec le pôle sud.
\end{enumerate}
\textbf{Protocole 2 (deux bobines) :}
\begin{enumerate}
\item Monter $B_1$ en série avec le générateur, le rhéostat et l'interrupteur (circuit inducteur).
\item Relier $B_2$ au galvanomètre (circuit induit, sans générateur).
\item Fermer l'interrupteur : observer la déviation brève.
\item Laisser l'interrupteur fermé : observer.
\item Ouvrir l'interrupteur : observer la déviation brève en sens inverse.
\item Varier le rhéostat avec l'interrupteur fermé : observer des déviations continues.
\end{enumerate}
\textbf{Observations :} Le galvanomètre du circuit induit ne dévie \emph{que pendant les variations} (mouvement de l'aimant, fermeture/ouverture/varition du courant dans $B_1$). Aucun courant permanent n'est observé pour un aimant immobile ou un courant continu constant dans $B_1$.
\textbf{Interprétation :} Un \cle{courant induit} apparaît dans un circuit fermé (circuit induit) lorsque le \cle{flux magnétique} $\Phi$ du champ magnétique $\vec{B}$ (créé par le circuit inducteur : aimant ou bobine $B_1$) à travers ce circuit \emph{varie}. C'est le phénomène d'\cle{induction électromagnétique}.
\end{experience}

\begin{methode}[Reconnaître et interpréter une expérience d'induction]
Pour analyser une expérience d'induction électromagnétique (aimant + bobine, ou deux bobines) :
\begin{enumerate}
\item \textbf{Repérer les deux circuits} : le \cle{circuit inducteur} (source du champ magnétique variable : aimant en mouvement, ou électroaimant alimenté par un courant variable) et le \cle{circuit induit} (bobine fermée sur un galvanomètre, \emph{sans générateur}).
\item \textbf{Identifier la cause de la variation} : déplacement de l'aimant, déplacement de la bobine, variation du courant dans l'inducteur (fermeture/ouverture d'un interrupteur, variation du rhéostat).
\item \textbf{Observer} : le galvanomètre dévie \emph{uniquement pendant la variation} ; si rien ne varie (aimant immobile, courant constant), aucun courant induit n'apparaît.
\item \textbf{Conclure} : un courant induit naît dans un circuit fermé chaque fois que le flux magnétique à travers ce circuit \emph{varie}.
\end{enumerate}
\textbf{Réflexe de rédaction} : toujours préciser « circuit fermé » et « variation de flux » dans la conclusion ; un champ constant, même intense, n'induit rien.
\end{methode}

\begin{exoresolu}[Aimant et bobine]
On approche le pôle nord d'un aimant droit d'une bobine de $N = 500$ spires reliée à un galvanomètre à zéro central. Le circuit induit est fermé et ne contient aucun générateur.
\begin{enumerate}
\item Décrire ce qu'observe l'expérimentateur pendant l'approche, puis lorsque l'aimant est immobile devant la bobine.
\item Identifier le circuit inducteur et le circuit induit.
\item Que se passe-t-il si on éloigne l'aimant ? Conclure.
\end{enumerate}
\tcblower
\textbf{Solution.}
\begin{enumerate}
\item Pendant l'approche de l'aimant, le champ magnétique $\vec{B}$ au niveau de la bobine augmente : le flux à travers les spires varie. Le galvanomètre dévie dans un sens : un \cle{courant induit} circule. Lorsque l'aimant est immobile, le flux ne varie plus : l'aiguille revient à zéro, il n'y a plus de courant.
\item Le \textbf{circuit inducteur} est l'aimant (source du champ magnétique). Le \textbf{circuit induit} est la bobine de $N = 500$ spires fermée sur le galvanomètre : c'est en son sein que naît la f.é.m. induite.
\item En éloignant l'aimant, le flux diminue : le galvanomètre dévie \emph{dans le sens opposé}. \textbf{Conclusion} : un courant induit n'apparaît dans un circuit fermé que lorsque le flux magnétique qui le traverse varie ; le sens du courant dépend du sens de la variation du flux.
\end{enumerate}
\end{exoresolu}

\courssub{Flux magnétique}

\begin{definition}[Flux magnétique à travers une surface]
Soit une surface plane $S$ traversée par un champ magnétique \cle{uniforme} $\vec{B}$. On munit la surface d'un vecteur normal $\vec{n}$ (lié au sens positif de parcours du circuit par la règle du tire-bouchon). L'angle entre $\vec{n}$ et $\vec{B}$ est noté $\alpha$.
Le \cle{flux magnétique} de $\vec{B}$ à travers cette surface est la grandeur scalaire :
\[ \Phi = B\,S\cos\alpha \]
Pour une bobine de $N$ spires identiques et parfaitement superposées (bobine plate), le flux total est :
\[ \Phi = N\,B\,S\cos\alpha \]
\end{definition}

\begin{propriete}[Unité du flux magnétique]
L'unité SI du flux magnétique est le \cle{weber} (symbole \textbf{Wb}).
\[ 1\,\mathrm{Wb} = 1\,\mathrm{T}\cdot\mathrm{m}^2 = 1\,\mathrm{V}\cdot\mathrm{s} \]
\textbf{Analyse dimensionnelle :} $[\Phi] = [B][S] = \mathrm{T}\cdot\mathrm{m}^2$. Puisque $1\,\mathrm{T} = 1\,\mathrm{N}\cdot\mathrm{A}^{-1}\cdot\mathrm{m}^{-1} = 1\,\mathrm{V}\cdot\mathrm{s}\cdot\mathrm{m}^{-2}$, on a bien $[\Phi] = \mathrm{V}\cdot\mathrm{s}$.
\end{propriete}

\begin{aretenir}[Cas particuliers du flux]
\begin{itemize}
\item $\vec{B} \perp$ plan de la spire ($\vec{n} \parallel \vec{B}$, $\alpha = 0$) : $\Phi = NBS$ (flux maximal, positif si $\vec{n}$ et $\vec{B}$ même sens).
\item $\vec{B}$ dans le plan de la spire ($\vec{n} \perp \vec{B}$, $\alpha = 90^\circ$) : $\Phi = 0$.
\item $\vec{B} \perp$ plan mais $\vec{n}$ opposé à $\vec{B}$ ($\alpha = 180^\circ$) : $\Phi = -NBS$ (flux minimal, négatif).
\end{itemize}
Le flux est une grandeur \emph{algébrique} : son signe dépend de l'orientation choisie pour $\vec{n}$ (sens positif du courant).
\end{aretenir}

\begin{methode}[Calcul du flux magnétique]
\begin{enumerate}
\item Vérifier que le champ $\vec{B}$ est \cle{uniforme} sur la surface et que la surface est \emph{plane}.
\item Orienter la surface : choisir un vecteur normal $\vec{n}$ (lié au sens positif choisi sur le circuit par la règle du tire-bouchon).
\item Déterminer l'angle $\alpha = (\vec{n}, \vec{B})$.
\item Appliquer, pour une spire : $\Phi = B\,S\cos\alpha$, et pour une bobine de $N$ spires :
\[ \Phi = N\,B\,S\cos\alpha \]
\item Vérifier l'homogénéité : $[\Phi] = \mathrm{T}\times\mathrm{m}^2 = \mathrm{Wb}$ (weber).
\end{enumerate}
\textbf{Attention} : $\alpha$ est l'angle entre $\vec{B}$ et la \emph{normale} $\vec{n}$ à la surface, \textbf{pas} l'angle entre $\vec{B}$ et le plan de la spire. Si l'énoncé donne l'angle $\beta$ avec le plan, alors $\alpha = 90^\circ - \beta$ et $\cos\alpha = \sin\beta$.
\end{methode}

\begin{exoresolu}[Flux à travers une bobine]
Une bobine plate de $N = 200$ spires circulaires de rayon $r = \SI{5,0}{cm}$ est placée dans un champ magnétique uniforme $B = \SI{0,40}{T}$. Le vecteur normal au plan des spires fait un angle $\alpha = 60^\circ$ avec $\vec{B}$.
\begin{enumerate}
\item Calculer le flux de $\vec{B}$ à travers la bobine.
\item On fait tourner la bobine jusqu'à ce que son plan soit perpendiculaire à $\vec{B}$. Calculer le nouveau flux et la variation de flux $\Delta\Phi$.
\end{enumerate}
\tcblower
\textbf{Solution.}
\begin{enumerate}
\item Surface d'une spire :
\[ S = \pi r^2 = \pi \times (\num{5,0}\times 10^{-2})^2 = \num{7,85}\times 10^{-3}\ \mathrm{m}^2 \]
Flux à travers la bobine :
\[ \Phi = N B S \cos\alpha = 200 \times \num{0,40} \times \num{7,85}\times 10^{-3} \times \cos 60^\circ \]
\[ \Phi = 200 \times \num{0,40} \times \num{7,85}\times 10^{-3} \times \num{0,50} = \num{3,1}\times 10^{-1}\ \mathrm{Wb} \]
\item Plan perpendiculaire à $\vec{B}$ signifie $\vec{n}$ parallèle à $\vec{B}$, donc $\alpha' = 0$ :
\[ \Phi' = NBS = 200 \times \num{0,40} \times \num{7,85}\times 10^{-3} = \num{6,3}\times 10^{-1}\ \mathrm{Wb} \]
Variation de flux :
\[ \Delta\Phi = \Phi' - \Phi = \num{0,63} - \num{0,31} = \num{3,1}\times 10^{-1}\ \mathrm{Wb} \]
\textbf{Conclusion} : le flux double lorsque la bobine se place face au champ ; cette variation de flux est la source d'une f.é.m. induite pendant la rotation.
\end{enumerate}
\end{exoresolu}

\courssub{Loi de Lenz et sens du courant induit}

\begin{propriete}[Loi de Lenz (1834)]
Le sens du courant induit dans un circuit fermé est tel que le champ magnétique $\vec{B}_i$ qu'il crée (appelé \cle{champ induit}) \textbf{s'oppose à la variation du flux magnétique} $\Phi$ qui l'a engendré.
\begin{itemize}
\item Si le flux $\Phi$ \textbf{augmente} ($d\Phi/dt > 0$) : $\vec{B}_i$ est de sens \emph{opposé} à $\vec{B}$ (il tente de diminuer le flux).
\item Si le flux $\Phi$ \textbf{diminue} ($d\Phi/dt < 0$) : $\vec{B}_i$ est de \emph{même sens} que $\vec{B}$ (il tente de maintenir le flux).
\end{itemize}
\textbf{Signification physique :} La loi de Lenz est une conséquence de la \textbf{conservation de l'énergie}. Les effets du courant induit (force de Laplace, champ induit) s'opposent toujours à la cause qui lui a donné naissance (mouvement de l'aimant, variation de courant). Il faut fournir un travail mécanique (ou électrique) pour maintenir la variation de flux ; ce travail se transforme en énergie électrique (effet Joule) dans le circuit induit.
\end{propriete}

\begin{methode}[Appliquer la loi de Lenz pour trouver le sens du courant induit]
\begin{enumerate}
\item Déterminer le sens du champ inducteur $\vec{B}$ à travers le circuit induit (lignes de champ : sortent du Nord, entrent dans le Sud).
\item Dire si le flux $\Phi$ \emph{augmente} ou \emph{diminue} (en valeur algébrique par rapport au sens positif choisi).
\item Appliquer la loi de Lenz pour trouver le sens du \cle{champ induit} $\vec{B}_i$ :
\begin{itemize}
\item $\Phi$ augmente $\Rightarrow$ $\vec{B}_i$ opposé à $\vec{B}$.
\item $\Phi$ diminue $\Rightarrow$ $\vec{B}_i$ même sens que $\vec{B}$.
\end{itemize}
\item En déduire le sens du courant induit par la \textbf{règle du tire-bouchon} (pouce = sens de $\vec{B}_i$ à l'intérieur de la bobine, doigts = sens du courant).
\end{enumerate}
\textbf{Réflexe} : Dessiner toujours un schéma avec $\vec{B}$, $\vec{B}_i$, le sens positif (flèche sur le fil) et le sens réel du courant induit.
\end{methode}

\begin{exoresolu}[Sens du courant induit]
On approche le pôle nord d'un aimant de la face d'une spire circulaire fermée sur elle-même.
\begin{enumerate}
\item Le flux à travers la spire augmente-t-il ou diminue-t-il ?
\item Déterminer le sens du champ induit $\vec{B}_i$ créé par le courant induit.
\item En déduire le sens du courant induit vu par l'observateur qui regarde l'aimant depuis la spire, et la nature (nord ou sud) de la face de la spire côté aimant.
\end{enumerate}
\tcblower
\textbf{Solution.}
\begin{enumerate}
\item En approchant le pôle nord, le champ $\vec{B}$ (dirigé du nord vers le sud, donc \emph{vers la spire}) devient plus intense : le flux $\Phi$ \textbf{augmente}.
\item D'après la loi de Lenz, le champ induit s'oppose à cette augmentation : $\vec{B}_i$ est de sens \textbf{opposé} à $\vec{B}$, c'est-à-dire dirigé \emph{vers l'aimant} (de la spire vers l'aimant).
\item Par la règle du tire-bouchon, pour créer un champ $\vec{B}_i$ dirigé vers l'aimant (à l'intérieur de la spire), le courant induit circule dans le \textbf{sens inverse des aiguilles d'une montre} (sens trigonométrique) vu par l'observateur placé côté aimant. La face de la spire côté aimant se comporte alors comme un pôle \textbf{nord} : elle \emph{repousse} l'aimant qu'on approche, ce qui illustre bien l'opposition à la cause (il faut fournir un travail pour avancer l'aimant).
\end{enumerate}
\textbf{Conclusion} : le courant induit s'oppose, par ses effets, au déplacement de l'aimant : c'est la loi de Lenz.
\end{exoresolu}

\courssub{F.é.m. induite : loi de Faraday}

\begin{propriete}[Loi de Faraday (1831) — F.é.m. d'induction]
La force électromotrice (f.é.m.) $e$ induite dans un circuit fermé est égale à l'opposé de la dérivée temporelle du flux magnétique total $\Phi$ qui traverse ce circuit :
\[ e = -\dfrac{d\Phi}{dt} \]
Si la variation du flux est régulière sur une durée $\Delta t$ :
\[ e = -\dfrac{\Delta\Phi}{\Delta t} \]
Le signe « $-$ » traduit la loi de Lenz (opposition à la variation). En valeur absolue (souvent demandée pour l'intensité du courant) :
\[ |e| = \left|\dfrac{\Delta\Phi}{\Delta t}\right| \]
\textbf{Unité :} $\mathrm{Wb}/\mathrm{s} = \mathrm{V}$ (volt).
\end{propriete}

\begin{methode}[Calculer une f.é.m. induite (loi de Faraday)]
\begin{enumerate}
\item Exprimer le flux $\Phi(t)$ à travers le circuit induit (attention au nombre de spires $N$ : $\Phi = NBS\cos\alpha$).
\item Identifier la variation : $\Delta\Phi = \Phi_{\text{final}} - \Phi_{\text{initial}}$ sur $\Delta t$.
\item Appliquer la loi de Faraday :
\[ e = -\dfrac{\Delta\Phi}{\Delta t} \quad \text{(variation régulière)} \qquad \text{ou} \qquad e = -\dfrac{d\Phi}{dt} \quad \text{(général)} \]
\item Le signe « $-$ » traduit la loi de Lenz ; si l'on demande une \emph{valeur} ou une \emph{intensité}, donner $|e|$ en volts.
\item Si le circuit induit est fermé et de résistance totale $R$, le courant induit vaut (loi d'Ohm) :
\[ i = \dfrac{|e|}{R} \]
\item Vérifier les unités : $\mathrm{Wb}/\mathrm{s} = \mathrm{V}$. Convertir systématiquement les durées en secondes ($\SI{50}{ms} = \SI{5,0e-2}{s}$) et les surfaces en $\mathrm{m}^2$.
\end{enumerate}
\end{methode}

\begin{exoresolu}[F.é.m. induite dans une bobine]
Une bobine de $N = 100$ spires, de section totale équivalente $S = \SI{2,0e-3}{m^2}$, est placée dans un champ uniforme perpendiculaire à son plan. Le champ passe de $B_1 = \SI{0,10}{T}$ à $B_2 = \SI{0,50}{T}$ en $\Delta t = \SI{0,20}{s}$, de façon régulière. La résistance totale du circuit induit est $R = \SI{5,0}{\Omega}$.
\begin{enumerate}
\item Calculer la valeur absolue de la f.é.m. induite.
\item Calculer l'intensité du courant induit.
\item Calculer la quantité d'électricité ayant traversé le circuit pendant $\Delta t$.
\end{enumerate}
\tcblower
\textbf{Solution.}
\begin{enumerate}
\item Variation de flux (le champ est perpendiculaire au plan, $\cos\alpha = 1$) :
\[ \Delta\Phi = N S\,\Delta B = 100 \times \num{2,0}\times 10^{-3} \times (\num{0,50} - \num{0,10}) = \num{8,0}\times 10^{-2}\ \mathrm{Wb} \]
Loi de Faraday en valeur absolue :
\[ |e| = \dfrac{|\Delta\Phi|}{\Delta t} = \dfrac{\num{8,0}\times 10^{-2}}{\num{0,20}} = \SI{0,40}{V} \]
\item Intensité du courant induit :
\[ i = \dfrac{|e|}{R} = \dfrac{\num{0,40}}{\num{5,0}} = \SI{8,0e-2}{A} = \SI{80}{mA} \]
\item Quantité d'électricité :
\[ q = i\,\Delta t = \num{8,0}\times 10^{-2} \times \num{0,20} = \SI{1,6e-2}{C} \]
\textbf{Conclusion} : la f.é.m. induite vaut $\SI{0,40}{V}$ ; elle fait circuler un courant de $\SI{80}{mA}$ pendant la variation du champ.
\end{enumerate}
\end{exoresolu}

\courssub{Auto-induction et inductance}

\begin{definition}[Inductance propre d'une bobine]
Lorsqu'un courant $i$ traverse une bobine, elle crée un champ magnétique $\vec{B}$ proportionnel à $i$. Le flux magnétique $\Phi_p$ que ce champ crée \emph{à travers la bobine elle-même} (flux propre) est proportionnel à l'intensité du courant :
\[ \Phi_p = L\,i \]
Le coefficient de proportionnalité $L$ est l'\cle{inductance} (ou inductance propre) de la bobine. Elle ne dépend que de la géométrie de la bobine (nombre de spires $N$, rayon, longueur, présence d'un noyau ferromagnétique) et non du courant.
\end{definition}

\begin{propriete}[Unité de l'inductance : le henry (H)]
\[ 1\,\mathrm{H} = 1\,\dfrac{\mathrm{Wb}}{\mathrm{A}} = 1\,\dfrac{\mathrm{V}\cdot\mathrm{s}}{\mathrm{A}} = 1\,\dfrac{\mathrm{J}}{\mathrm{A}^2} \]
\textbf{Analyse dimensionnelle :} $[L] = [\Phi]/[i] = \mathrm{Wb}\cdot\mathrm{A}^{-1} = \mathrm{V}\cdot\mathrm{s}\cdot\mathrm{A}^{-1}$.
\end{propriete}

\begin{propriete}[F.é.m. d'auto-induction]
Si le courant $i$ traversant la bobine varie, le flux propre $\Phi_p = Li$ varie. D'après la loi de Faraday, une f.é.m. d'\cle{auto-induction} apparaît aux bornes de la bobine :
\[ e = -\dfrac{d\Phi_p}{dt} = -L\,\dfrac{di}{dt} \]
Pour une variation régulière sur $\Delta t$ :
\[ e = -L\,\dfrac{\Delta i}{\Delta t} \]
Le signe « $-$ » (loi de Lenz) indique que cette f.é.m. \textbf{s'oppose à la variation du courant} :
\begin{itemize}
\item À l'établissement ($di/dt > 0$) : $e < 0$, elle freine la montée du courant.
\item À la disparition ($di/dt < 0$) : $e > 0$, elle tend à maintenir le courant (phénomène d'\cle{étincelle de rupture} à l'ouverture d'un circuit inductif).
\end{itemize}
\end{propriete}

\begin{propriete}[Énergie magnétique emmagasinée dans une bobine]
Le travail fourni par le générateur pour vaincre la f.é.m. d'auto-induction lors de l'établissement du courant est stocké sous forme d'énergie magnétique dans le champ créé par la bobine :
\[ E_m = \dfrac{1}{2}\,L\,i^2 \]
Cette énergie est restituée lors de la disparition du courant (échauffement par effet Joule dans la résistance du circuit, ou étincelle).
\end{propriete}

\begin{methode}[Auto-induction : reconnaître et calculer]
\begin{enumerate}
\item \textbf{Reconnaître} : un seul circuit, comportant une bobine ; c'est la variation du courant \emph{dans le circuit lui-même} qui crée la f.é.m. induite (le circuit est à la fois inducteur et induit).
\item \textbf{Flux propre} : $\Phi_p = L\,i$, avec $L$ l'inductance en henry (H).
\item \textbf{F.é.m. d'auto-induction} :
\[ e = -L\,\dfrac{di}{dt} \qquad \text{(variation régulière : } e = -L\,\dfrac{\Delta i}{\Delta t}\text{)} \]
\item \textbf{Conséquences expérimentales} : retard à l'établissement du courant à la fermeture (une lampe en série avec la bobine s'allume avec retard), étincelle de rupture à l'ouverture.
\item Vérifier l'homogénéité : $[L] = \dfrac{\mathrm{V}\cdot\mathrm{s}}{\mathrm{A}} = \mathrm{H}$.
\end{enumerate}
\end{methode}

\begin{exoresolu}[F.é.m. d'auto-induction]
Une bobine d'inductance $L = \SI{0,20}{H}$ est parcourue par un courant dont l'intensité passe de $i_1 = \SI{2,0}{A}$ à $i_2 = \SI{0,50}{A}$ en $\Delta t = \SI{50}{ms}$, de façon régulière.
\begin{enumerate}
\item Calculer la valeur absolue de la f.é.m. d'auto-induction.
\item Quel est le rôle de cette f.é.m. vis-à-vis de la diminution du courant ?
\item Calculer l'énergie magnétique emmagasinée dans la bobine lorsque $i = \SI{2,0}{A}$.
\end{enumerate}
\tcblower
\textbf{Solution.}
\begin{enumerate}
\item Variation d'intensité : $\Delta i = i_2 - i_1 = \num{0,50} - \num{2,0} = \SI{-1,5}{A}$, avec $\Delta t = \SI{5,0e-2}{s}$.
\[ |e| = L\,\left|\dfrac{\Delta i}{\Delta t}\right| = \num{0,20} \times \dfrac{\num{1,5}}{\num{5,0}\times 10^{-2}} = \num{0,20} \times 30 = \SI{6,0}{V} \]
\item D'après la loi de Lenz, cette f.é.m. \textbf{s'oppose à la diminution du courant} : elle tend à prolonger le courant (c'est elle qui provoque l'étincelle de rupture à l'ouverture d'un circuit inductif).
\item Énergie magnétique emmagasinée :
\[ E_m = \dfrac{1}{2}L\,i^2 = \dfrac{1}{2}\times \num{0,20} \times (\num{2,0})^2 = \SI{0,40}{J} \]
\textbf{Conclusion} : la bobine développe une f.é.m. de $\SI{6,0}{V}$ qui freine la disparition du courant ; elle restitue l'énergie $\SI{0,40}{J}$ stockée sous forme magnétique.
\end{enumerate}
\end{exoresolu}

\courssub{Applications et bilan énergétique}

\begin{methode}[Distinguer circuit inducteur et circuit induit dans un problème]
\begin{enumerate}
\item Chercher la \textbf{source du champ magnétique variable} : aimant en mouvement, bobine alimentée par un générateur dont le courant varie $\Rightarrow$ c'est le \cle{circuit inducteur}.
\item Chercher le circuit \textbf{dépourvu de générateur} où apparaît un courant (détecté par un galvanomètre, une lampe) $\Rightarrow$ c'est le \cle{circuit induit}.
\item Vérifier que le circuit induit est \textbf{fermé} (sinon f.é.m. induite mais pas de courant).
\item Rédiger la chaîne causale : variation de $i$ dans l'inducteur $\rightarrow$ variation de $\vec{B}$ $\rightarrow$ variation de $\Phi$ dans l'induit $\rightarrow$ f.é.m. induite $\rightarrow$ courant induit.
\end{enumerate}
\end{methode}

\begin{exoresolu}[Deux bobines en regard — Principe du transformateur]
Deux bobines $B_1$ et $B_2$ sont placées face à face. $B_1$ est reliée à un générateur, un rhéostat et un interrupteur ; $B_2$ est reliée à un galvanomètre. On ferme l'interrupteur, puis on le laisse fermé, puis on l'ouvre.
\begin{enumerate}
\item Identifier le circuit inducteur et le circuit induit.
\item Décrire les indications du galvanomètre aux trois phases de l'expérience.
\item Citer une application technologique majeure de ce dispositif au Cameroun.
\end{enumerate}
\tcblower
\textbf{Solution.}
\begin{enumerate}
\item Le \textbf{circuit inducteur} est $B_1$ (alimenté par le générateur, source du champ magnétique variable). Le \textbf{circuit induit} est $B_2$ (sans générateur, fermé sur le galvanomètre).
\item \textbf{Fermeture} : le courant dans $B_1$ passe de 0 à une valeur maximale, le flux dans $B_2$ varie : le galvanomètre dévie \emph{brièvement} dans un sens. \textbf{Interrupteur maintenu fermé} : le courant est constant, le flux ne varie plus : le galvanomètre indique \emph{zéro}. \textbf{Ouverture} : le flux diminue brutalement : le galvanomètre dévie \emph{brièvement dans le sens opposé}.
\item C'est le principe du \textbf{transformateur} statique. Au Cameroun, \textbf{ENEO} (ex-SONEL/SONATREL) utilise des transformateurs abaisseurs dans les postes de distribution (ex: poste de Nachtigal, Lom Pangar, Edéa) pour faire passer la tension des lignes de transport (\SI{90}{kV}, \SI{225}{kV}) à la tension de distribution (\SI{15}{kV}, \SI{30}{kV}) puis au \SI{220}{V} domestique. C'est aussi le principe des chargeurs à induction des téléphones (MTN/Orange) et des plaques à induction.
\end{enumerate}
\textbf{Conclusion} : seules les variations du courant inducteur (fermeture, ouverture, réglage du rhéostat) produisent un courant induit. Le courant continu constant n'induit rien.
\end{exoresolu}

\begin{savaistu}[De Faraday à l'électrification du Cameroun]
\begin{itemize}
\item \textbf{Histoire :} Michael Faraday (1831) et Joseph Henry (1832) découvrent indépendamment l'induction. Faraday construit le premier générateur (disque de Faraday) ; Henry découvre l'auto-induction. L'unité d'inductance, le \textbf{henry (H)}, honore Henry.
\item \textbf{Alternateurs :} Les barrages hydroélectriques de \textbf{Edéa}, \textbf{Lom Pangar} et \textbf{Nachtigal} utilisent des alternateurs : un rotor (aimants ou électroaimants) tourne dans un stator (bobines). La variation périodique du flux induit une tension alternative triphasée (fréquence \SI{50}{Hz} au Cameroun).
\item \textbf{Transport et distribution :} L'induction mutuelle dans les transformateurs permet d'élever la tension pour le transport (faibles pertes par effet Joule $P = R I^2$) et de l'abaisser pour la sécurité des usagers.
\item \textbf{Cuisson à induction :} Une bobine sous la plaque de verre crée un champ magnétique alternatif à haute fréquence (\SIrange{20}{100}{kHz}). Ce champ induit des \textbf{courants de Foucault} dans le fond ferromagnétique de la casserole (circuit induit), l'échauffant directement par effet Joule sans chauffer la plaque. Rendement > 90\%.
\item \textbf{Freins à courants de Foucault :} Utilisés sur les trains (bientôt le train Ngaoundéré-Yaoundé-Douala) et les poids lourds : un aimant puissant induit des courants dans un disque métallique tournant solidaire de la roue ; la force de Laplace freine le disque sans contact ni usure.
\end{itemize}
\end{savaistu}

\begin{aretenir}[Bilan des formules clés — Induction électromagnétique]
\begin{center}
\begin{tabular}{|l|c|c|}\hline
\textbf{Notion} & \textbf{Formule} & \textbf{Unité SI} \\ \hline
Flux magnétique (1 spire) & $\Phi = B S \cos\alpha$ & Weber (\textbf{Wb}) = $\mathrm{T}\cdot\mathrm{m}^2$ \\
Flux magnétique ($N$ spires) & $\Phi = N B S \cos\alpha$ & \textbf{Wb} \\
Inductance propre & $L = \dfrac{\Phi_p}{i}$ & Henry (\textbf{H}) = $\mathrm{V}\cdot\mathrm{s}\cdot\mathrm{A}^{-1}$ \\
F.é.m. d'induction (Faraday) & $e = -\dfrac{d\Phi}{dt} = -\dfrac{\Delta\Phi}{\Delta t}$ & Volt (\textbf{V}) \\
F.é.m. d'auto-induction & $e = -L\,\dfrac{di}{dt} = -L\,\dfrac{\Delta i}{\Delta t}$ & Volt (\textbf{V}) \\
Courant induit (circuit $R$) & $i = \dfrac{|e|}{R}$ & Ampère (\textbf{A}) \\
Énergie magnétique & $E_m = \dfrac{1}{2} L i^2$ & Joule (\textbf{J}) \\ \hline
\end{tabular}
\end{center}
\textbf{Rappel} : $\alpha = (\vec{n}, \vec{B})$ ; $\vec{n}$ orienté par la règle du tire-bouchon (sens positif du courant). Le signe « $-$ » = Loi de Lenz (opposition à la variation).
\end{aretenir}

\begin{attention}[Les 5 erreurs qui coûtent des points au Probatoire]
\begin{enumerate}
\item \textbf{Oublier le nombre de spires $N$} dans le flux : $\Phi = NBS\cos\alpha$ pour une bobine. Un flux calculé pour une seule spire fausse ensuite la f.é.m. $e$.
\item \textbf{Confondre l'angle} : $\alpha$ est l'angle entre $\vec{B}$ et la \emph{normale} $\vec{n}$ à la surface, pas l'angle entre $\vec{B}$ et le plan de la spire. Si l'énoncé donne l'angle avec le plan ($\beta$), prendre le complémentaire : $\cos\alpha = \sin\beta$.
\item \textbf{Croire qu'un champ constant induit un courant} : c'est la \emph{variation} du flux qui crée la f.é.m. induite. Aimant immobile ou courant continu constant $\Rightarrow$ $e = 0$.
\item \textbf{Mal gérer le signe et les durées} : le signe « $-$ » de $e = -\dfrac{d\Phi}{dt}$ traduit Lenz ; si l'on demande une \emph{valeur}, donner $|e|$ en volts. Convertir systématiquement les durées en secondes ($\SI{50}{ms} = \SI{5,0e-2}{s}$) et les surfaces en $\mathrm{m}^2$ ($1\,\mathrm{cm}^2 = 10^{-4}\,\mathrm{m}^2$).
\item \textbf{Confondre induction et auto-induction} : dans l'auto-induction, il n'y a qu'\emph{un seul} circuit et $e = -L\,\dfrac{di}{dt}$ ; ne pas écrire $e = -\dfrac{\Delta\Phi}{\Delta t}$ avec le flux d'un autre circuit. Préciser toujours qui est l'inducteur et qui est l'induit.
\end{enumerate}
\end{attention}

\begin{exercice}[Entraînement — Induction et auto-induction]
\begin{enumerate}
\item \textbf{Bobine tournante :} Une bobine de $N = 150$ spires, de surface $S = \SI{40}{cm^2}$, tourne à vitesse constante dans un champ magnétique uniforme $B = \SI{0,80}{T}$. Son axe de rotation est perpendiculaire à $\vec{B}$. À l'instant $t=0$, le plan de la bobine est perpendiculaire à $\vec{B}$ (flux maximal). La bobine effectue un quart de tour en $\Delta t = \SI{0,10}{s}$.
\begin{enumerate}
\item Calculer le flux initial $\Phi_i$ et le flux final $\Phi_f$.
\item Déterminer la valeur moyenne de la f.é.m. induite pendant ce quart de tour.
\item Si la résistance de la bobine est $R = \SI{12}{\Omega}$, calculer l'intensité moyenne du courant induit.
\end{enumerate}
\item \textbf{Auto-induction :} Un circuit série comprend une bobine d'inductance $L = \SI{50}{mH}$ et une résistance $R = \SI{10}{\Omega}$, alimenté par une tension continue $E = \SI{12}{V}$. À l'instant $t=0$, on ferme l'interrupteur.
\begin{enumerate}
\item Écrire l'équation différentielle du circuit (loi des mailles).
\item Donner l'expression de l'intensité $i(t)$ (sans redémontrer : $i(t) = \dfrac{E}{R}(1 - e^{-t/\tau})$ avec $\tau = L/R$).
\item Calculer la constante de temps $\tau$ et la f.é.m. d'auto-induction à l'instant $t=0^+$ (juste après la fermeture).
\item Calculer l'énergie magnétique stockée dans la bobine en régime permanent.
\end{enumerate}
\item \textbf{Chute d'aimant dans un tube conducteur :} On laisse tomber un aimant droit dans un tube vertical en cuivre (bon conducteur, non ferromagnétique). L'aimant tombe plus lentement que dans le vide.
\begin{enumerate}
\item Expliquer ce phénomène en utilisant la loi de Lenz et les courants de Foucault.
\item Que se passerait-il si le tube était en plastique ? En cuivre mais fendu longitudinalement ?
\end{enumerate}
\end{enumerate}
\end{exercice}

\begin{corrige}[Exercice n°1]
\begin{enumerate}
\item \textbf{Bobine tournante :}
\begin{enumerate}
\item $S = \SI{40}{cm^2} = \num{4,0e-3}{m^2}$. $\Phi_i = NBS\cos 0 = 150 \times 0,80 \times \num{4,0e-3} = \SI{0,48}{Wb}$. $\Phi_f = NBS\cos 90^\circ = 0$.
\item $\Delta\Phi = \Phi_f - \Phi_i = -\SI{0,48}{Wb}$. $|e| = \dfrac{|\Delta\Phi|}{\Delta t} = \dfrac{0,48}{0,10} = \SI{4,8}{V}$.
\item $i = \dfrac{|e|}{R} = \dfrac{4,8}{12} = \SI{0,40}{A}$.
\end{enumerate}
\item \textbf{Auto-induction :}
\begin{enumerate}
\item $E = u_R + u_L = Ri + L\dfrac{di}{dt}$.
\item $\tau = \dfrac{L}{R} = \dfrac{\num{50e-3}}{10} = \SI{5,0}{ms}$. $i(t) = \dfrac{12}{10}(1 - e^{-t/5\text{ms}}) = 1,2(1 - e^{-t/5\text{ms}})\,\mathrm{A}$.
\item À $t=0^+$, $i=0$, donc $u_R=0$. Loi des mailles : $E = u_L(0^+) \Rightarrow e(0^+) = \SI{12}{V}$ (la f.é.m. d'auto-induction s'oppose totalement à la tension du générateur au tout début).
\item Régime permanent : $i_{\text{max}} = E/R = \SI{1,2}{A}$. $E_m = \dfrac{1}{2} L i_{\text{max}}^2 = 0,5 \times \num{50e-3} \times 1,44 = \SI{36}{mJ}$.
\end{enumerate}
\item \textbf{Chute d'aimant :}
\begin{enumerate}
\item L'aimant en chute crée une variation de flux dans chaque section du tube cuivre. Des \textbf{courants de Foucault} (courants induits en boucles fermées dans l'épaisseur du cuivre) apparaissent. D'après Lenz, ces courants créent un champ qui s'oppose à la cause : la chute de l'aimant. Une force de Laplace verticale vers le haut freine l'aimant. Il atteint une vitesse limite (vitesse terminale) où force magnétique = poids.
\item Tube plastique : isolant, pas de courants de Foucault $\Rightarrow$ chute libre (sans frottements magnétiques). Tube cuivre fendu : la fente empêche les boucles de courant fermées (circuit ouvert) $\Rightarrow$ pas de courants de Foucault $\Rightarrow$ chute libre.
\end{enumerate}
\end{enumerate}
\end{corrige}

\newpage

\coursec{Exercices}

\courssub{Je vérifie mes connaissances}

\begin{exercice}[QCM : Unités et définitions]
Parmi les affirmations suivantes, une seule est correcte. Laquelle ?
\begin{enumerate}
    \item L'unité du flux magnétique dans le SI est le tesla (T).
    \item L'unité de l'inductance dans le SI est le weber (Wb).
    \item Le flux magnétique à travers une surface plane de vecteur aire $\vec{S}$ plongée dans un champ magnétique uniforme $\vec{B}$ vaut $\Phi = \vec{B} \cdot \vec{S}$.
    \item La f.é.m. d'auto-induction dans une bobine d'inductance $L$ parcourue par un courant $i$ vaut $e = L\,i$.
\end{enumerate}
\end{exercice}

\begin{exercice}[Vrai ou Faux ? Justifiez]
Déterminez si les affirmations suivantes sont vraies ou fausses. Justifiez brièvement votre réponse par une loi ou un contre-exemple.
\begin{enumerate}
    \item Un aimant immobile à l'intérieur d'une bobine fermée crée un courant induit dans la bobine.
    \item La f.é.m. induite dans un circuit dépend uniquement de la variation du flux magnétique, pas de la rapidité de cette variation.
    \item Selon la loi de Lenz, le courant induit tend à s'opposer à la cause qui le produit.
    \item L'inductance $L$ d'une bobine dépend du courant $i$ qui la traverse.
\end{enumerate}
\end{exercice}

\begin{exercice}[Calcul direct de flux]
Une bobine circulaire de $N = 250$ spires a un rayon $r = 4,0\,\text{cm}$. Elle est placée dans un champ magnétique uniforme de norme $B = 0,80\,\text{T}$. Le vecteur normal à la bobine forme un angle de $30^\circ$ avec les lignes de champ.
\begin{enumerate}
    \item Calculez le flux magnétique $\Phi$ traversant une spire (en Wb).
    \item Calculez le flux total $\Phi_{\text{tot}}$ concatené par la bobine (en Wb).
\end{enumerate}
\end{exercice}

\begin{exercice}[F.é.m. induite : application numérique]
Le flux magnétique traversant une bobine de $N = 150$ spires varie linéairement de $\Phi_1 = 2,0 \times 10^{-3}\,\text{Wb}$ à $\Phi_2 = 5,0 \times 10^{-3}\,\text{Wb}$ en un temps $\Delta t = 0,10\,\text{s}$.
\begin{enumerate}
    \item Calculez la valeur moyenne de la f.é.m. induite $e$ (en V) pendant cet intervalle.
    \item Précisez le sens du courant induit si le flux est orienté dans le sens de la normale choisie et qu'il augmente.
\end{enumerate}
\end{exercice}

\courssub{Je m'entraîne}

\begin{exercice}[Barreau glissant sur rails : générateur linéaire]
Un barreau conducteur $MN$ de longueur $\ell = 12\,\text{cm}$ et de résistance $R = 2,0\,\Omega$ glisse sans frottement à la vitesse constante $v = 3,0\,\text{m/s}$ sur deux rails parallèles horizontaux reliés à leurs extrémités par un conducteur ohmique de résistance négligeable, formant un circuit rectangulaire fermé. L'ensemble est plongé dans un champ magnétique uniforme vertical $\vec{B}$ de norme $B = 0,50\,\text{T}$, dirigé vers le bas.
\begin{enumerate}
    \item En utilisant la force de Laplace sur les porteurs de charge, montrez que la f.é.m. induite aux bornes du barreau vaut $e = B\,\ell\,v$.
    \item Calculez la valeur de cette f.é.m. (en V).
    \item Déterminez l'intensité du courant induit (en A) et son sens (de $M$ vers $N$ ou de $N$ vers $M$ ?) en appliquant la loi de Lenz.
    \item Quelle force extérieure $\vec{F}_{\text{ext}}$ faut-il appliquer au barreau pour maintenir sa vitesse constante ? Donnez sa valeur et son sens.
\end{enumerate}
\end{exercice}

\begin{exercice}[Variation d'orientation : alternateur élémentaire]
Une bobine plane de $N = 200$ spires, de surface $S = 50\,\text{cm}^2$, tourne à vitesse angulaire constante $\omega = 100\pi\,\text{rad/s}$ (soit 50 tours/s) dans un champ magnétique uniforme horizontal $\vec{B}$ de norme $B = 5,0 \times 10^{-2}\,\text{T}$. L'axe de rotation est vertical. À l'instant $t=0$, le vecteur normal $\vec{n}$ à la bobine est aligné avec $\vec{B}$.
\begin{enumerate}
    \item Exprimez le flux total $\Phi_{\text{tot}}(t)$ en fonction du temps.
    \item En déduisez l'expression de la f.é.m. induite $e(t)$. Quelle est sa valeur maximale $E_{\text{max}}$ (en V) ?
    \item La bobine alimente une résistance $R = 10\,\Omega$. Exprimez l'intensité du courant $i(t)$ et sa valeur maximale.
\end{enumerate}
\end{exercice}

\begin{exercice}[Auto-induction : établissement du courant]
Un circuit série comprend une bobine d'inductance $L = 0,20\,\text{H}$ et de résistance $r = 5,0\,\Omega$, un interrupteur $K$ et un générateur de tension continue $E = 12\,\text{V}$. L'interrupteur est fermé à l'instant $t=0$. Le courant est nul pour $t<0$.
\begin{enumerate}
    \item Écrivez l'équation différentielle régissant l'intensité $i(t)$ pour $t \ge 0$ (loi des mailles).
    \item Donnez l'expression de $i(t)$ (sans redémontrer) et calculez la constante de temps $\tau$ (en ms).
    \item Calculez la f.é.m. d'auto-induction $e_L(t)$ à l'instant $t = 0^+$ et à l'instant $t = \tau$.
    \item Vérifiez que la somme des tensions aux bornes de la résistance et de la bobine vaut bien $E$ à tout instant.
\end{enumerate}
\end{exercice}

\begin{exercice}[Rupture de circuit : l'étincelle]
Une bobine d'inductance $L = 0,15\,\text{H}$ et de résistance interne $r = 8,0\,\Omega$ est parcourue par un courant permanent $I_0 = 1,2\,\text{A}$ fourni par une pile. On ouvre brutalement l'interrupteur en série. On admet que le courant s'annule selon une exponentielle décroissante de constante de temps $\tau' = 0,50\,\text{ms}$ (due à la capacité parasite du circuit).
\begin{enumerate}
    \item Calculez l'énergie magnétique stockée dans la bobine juste avant l'ouverture (en mJ).
    \item Déterminez la valeur maximale de la f.é.m. d'auto-induction $|e_L|_{\text{max}}$ (en V) apparaissant aux bornes de la bobine lors de la rupture.
    \item Cette surtension peut-elle provoquer un arc électrique (rupture diélectrique de l'air) si l'écartement des contacts de l'interrupteur est de $1,0\,\text{mm}$ ? (Champ de rupture de l'air $\approx 3 \times 10^6\,\text{V/m}$).
    \item Quel composant électronique place-t-on habituellement en parallèle sur une bobine pour protéger l'interrupteur ? Expliquez son rôle.
\end{enumerate}
\end{exercice}

\courssub{Je me prépare au Probatoire}

\begin{exercice}[Probatoire 2022 type : Bobine tournante et éclairage LED]
Un élève souhaite allumer une diode électroluminescente (DEL) rouge en faisant tourner manuellement une bobine dans le champ magnétique terrestre. La bobine a $N = 300$ spires, une section $S = 20\,\text{cm}^2$ et une résistance interne $r = 15\,\Omega$. Le champ magnétique terrestre local a une intensité $B_{\text{T}} = 4,5 \times 10^{-5}\,\text{T}$ ; on suppose qu'il est horizontal. La DEL rouge se met à conduire pour une tension seuil $U_{\text{seuil}} = 1,8\,\text{V}$ et un courant minimal $I_{\text{min}} = 5,0\,\text{mA}$.
\begin{enumerate}
    \item L'élève fait tourner la bobine à une fréquence de rotation $f = 2,0\,\text{rev/s}$ (tours par seconde), l'axe de rotation étant vertical. Établissez l'expression de la f.é.m. induite $e(t)$ en fonction du temps. Calculez sa valeur efficace $E_{\text{eff}}$ et sa valeur maximale $E_{\text{max}}$.
    \item La tension aux bornes de la DEL vaut-elle $e(t)$ ? Justifiez. La DEL s'allume-t-elle ? (On néglige la self-induction de la bobine).
    \item L'élève décide d'augmenter le nombre de spires à $N' = 1200$ en gardant le même fil (donc même section de fil, même résistivité). La résistance devient $r' = 60\,\Omega$. Recalculez $E'_{\text{max}}$. La DEL s'allume-t-elle maintenant ?
    \item Pour améliorer le rendement, l'élève insère un noyau de fer doux dans la bobine, ce qui multiplie le champ magnétique à l'intérieur par un facteur $\mu_r = 200$ (perméabilité relative). Recalculez $E''_{\text{max}}$ pour $N' = 1200$ spires. Concluez sur la faisabilité de l'éclairage.
\end{enumerate}
\end{exercice}

\begin{exercice}[Probatoire 2021 type : Freinage magnétique et chute d'aimant]
On laisse tomber un aimant cylindrique de masse $m = 50\,\text{g}$ et de moment dipolaire $\vec{\mu}$ (orienté verticalement, pôle Nord en bas) dans un tube vertical en cuivre de longueur $L = 1,0\,\text{m}$, de rayon intérieur $R = 1,2\,\text{cm}$ et d'épaisseur $e = 1,0\,\text{mm}$. La conductivité du cuivre est $\sigma = 5,9 \times 10^7\,\text{S/m}$. L'aimant atteint une vitesse limite $v_{\text{lim}} = 0,15\,\text{m/s}$.
\begin{enumerate}
    \item Expliquez qualitativement, en utilisant la loi de Lenz, pourquoi la vitesse de l'aimant se stabilise à une valeur limite. Dessinez les courants de Foucault induits dans une section transversale du tube au passage de l'aimant.
    \item On modélise l'interaction par une force de freinage magnétique $\vec{F}_{\text{mag}} = -k\,\vec{v}$ proportionnelle à la vitesse (visqueuse). Exprimez la constante $k$ en fonction de $m$, $g$ et $v_{\text{lim}}$. Calculez $k$ (en N$\cdot$s/m).
    \item L'énergie potentielle de pesanteur perdue par l'aimant pendant sa chute à vitesse limite sur la longueur $L$ est entièrement dissipée par effet Joule dans le tube. Calculez la puissance moyenne dissipée $P_{\text{Joule}}$ (en W).
    \item Si on fend le tube longitudinalement sur toute sa hauteur (fente de $1\,\text{mm}$ de large), que devient la vitesse limite ? Justifiez.
\end{enumerate}
\end{exercice}

\begin{exercice}[Probatoire 2023 type : Transformateur réel et auto-induction]
Un transformateur monophasé alimente un poste de soudage. Le primaire a $N_1 = 500$ spires, une inductance propre $L_1 = 2,0\,\text{H}$ et une résistance $r_1 = 2,0\,\Omega$. Le secondaire a $N_2 = 50$ spires, une inductance propre $L_2 = 0,02\,\text{H}$ et une résistance $r_2 = 0,05\,\Omega$. Le coefficient de couplage est $k = 0,98$. Le primaire est branché sur le secteur $u_1(t) = U_1\sqrt{2}\cos(\omega t)$ avec $U_1 = 230\,\text{V}$ (efficace) et $f = 50\,\text{Hz}$. Le secondaire est en circuit ouvert ($i_2 = 0$).
\begin{enumerate}
    \item Calculez la réactance d'induction du primaire $X_1 = L_1\omega$ (en $\Omega$). Comparez $X_1$ à $r_1$. Le primaire est-il un bon inducteur ?
    \item Déterminez l'intensité efficace $I_{1,0}$ absorbée par le primaire à vide (courant de magnétisation). Donnez son expression instantanée $i_{1,0}(t)$ en choisissant l'origine des temps pour que $u_1(0) = U_1\sqrt{2}$.
    \item Calculez la tension efficace $U_{2,0}$ aux bornes du secondaire à vide. Vérifiez le rapport de transformation.
    \item On met maintenant le secondaire en court-circuit ($u_2 = 0$) par une tige de soudage de résistance négligeable. On admet que les flux de fuite sont modélisés par des inductances de fuite $L_{\sigma1} = (1-k)L_1$ et $L_{\sigma2} = (1-k)L_2$ ramenées au primaire. Calculez l'inductance de fuite totale ramenée au primaire $L_{\sigma}$ et la tension de fuite $U_{\sigma}$ pour un courant de court-circuit efficace $I_{1,\text{cc}} = 100\,\text{A}$. Commenter l'intérêt de la fuite pour le soudage à l'arc.
\end{enumerate}
\end{exercice}

\newpage

\coursec{Corrigés des exercices}

\begin{corrige}[Exercice 1 — QCM : Unités et définitions]
\textbf{Réponse correcte : 3.}

Analysons chaque affirmation.
\begin{enumerate}
    \item \textbf{Fausse.} Le tesla (\SI{}{\tesla}) est l'unité du \cle{champ magnétique} $\vec{B}$. L'unité du flux magnétique est le \cle{weber} (\SI{}{\weber}), avec $1\ \mathrm{Wb} = 1\ \mathrm{T}\cdot\mathrm{m}^2$.
    \item \textbf{Fausse.} C'est l'inverse : l'unité de l'inductance $L$ est le \cle{henry} (\SI{}{\henry}). Le weber est l'unité du flux.
    \item \textbf{Vraie.} Pour une surface plane de vecteur aire $\vec{S} = S\,\vec{n}$ plongée dans un champ uniforme $\vec{B}$, le flux est le produit scalaire $\Phi = \vec{B}\cdot\vec{S} = B\,S\cos\theta$, où $\theta = (\vec{B},\vec{n})$.
    \item \textbf{Fausse.} La f.é.m. d'auto-induction est $e = -L\,\dfrac{di}{dt}$ : elle dépend de la \emph{dérivée} du courant, non du courant lui-même. En régime permanent ($i$ constant), $e = 0$.
\end{enumerate}
\end{corrige}

\begin{corrige}[Exercice 2 — Vrai ou Faux ? Justifiez]
\begin{enumerate}
    \item \textbf{Faux.} Un courant induit n'apparaît que si le flux magnétique \emph{varie}. Aimant immobile $\Rightarrow$ flux constant $\Rightarrow$ $e = -\dfrac{d\Phi}{dt} = 0$ : aucun courant induit. C'est le \emph{mouvement relatif} aimant/bobine qui crée l'induction.
    \item \textbf{Faux.} La loi de Faraday $e = -\dfrac{d\Phi}{dt}$ montre que la f.é.m. est proportionnelle à la \emph{vitesse de variation} du flux : plus la variation est rapide (petit $\Delta t$), plus $|e|$ est grande. À variation $\Delta\Phi$ identique, diviser $\Delta t$ par 10 multiplie $|e|$ par 10.
    \item \textbf{Vrai.} C'est l'énoncé même de la \cle{loi de Lenz} : par ses effets (champ induit, force de Laplace), le courant induit s'oppose à la cause qui lui a donné naissance (approche de l'aimant, augmentation du courant, etc.). C'est une conséquence de la conservation de l'énergie.
    \item \textbf{Faux.} L'inductance $L$ est une caractéristique \emph{géométrique} de la bobine (nombre de spires, dimensions, nature du noyau) : $L = \mu_0\,\dfrac{N^2 S}{\ell}$ pour un solénoïde long. Elle ne dépend pas du courant $i$ (tant que le matériau n'est pas saturé).
\end{enumerate}
\end{corrige}

\begin{corrige}[Exercice 3 — Calcul direct de flux]
\textbf{Données :} $N = 250$, $r = \SI{4,0}{\centi\meter} = \SI{4,0e-2}{\meter}$, $B = \SI{0,80}{\tesla}$, $\theta = \SI{30}{\degree}$.

\textbf{1. Flux à travers une spire.}

Surface d'une spire :
\[
S = \pi r^2 = \pi \times (\SI{4,0e-2}{\meter})^2 = \pi \times \SI{1,6e-3}{\square\meter} \approx \SI{5,03e-3}{\square\meter}.
\]
Flux à travers une spire :
\[
\Phi = B\,S\cos\theta = \SI{0,80}{\tesla} \times \SI{5,03e-3}{\square\meter} \times \cos\SI{30}{\degree}
= \SI{0,80}{\tesla} \times \SI{5,03e-3}{\square\meter} \times 0,866.
\]
\[
\boxed{\Phi \approx \SI{3,5e-3}{\weber} = \SI{3,5}{\milli\weber}}
\]

\textbf{2. Flux total concaténé.}

Le même flux traverse les $N$ spires :
\[
\Phi_{\text{tot}} = N\,\Phi = 250 \times \SI{3,48e-3}{\weber}.
\]
\[
\boxed{\Phi_{\text{tot}} \approx \SI{0,87}{\weber}}
\]

\textbf{Vérification :} si la bobine était perpendiculaire au champ ($\theta = \SI{90}{\degree}$), le flux serait nul ; ici $\cos\SI{30}{\degree} < 1$, le flux est bien inférieur au flux maximal $B S = \SI{4,0}{\milli\weber}$.
\end{corrige}

\begin{corrige}[Exercice 4 — F.é.m. induite : application numérique]
\textbf{Données :} $N = 150$, $\Phi_1 = \SI{2,0e-3}{\weber}$, $\Phi_2 = \SI{5,0e-3}{\weber}$, $\Delta t = \SI{0,10}{\second}$.

\textbf{1. F.é.m. moyenne.}

Variation de flux par spire : $\Delta\Phi = \Phi_2 - \Phi_1 = \SI{3,0e-3}{\weber}$.

Loi de Faraday (valeur moyenne, en tenant compte des $N$ spires) :
\[
|e| = N\,\left|\dfrac{\Delta\Phi}{\Delta t}\right| = 150 \times \dfrac{\SI{3,0e-3}{\weber}}{\SI{0,10}{\second}} = 150 \times \SI{3,0e-2}{\volt}.
\]
\[
\boxed{|e| = \SI{4,5}{\volt}}
\]
Avec la convention d'orientation choisie, $e = -\SI{4,5}{\volt}$ (le signe moins traduit la loi de Lenz).

\textbf{2. Sens du courant induit.}

Le flux, orienté dans le sens de la normale $\vec{n}$, \emph{augmente}. D'après la loi de Lenz, le courant induit crée un champ induit $\vec{B}_{\text{induit}}$ \emph{opposé} à $\vec{B}$, donc dirigé selon $-\vec{n}$. Par la règle du tire-bouchon (ou de la main droite), le courant induit circule dans le \textbf{sens inverse du sens d'orientation choisi} sur la spire (sens négatif).
\end{corrige}

\begin{corrige}[Exercice 5 — Barreau glissant sur rails : générateur linéaire]
\textbf{Données :} $\ell = \SI{12}{\centi\meter} = \SI{0,12}{\meter}$, $R = \SI{2,0}{\ohm}$, $v = \SI{3,0}{\meter\per\second}$, $B = \SI{0,50}{\tesla}$ (vertical, vers le bas).

\textbf{1. Démonstration de $e = B\ell v$.}

Les porteurs de charge (électrons, charge $q = -e_0$) du barreau sont entraînés à la vitesse $\vec{v}$. Chacun subit la force magnétique de Lorentz :
\[
\vec{F}_m = q\,\vec{v}\wedge\vec{B}.
\]
Comme $\vec{v}\perp\vec{B}$, cette force est dirigée \emph{le long du barreau} et vaut en norme $F_m = |q|\,v\,B$. Elle sépare les charges : les électrons s'accumulent à une extrémité, l'autre se charge positivement. Il apparaît un champ électrique $\vec{E}$ dans le barreau. À l'équilibre (circuit ouvert) :
\[
q\vec{E} + q\,\vec{v}\wedge\vec{B} = \vec{0} \quad\Rightarrow\quad E = v\,B.
\]
La f.é.m. est la tension entre les extrémités du barreau :
\[
e = E\,\ell = B\,\ell\,v. \qquad \blacksquare
\]

\textbf{2. Valeur numérique.}
\[
e = \SI{0,50}{\tesla} \times \SI{0,12}{\meter} \times \SI{3,0}{\meter\per\second} \quad\Rightarrow\quad \boxed{e = \SI{0,18}{\volt}}
\]

\textbf{3. Courant induit.}

Le circuit est fermé sur une résistance totale $R = \SI{2,0}{\ohm}$ :
\[
i = \dfrac{e}{R} = \dfrac{\SI{0,18}{\volt}}{\SI{2,0}{\ohm}} \quad\Rightarrow\quad \boxed{i = \SI{0,090}{\ampere} = \SI{90}{\milli\ampere}}
\]

\emph{Sens :} supposons le barreau se déplaçant vers la droite. La surface du circuit augmente, donc le flux du champ (vers le bas) augmente en valeur absolue. D'après la loi de Lenz, le courant induit crée un champ dirigé \emph{vers le haut} : vu de dessus, il circule dans le sens trigonométrique. Les porteurs positifs sont poussés par $\vec{v}\wedge\vec{B}$ vers une extrémité : si $\vec{v}$ pointe vers la droite et $\vec{B}$ vers le bas, $\vec{v}\wedge\vec{B}$ dirige la force sur les charges positives de $M$ vers $N$ (avec $M$ côté gauche). Le courant induit circule donc \textbf{de $M$ vers $N$ dans le barreau} (le barreau joue le rôle de générateur).

\textbf{4. Force extérieure.}

Le barreau parcouru par $i$ dans le champ $\vec{B}$ subit la force de Laplace :
\[
\vec{F}_{\text{Lap}} = i\,\vec{\ell}\wedge\vec{B}, \qquad F_{\text{Lap}} = B\,i\,\ell = \SI{0,50}{\tesla} \times \SI{0,090}{\ampere} \times \SI{0,12}{\meter} = \SI{5,4e-3}{\newton}.
\]
Par la loi de Lenz, cette force s'oppose au mouvement (dirigée vers la gauche). La vitesse étant constante, la somme des forces est nulle (principe d'inertie) :
\[
\vec{F}_{\text{ext}} = -\vec{F}_{\text{Lap}} \quad\Rightarrow\quad \boxed{F_{\text{ext}} = \SI{5,4e-3}{\newton}, \text{ dans le sens du mouvement}}
\]

\textbf{Vérification énergétique :} puissance mécanique fournie $P = F_{\text{ext}}\,v = \SI{5,4e-3}{\newton}\times \SI{3,0}{\meter\per\second} = \SI{1,62e-2}{\watt}$ ; puissance Joule $P_J = R i^2 = \SI{2,0}{\ohm}\times(\SI{0,090}{\ampere})^2 = \SI{1,62e-2}{\watt}$. Égalité parfaite : l'énergie se conserve.
\end{corrige}

\begin{corrige}[Exercice 6 — Variation d'orientation : alternateur élémentaire]
\textbf{Données :} $N = 200$, $S = \SI{50}{\centi\meter\squared} = \SI{5,0e-3}{\square\meter}$, $\omega = \SI{100\pi}{\radian\per\second}$, $B = \SI{5,0e-2}{\tesla}$, $\theta(0) = 0$.

\textbf{1. Flux total.}

L'angle entre $\vec{n}$ et $\vec{B}$ est $\theta(t) = \omega t$. Le flux à travers une spire est $\Phi_1 = B S\cos(\omega t)$, donc :
\[
\boxed{\Phi_{\text{tot}}(t) = N B S \cos(\omega t) = 200\times \SI{5,0e-2}{\tesla}\times \SI{5,0e-3}{\square\meter}\cos(\SI{100\pi}{\radian\per\second} t) = \SI{5,0e-2}{\weber}\cos(\SI{100\pi}{\radian\per\second} t)}
\]

\textbf{2. F.é.m. induite.}
\[
e(t) = -\dfrac{d\Phi_{\text{tot}}}{dt} = N B S\,\omega\,\sin(\omega t).
\]
\[
\boxed{e(t) = \SI{5,0e-2}{\weber}\times \SI{100\pi}{\radian\per\second}\,\sin(\SI{100\pi}{\radian\per\second} t) = 5\pi\,\sin(\SI{100\pi}{\radian\per\second} t)\ \mathrm{(V)} \approx \SI{15,7}{\volt}\sin(\SI{100\pi}{\radian\per\second} t)}
\]
Valeur maximale :
\[
\boxed{E_{\text{max}} = N B S\,\omega = 5\pi \approx \SI{15,7}{\volt}}
\]

\textbf{3. Courant dans la résistance.}
\[
i(t) = \dfrac{e(t)}{R} = \dfrac{N B S\,\omega}{R}\sin(\omega t) = \dfrac{5\pi}{10}\sin(\SI{100\pi}{\radian\per\second} t) \approx \SI{1,57}{\ampere}\sin(\SI{100\pi}{\radian\per\second} t).
\]
\[
\boxed{I_{\text{max}} = \dfrac{E_{\text{max}}}{R} \approx \SI{1,57}{\ampere}}
\]
C'est le principe de l'\cle{alternateur} : la rotation de la bobine dans le champ produit une tension sinusoïdale de fréquence $f = \dfrac{\omega}{2\pi} = \SI{50}{\hertz}$ — exactement la fréquence du réseau ENEO !
\end{corrige}

\begin{corrige}[Exercice 7 — Auto-induction : établissement du courant]
\textbf{Données :} $L = \SI{0,20}{\henry}$, $r = \SI{5,0}{\ohm}$, $E = \SI{12}{\volt}$, $i(0^-) = 0$.

\textbf{1. Équation différentielle.}

Loi des mailles (la bobine est un dipôle de tension $u_L = r\,i + L\dfrac{di}{dt}$, convention récepteur) :
\[
E = r\,i + L\,\dfrac{di}{dt} \quad\Longleftrightarrow\quad \boxed{\dfrac{di}{dt} + \dfrac{r}{L}\,i = \dfrac{E}{L}}
\]

\textbf{2. Solution et constante de temps.}
\[
i(t) = \dfrac{E}{r}\left(1 - e^{-t/\tau}\right) \quad\text{avec}\quad \tau = \dfrac{L}{r} = \dfrac{\SI{0,20}{\henry}}{\SI{5,0}{\ohm}} = \SI{4,0e-2}{\second}.
\]
\[
\boxed{\tau = \SI{40}{\milli\second}} \qquad i(t) = \SI{2,4}{\ampere}\left(1 - e^{-t/\SI{0,040}{\second}}\right)
\]
Le courant permanent est $I_0 = \dfrac{E}{r} = \SI{2,4}{\ampere}$.

\textbf{3. F.é.m. d'auto-induction.}
\[
e_L(t) = -L\,\dfrac{di}{dt} = -L\times\dfrac{E}{r}\times\dfrac{1}{\tau}\,e^{-t/\tau} = -E\,e^{-t/\tau}.
\]
\begin{itemize}
    \item À $t = 0^+$ : $e_L(0^+) = -E = \boxed{-\SI{12}{\volt}}$ : la bobine s'oppose totalement à l'établissement du courant (le courant ne peut pas sauter de 0 à $\SI{2,4}{\ampere}$ instantanément).
    \item À $t = \tau$ : $e_L(\tau) = -E\,e^{-1} = -\dfrac{12}{e} \approx \boxed{-\SI{4,4}{\volt}}$.
\end{itemize}

\textbf{4. Vérification de la loi des mailles.}

Tension aux bornes de la résistance : $u_r = r\,i = E\left(1 - e^{-t/\tau}\right)$.
Tension aux bornes de la bobine (partie inductive) : $u_L = L\dfrac{di}{dt} = E\,e^{-t/\tau}$.
\[
u_r + u_L = E\left(1 - e^{-t/\tau}\right) + E\,e^{-t/\tau} = E \quad \forall t. \qquad \blacksquare
\]
La somme vaut bien $\SI{12}{\volt}$ à tout instant.
\end{corrige}

\begin{corrige}[Exercice 8 — Rupture de circuit : l'étincelle]
\textbf{Données :} $L = \SI{0,15}{\henry}$, $r = \SI{8,0}{\ohm}$, $I_0 = \SI{1,2}{\ampere}$, $\tau' = \SI{0,50}{\milli\second} = \SI{5,0e-4}{\second}$.

\textbf{1. Énergie magnétique stockée.}
\[
W_m = \dfrac{1}{2}L I_0^2 = \dfrac{1}{2}\times \SI{0,15}{\henry}\times (\SI{1,2}{\ampere})^2 = 0,5\times 0,15\times 1,44.
\]
\[
\boxed{W_m = \SI{0,108}{\joule} = \SI{108}{\milli\joule}}
\]

\textbf{2. F.é.m. maximale à la rupture.}

Le courant décroît selon $i(t) = I_0\,e^{-t/\tau'}$, donc $\left|\dfrac{di}{dt}\right|_{\max} = \dfrac{I_0}{\tau'}$ (à $t = 0$) :
\[
|e_L|_{\max} = L\,\dfrac{I_0}{\tau'} = \SI{0,15}{\henry}\times\dfrac{\SI{1,2}{\ampere}}{\SI{5,0e-4}{\second}} = \SI{0,15}{\henry}\times \SI{2,4e3}{\ampere\per\second}.
\]
\[
\boxed{|e_L|_{\max} = \SI{360}{\volt}}
\]
Soit 30 fois la tension de la pile ($\approx \SI{12}{\volt}$) : c'est la \cle{surtension de rupture}.

\textbf{3. Arc électrique ?}

Champ électrique entre les contacts espacés de $d = \SI{1,0}{\milli\meter} = \SI{1e-3}{\meter}$ :
\[
E_{\text{champ}} = \dfrac{|e_L|_{\max}}{d} = \dfrac{\SI{360}{\volt}}{\SI{1e-3}{\meter}} = \SI{3,6e5}{\volt\per\meter}.
\]
Ce champ est \emph{inférieur} au champ de rupture de l'air ($\approx \SI{3e6}{\volt\per\meter}$) : \textbf{l'arc ne s'amorce pas immédiatement}. En pratique, l'écartement au tout début de l'ouverture est bien plus petit ($d \ll \SI{1}{\milli\meter}$), le champ dépasse alors le champ de rupture et une étincelle jaillit ; c'est pourquoi on observe une étincelle à l'ouverture.

\textbf{4. Protection.}

On place une \cle{diode} (dite \emph{diode de roue libre}) en parallèle sur la bobine, montée en inverse vis-à-vis de l'alimentation. En fonctionnement normal, elle est bloquée et ne conduit pas. À l'ouverture de l'interrupteur, la f.é.m. d'auto-induction la rend passante : le courant de la bobine s'écoule alors en boucle fermée à travers la diode et décroît exponentiellement avec la constante $\tau = L/r$, \emph{sans surtension} aux bornes de l'interrupteur. L'énergie magnétique est dissipée par effet Joule dans la bobine et la diode.
\end{corrige}

\begin{corrige}[Exercice 9 — Probatoire 2022 type : Bobine tournante et éclairage LED]
\textbf{Données :} $N = 300$, $S = \SI{20}{\centi\meter\squared} = \SI{2,0e-3}{\square\meter}$, $r = \SI{15}{\ohm}$, $B_T = \SI{4,5e-5}{\tesla}$, $U_{\text{seuil}} = \SI{1,8}{\volt}$, $I_{\min} = \SI{5,0}{\milli\ampere}$, $f = \SI{2,0}{\hertz}$.

\textbf{1. F.é.m. induite.}

Vitesse angulaire : $\omega = 2\pi f = 4\pi\ \mathrm{rad/s} \approx \SI{12,6}{\radian\per\second}$.

Flux total : $\Phi_{\text{tot}}(t) = N B_T S\cos(\omega t)$, d'où :
\[
e(t) = -\dfrac{d\Phi_{\text{tot}}}{dt} = N B_T S\,\omega\sin(\omega t).
\]
Valeur maximale :
\[
E_{\max} = N B_T S\,\omega = 300\times \SI{4,5e-5}{\tesla}\times \SI{2,0e-3}{\square\meter}\times 4\pi.
\]
\[
E_{\max} = 300\times 4,5\times 10^{-5} = 1,35\times 10^{-2}\ ;\quad \times 2,0\times 10^{-3} = 2,7\times 10^{-5}\ ;\quad \times 12,57 \approx 3,4\times 10^{-4}\ \mathrm{V}.
\]
\[
\boxed{E_{\max} \approx \SI{0,34}{\milli\volt}} \qquad \boxed{E_{\text{eff}} = \dfrac{E_{\max}}{\sqrt{2}} \approx \SI{0,24}{\milli\volt}}
\]

\textbf{2. La DEL s'allume-t-elle ?}

La bobine a une résistance interne $r$ : la tension aux bornes de la DEL est $u_{\text{DEL}} = e - r\,i < e$ (modèle du générateur réel). Or même $E_{\max} \approx \SI{0,34}{\milli\volt} \ll U_{\text{seuil}} = \SI{1,8}{\volt}$. \textbf{La DEL ne s'allume pas} : la tension est environ 5000 fois trop faible.

\textbf{3. Avec $N' = 1200$ spires.}

$E_{\max}$ est proportionnelle à $N$ :
\[
E'_{\max} = E_{\max}\times\dfrac{1200}{300} = 4\times \SI{0,34}{\milli\volt} \approx \SI{1,36}{\milli\volt}.
\]
Toujours très inférieur à $\SI{1,8}{\volt}$ : \textbf{la DEL ne s'allume toujours pas}. (La résistance $r' = \SI{60}{\ohm}$ aggraverait encore la chute de tension interne.)

\textbf{4. Avec noyau de fer doux ($\mu_r = 200$).}

Le champ dans la bobine devient $B' = \mu_r B_T = 200\times \SI{4,5e-5}{\tesla} = \SI{9,0e-3}{\tesla}$ :
\[
E''_{\max} = \mu_r\,E'_{\max} = 200\times \SI{1,36e-3}{\volt} \approx \SI{0,27}{\volt}.
\]
Encore inférieur à $\SI{1,8}{\volt}$ ! \textbf{Conclusion :} même avec 1200 spires et un noyau de fer doux, le champ terrestre est trop faible et la rotation manuelle trop lente pour allumer une DEL. Il faudrait un aimant puissant (aimant néodyme, $B \sim \SI{0,3}{\tesla}$) : c'est le principe des lampes de poche dynamo. L'exercice illustre l'importance quantitative des ordres de grandeur.

\textbf{Barème indicatif (sur 10) :} expression de $e(t)$ : 2 pts ; $E_{\max}$ et $E_{\text{eff}}$ : 2 pts ; modèle générateur réel et conclusion : 2 pts ; proportionalité $E_{\max}\propto N$ : 1,5 pt ; effet de $\mu_r$ et conclusion argumentée : 2,5 pts.

\textbf{Points de vigilance :} convertir $S$ en $\mathrm{m}^2$ ; ne pas confondre valeur maximale et valeur efficace ($E_{\text{eff}} = E_{\max}/\sqrt{2}$, valable uniquement en sinusoïdal) ; justifier l'écart tension/courant avec le seuil de la DEL ; citer la loi de Faraday avant tout calcul.
\end{corrige}

\begin{corrige}[Exercice 10 — Probatoire 2021 type : Freinage magnétique et chute d'aimant]
\textbf{Données :} $m = \SI{50}{\gram} = \SI{0,050}{\kilogram}$, $L = \SI{1,0}{\meter}$, $v_{\text{lim}} = \SI{0,15}{\meter\per\second}$, $g = \SI{9,8}{\meter\per\second\squared}$.

\textbf{1. Stabilisation de la vitesse (loi de Lenz).}

Quand l'aimant descend, le flux de son champ magnétique varie à travers chaque section du tube de cuivre (conducteur). Des courants induits — les \cle{courants de Foucault} — naissent dans la paroi du tube, en boucles fermées autour de l'axe du tube. D'après la loi de Lenz :
\begin{itemize}
    \item \emph{au-dessous} de l'aimant, le flux (dirigé vers le bas, pôle Nord en bas) augmente : le courant induit crée un champ vers le haut, repoussant l'aimant ;
    \item \emph{au-dessus} de l'aimant, le flux diminue : le courant induit crée un champ vers le bas, attirant l'aimant vers le haut.
\end{itemize}
Les deux effets produisent une force globale dirigée \emph{vers le haut}, opposée à la chute. Cette force croît avec la vitesse (la variation de flux est plus rapide), jusqu'à compenser exactement le poids : la vitesse se stabilise à $v_{\text{lim}}$.

\begin{popfigure}
\begin{tikzpicture}[scale=0.9]
\draw[thick, popblue] (0,0) rectangle (4,6);
\draw[fill=poporangeL, draw=popdark] (1.6,3.4) rectangle (2.4,4.2);
\node at (2,3.8) {\small N};
\draw[->, thick, popdark] (2,3.3) -- (2,2.6) node[midway, right]{$\vec{v}$};
\draw[->, thick, popgreen] (0.7,4.8) arc (150:30:0.55);
\draw[->, thick, popgreen] (3.3,4.8) arc (30:-30:0.55);
\draw[->, thick, poppink] (0.7,2.2) arc (210:330:0.55);
\draw[->, thick, poppink] (3.3,2.2) arc (-30:-150:0.55);
\node[popgreen] at (0.7,5.3) {\small $i$ au-dessus};
\node[poppink] at (3.3,1.7) {\small $i$ au-dessous};
\end{tikzpicture}
\legende{Courants de Foucault induits dans le tube au passage de l'aimant (sens vus selon la loi de Lenz).}
\end{popfigure}

\textbf{2. Constante de freinage $k$.}

À vitesse limite, le mouvement est rectiligne uniforme : d'après le principe d'inertie, $\vec{P} + \vec{F}_{\text{mag}} = \vec{0}$, soit $m g = k\,v_{\text{lim}}$ :
\[
\boxed{k = \dfrac{m g}{v_{\text{lim}}} = \dfrac{\SI{0,050}{\kilogram}\times \SI{9,8}{\meter\per\second\squared}}{\SI{0,15}{\meter\per\second}} \approx \SI{3,3}{\newton\second\per\meter}}
\]

\textbf{3. Puissance dissipée par effet Joule.}

À vitesse limite, l'énergie potentielle perdue est intégralement convertie en chaleur. La puissance perdue par la pesanteur vaut :
\[
P_{\text{Joule}} = m\,g\,v_{\text{lim}} = \SI{0,050}{\kilogram}\times \SI{9,8}{\meter\per\second\squared}\times \SI{0,15}{\meter\per\second}.
\]
\[
\boxed{P_{\text{Joule}} \approx \SI{7,4e-2}{\watt} = \SI{74}{\milli\watt}}
\]
(Vérification : $P = F_{\text{mag}}\,v = k v_{\text{lim}}^2 = 3,27\times 0,0225 \approx \SI{74}{\milli\watt}$. Cohérent.)

\textbf{4. Tube fendu longitudinalement.}

La fente coupe les boucles de courant autour de l'axe : les courants de Foucault ne peuvent plus circuler en boucle fermée dans la paroi. La force de freinage magnétique devient quasi nulle, donc $k \to 0$ et :
\[
\boxed{v_{\text{lim}} \to \text{très grande (chute quasi libre)}}
\]
L'aimant tombe alors presque comme dans un tube isolant. C'est pourquoi on fend les noyaux et on feuillette les circuits magnétiques pour limiter les courants de Foucault.

\textbf{Barème indicatif (sur 10) :} explication par la loi de Lenz + schéma : 3 pts ; expression et calcul de $k$ : 2,5 pts ; bilan énergétique et $P_{\text{Joule}}$ : 2,5 pts ; effet de la fente justifié : 2 pts.

\textbf{Points de vigilance :} invoquer explicitement la loi de Lenz pour le \emph{sens} des courants ; citer le principe d'inertie à vitesse limite ; ne pas confondre énergie et puissance ($P = mgv$, pas $mgL$) ; préciser que la fente supprime les \emph{boucles} de courant.
\end{corrige}

\begin{corrige}[Exercice 11 — Probatoire 2023 type : Transformateur réel et auto-induction]
\textbf{Données :} $N_1 = 500$, $L_1 = \SI{2,0}{\henry}$, $r_1 = \SI{2,0}{\ohm}$, $N_2 = 50$, $L_2 = \SI{0,02}{\henry}$, $r_2 = \SI{0,05}{\ohm}$, $k = 0,98$, $U_1 = \SI{230}{\volt}$, $f = \SI{50}{\hertz}$, secondaire ouvert.

\textbf{1. Réactance du primaire.}
\[
\omega = 2\pi f = 100\pi \approx \SI{314}{\radian\per\second}.
\]
\[
\boxed{X_1 = L_1\omega = \SI{2,0}{\henry}\times \SI{314}{\radian\per\second} \approx \SI{628}{\ohm}}
\]
Comparaison : $X_1 \approx \SI{628}{\ohm} \gg r_1 = \SI{2,0}{\ohm}$ (rapport $\approx 314$). La chute ohmique est négligeable devant la f.é.m. d'auto-induction : \textbf{le primaire est un excellent inducteur} (bobine presque parfaite).

\textbf{2. Courant de magnétisation à vide.}

À vide ($i_2 = 0$), le primaire se comporte comme une bobine seule. Comme $X_1 \gg r_1$ :
\[
\boxed{I_{1,0} \approx \dfrac{U_1}{X_1} = \dfrac{\SI{230}{\volt}}{\SI{628}{\ohm}} \approx \SI{0,37}{\ampere}}
\]
(Le calcul exact avec $Z_1 = \sqrt{r_1^2 + X_1^2} \approx \SI{628}{\ohm}$ donne le même résultat à 0,001 \% près.)

Dans une inductance pure, le courant est en retard de $\dfrac{\pi}{2}$ sur la tension. Avec $u_1(t) = U_1\sqrt{2}\cos(\omega t)$ :
\[
\boxed{i_{1,0}(t) = I_{1,0}\sqrt{2}\cos\left(\omega t - \dfrac{\pi}{2}\right) = \SI{0,52}{\ampere}\cos\left(\SI{100\pi}{\radian\per\second} t - \dfrac{\pi}{2}\right)}
\]

\textbf{3. Tension secondaire à vide.}

Le flux commun étant le même pour les deux enroulements (couplage fort), les f.é.m. induites sont proportionnelles aux nombres de spires :
\[
\dfrac{U_{2,0}}{U_1} = \dfrac{N_2}{N_1} \quad\Rightarrow\quad U_{2,0} = U_1\,\dfrac{N_2}{N_1} = \SI{230}{\volt}\times\dfrac{50}{500}.
\]
\[
\boxed{U_{2,0} = \SI{23}{\volt}}
\]
Rapport de transformation : $m = \dfrac{N_2}{N_1} = 0,10$ : c'est un \textbf{abaisseur} de tension (adapté au soudage : forte intensité sous faible tension).

\textbf{4. Court-circuit et inductances de fuite.}

Inductances de fuite : $L_{\sigma 1} = (1-k)L_1 = 0,02\times \SI{2,0}{\henry} = \SI{0,040}{\henry}$ ; $L_{\sigma 2} = (1-k)L_2 = 0,02\times \SI{0,02}{\henry} = \SI{4,0e-4}{\henry}$.

Ramenée au primaire, l'inductance de fuite secondaire est multipliée par $\left(\dfrac{N_1}{N_2}\right)^2 = 100$ :
\[
L_{\sigma} = L_{\sigma 1} + L_{\sigma 2}\left(\dfrac{N_1}{N_2}\right)^2 = \SI{0,040}{\henry} + \SI{4,0e-4}{\henry}\times 100 = \SI{0,040}{\henry} + \SI{0,040}{\henry}.
\]
\[
\boxed{L_{\sigma} = \SI{0,080}{\henry}}
\]
Tension de fuite pour $I_{1,\text{cc}} = \SI{100}{\ampere}$ :
\[
U_\sigma = L_\sigma\,\omega\,I_{1,\text{cc}} = \SI{0,080}{\henry}\times \SI{314}{\radian\per\second}\times \SI{100}{\ampere} \approx \SI{2,5e3}{\volt}.
\]
\[
\boxed{U_\sigma \approx \SI{2,5}{\kilo\volt}}
\]

\emph{Commentaire :} cette valeur, supérieure à la tension secteur, montre qu'un courant de $\SI{100}{\ampere}$ ne peut pas réellement s'établir en court-circuit franc : l'impédance de fuite \emph{limite} le courant de court-circuit. C'est précisément l'intérêt pour le poste à souder : les fuites magnétiques (volontairement augmentées dans les transformateurs de soudage) stabilisent l'arc et empêchent la destruction du transformateur lorsque l'électrode touche la pièce (court-circuit). Le courant réel de court-circuit serait $I_{\text{cc}} \approx \dfrac{U_1}{L_\sigma\omega} = \dfrac{\SI{230}{\volt}}{\SI{25,1}{\ohm}} \approx \SI{9,2}{\ampere}$ au primaire, soit environ $\SI{92}{\ampere}$ au secondaire — courant typique de soudage à l'arc.

\textbf{Barème indicatif (sur 10) :} calcul de $X_1$ et comparaison : 2 pts ; courant de magnétisation (valeur + expression avec déphasage) : 3 pts ; rapport de transformation : 2 pts ; inductances de fuite, $U_\sigma$ et commentaire physique : 3 pts.

\textbf{Points de vigilance :} le déphasage de $\pi/2$ du courant sur la tension dans une inductance est exigible (le justifier par $u = L\,di/dt$) ; le rapport de transformation ne s'applique qu'\emph{à vide} (ou en charge pour un transformateur parfait) ; pour ramener une grandeur au primaire, multiplier par $(N_1/N_2)^2$ pour les impédances ; toujours commenter la cohérence physique du résultat.
\end{corrige}

\newpage

\coursec{Fiche bilan}

\begin{popfigure}
\begin{tikzpicture}[
  mindmap,
  grow cyclic,
  every node/.style={concept, text=white, font=\small\bfseries},
  root concept/.append style={concept color=popdark, font=\large\bfseries, minimum size=3.2cm},
  level 1 concept/.append style={sibling angle=60, level distance=4.2cm, minimum size=2.6cm},
  level 2 concept/.append style={sibling angle=45, level distance=2.8cm, minimum size=2.2cm, font=\footnotesize},
]
\node[root concept]{Induction électromagnétique}
  child[concept color=popblue]{ node {Mise en évidence}
    child[concept color=popblue!70]{ node {Aimant + bobine}}
    child[concept color=popblue!70]{ node {Inducteur / induit}}
  }
  child[concept color=poporange]{ node {Flux magnétique}
    child[concept color=poporange!70]{ node {$\Phi = NBS\cos\theta$}}
    child[concept color=poporange!70]{ node {Unité : weber (Wb)}}
  }
  child[concept color=popgreen]{ node {Loi de Lenz}
    child[concept color=popgreen!70]{ node {Effets opposés à la cause}}
    child[concept color=popgreen!70]{ node {Sens du courant induit}}
  }
  child[concept color=popred]{ node {Loi de Faraday}
    child[concept color=popred!70]{ node {$e = -\dfrac{d\Phi}{dt}$}}
    child[concept color=popred!70]{ node {$e$ en volts}}
  }
  child[concept color=poppurple]{ node {Auto-induction}
    child[concept color=poppurple!70]{ node {$\Phi = Li$ ; $e = -L\dfrac{di}{dt}$}}
    child[concept color=poppurple!70]{ node {$L$ en henry (H)}}
  }
  child[concept color=popgold]{ node {Énergie et applications}
    child[concept color=popgold!70]{ node {$E_L = \dfrac{1}{2}Li^2$}}
    child[concept color=popgold!70]{ node {Alternateur, transformateur}}
  };
\end{tikzpicture}
\legende{Carte mentale de la leçon : induction électromagnétique.}
\end{popfigure}

\begin{bilan}
\textbf{Définitions clés}
\begin{itemize}
  \item \cle{Induction électromagnétique} : apparition d'une f.é.m. (et d'un courant si le circuit est fermé) dans un circuit \emph{induit} soumis à une variation de flux magnétique produite par un circuit ou un aimant \emph{inducteur}.
  \item \cle{Flux magnétique} à travers $N$ spires planes de surface $S$ placées dans un champ uniforme $\vec{B}$ : $\Phi = NBS\cos\theta$, où $\theta$ est l'angle entre $\vec{B}$ et la normale $\vec{n}$ à la spire.
  \item \cle{Auto-induction} : un circuit parcouru par un courant variable est à la fois inducteur et induit ; il s'oppose aux variations de sa propre intensité.
  \item \cle{Inductance} $L$ : coefficient de proportionnalité entre le flux propre et l'intensité : $\Phi = Li$.
\end{itemize}

\textbf{Formules à connaître par cœur}
\begin{center}
\begin{tabularx}{\linewidth}{|l|X|l|}
\hline
\rowcolor{popblueL} \textbf{Grandeur / loi} & \textbf{Expression} & \textbf{Unités SI} \\
\hline
Flux magnétique & $\Phi = NBS\cos\theta$ & weber (Wb) ; $1~\text{Wb} = 1~\text{T}\cdot\text{m}^2$ \\
\hline
Loi de Faraday & $e = -\dfrac{d\Phi}{dt}$ & volt (V) \\
\hline
Loi de Lenz & le courant induit s'oppose, par ses effets, à la cause qui lui a donné naissance & --- \\
\hline
Flux propre & $\Phi = Li$ & $L$ en henry (H) \\
\hline
F.é.m. d'auto-induction & $e = -L\dfrac{di}{dt}$ & volt (V) ; $1~\text{H} = 1~\text{V}\cdot\text{s}\cdot\text{A}^{-1}$ \\
\hline
Énergie d'une bobine & $E_L = \dfrac{1}{2}Li^2$ & joule (J) \\
\hline
\end{tabularx}
\end{center}

\textbf{Méthodes express}
\begin{itemize}
  \item \textbf{Calculer un flux} : repérer $N$, $S$, $B$, l'angle $\theta$ avec la \emph{normale} (pas avec le plan !), puis appliquer $\Phi = NBS\cos\theta$.
  \item \textbf{Trouver le sens du courant induit (Lenz)} : (1) sens de $\vec{B}$ inducteur ; (2) le flux augmente-t-il ou diminue-t-il ? (3) le champ induit $\vec{B}_i$ s'oppose à cette variation ; (4) la règle du tire-bouchon (ou de la main droite) donne le sens de $i$.
  \item \textbf{Calculer une f.é.m.} : exprimer $\Phi(t)$, dériver par rapport au temps, appliquer $e = -d\Phi/dt$ ; la valeur absolue donne la grandeur, le signe (avec Lenz) donne le sens.
  \item \textbf{Auto-induction} : à la fermeture d'un circuit (bobine), le courant s'établit avec retard ; à l'ouverture, une étincelle peut naître (surintensité due à $di/dt$ très grand).
\end{itemize}
\end{bilan}

\begin{aretenir}[Checklist avant le Probatoire]
\begin{itemize}
  \item[$\square$] Je sais définir l'induction électromagnétique et citer une expérience qui la met en évidence.
  \item[$\square$] Je sais distinguer le circuit \cle{inducteur} (source du champ variable) du circuit \cle{induit} (siège de la f.é.m.).
  \item[$\square$] Je connais l'expression du flux $\Phi = NBS\cos\theta$ et je sais que $\theta$ se mesure entre $\vec{B}$ et la normale à la spire.
  \item[$\square$] Je connais l'unité du flux : le \cle{weber} (Wb), avec $1~\text{Wb} = 1~\text{T}\cdot\text{m}^2$.
  \item[$\square$] Je sais énoncer la loi de \cle{Lenz} et l'utiliser pour trouver le sens du courant induit.
  \item[$\square$] Je connais la loi de \cle{Faraday} : $e = -\dfrac{d\Phi}{dt}$, et je sais l'appliquer à un flux variable.
  \item[$\square$] Je sais que le signe moins traduit la loi de Lenz (opposition à la variation de flux).
  \item[$\square$] Je connais la f.é.m. d'auto-induction $e = -L\dfrac{di}{dt}$ et l'unité de $L$ : le \cle{henry} (H).
  \item[$\square$] Je sais expliquer le retard à l'établissement du courant dans une bobine et l'étincelle de rupture.
  \item[$\square$] Je connais l'énergie emmagasinée dans une bobine : $E_L = \dfrac{1}{2}Li^2$.
  \item[$\square$] Je sais vérifier l'homogénéité de chaque formule par analyse dimensionnelle.
  \item[$\square$] Je sais citer des applications : alternateur (barrages d'Edéa, Nachtigal), transformateur, cuisinière à induction.
\end{itemize}
\end{aretenir}

\findecours

\end{document}
